% source: https://learning.edx.org/course/course-v1:HarvardX+TinyML1+1T2022/home
% Fundamentals of TinyML
% google log in


% Bilder TinyML/Fundamentals

% Reihenfolge
%1. TinyML/Fundamentals
%2. TinyML/Applications
%3. TinyML/deployment
%4. TinyML/MLOps for Scaling TinyML

\chapter{Fundamentals of TinyML}

\section{Objectives}

Introduce background to TinyML and discuss the range of topics covered in this course and the follow up courses.

\section{What is this specialization all about?}

 Welcome to TinyML.
 
Tiny machine learning, or TinyML, is an emerging field that is at the intersection of machine learning and embedded systems.
An embedded system is a computing device that usually is small, or tiny, like the one that I'm holding in my hand, that operates with low power, extremely low power.
So much so that some of these devices can run for days, weeks, months, sometimes even years on something like a coin cell battery.
And TinyML is all about taking just a small, little computing device and enabling it with machine learning smarts.

Now, what is machine learning?

Machine learning is a subfield within AI, or artificial intelligence, that uses a lot of data to infer interesting patterns about new data that it's seeing.
Together, the intersection of embedded systems and machine learning is going to unlock a whole new world.



Before we get started, let's level the playing field and start with a crisp definition of what is TinyML.

TinyML, or tiny machine learning, is a fast-growing field of machine learning technologies and applications that include algorithms, hardware, and software that are capable of performing on-device sensor data analytics, doing so at extremely low power, typically in the order of milliwatts, so that we can enable always-on machine learning use cases on battery-powered devices.

A more intuitive example:

How many of you have said, OK, Google?
What happens when you say that?
The machine wakes up and it knows to respond to you.
OK, fine.
You might have said, hey, Siri.
Same thing, the machine wakes up and responds back to you.
All right, all right, fine.
Perhaps you have said Alexa.
No matter what you have said, they all represent an important application of what's in TinyML.


It's called keyword spotting or hotword detection.
Now, the key difference is that while they resemble the application of TinyML, these devices are largely plugged into the wall, or they require charging.
I'm talking about putting TinyML onto these tiny, tiny little processors that are on this board and then forgetting about them.
There's no power being given to them outside of that little battery.
That's the difference between the applications that we're seeing today and where we're going into the future.
Now, while this is an example application of TinyML, let me talk about some of the more forward-looking examples.
Let's talk about autonomous drones, drones that are completely independent, that operate fully by themselves.
Imagine autonomous drones that are able to do video surveillance in order to keep us safe by figuring out what's a threat, what's normal, and so forth.
Or they might be able to venture really far out into places where humans cannot go.
For instance, later on, we'll talk about how drones are being used in the Amazon forest for preserving our wildlife.
So for instance, when the forests are deforested, we actually can go in and autonomously plant new seeds so that we can keep the forest growing.
Or we can talk about mobile robots, which are extremely important in the service industry for taking care of the elderly in the future.
Well, these mobile robots need to be moving around, constantly sensing the information around them.
And machine learning is very good at that problem.
And we need to be able to pack in lots of sensors into these robots and do so with extreme efficiency, so they can operate on battery power.

Or think about smart glasses, right?
These are passive in the sense that they don't proactively do anything.
But imagine smart glasses in the future.
Imagine a situation where you go out into the mall, or marketplace, or you're at a party.
And how many times have you had to say, I'm sorry, could you repeat that again, right?
Imagine if these glasses had contextual hearing.
In other words, they're able to see where you're glancing.
And then based on that, be able to pick up the voice signals or the sound signals that are coming there and then amplify them so that only the things that are in your context will actually be richer.
That way, you have a better experience, right?
Or imagine, for instance, you're in a new country and someone's talking to you in a language that you don't know.
And imagine your glasses are able to do real-time machine translation and then be able to feed it into your brain.
These are some remarkable applications.
Or let's talk about TinyML in the medical area.
Here, I'm showing you Micra, which is from Medtronic.
It's the world's smallest pacemaker that goes into your heart.
It's small, it's non-invasive, and it's completely self-contained.
And if this device had machine learning in it, it would be able to, perhaps, proactively tell you how your heart's doing or be able to do things for you so that you have a better, healthy life choice.
Or let's talk about what Elon Musk is doing, where he's implanting chips inside the brain so that they can send signals to understand brain activity.
By measuring the electrical signals that are emitted by neurons, we might be able to learn things about our brain.
And in doing so, we might be able to cure diseases like treat depression, insomnia, or any other disease.
Now, let me take a step back and remind you that technology is not always about looking into the future.
Sometimes it is about what you can do here and now, and I'm going to tell you that we can use TinyML to make this world a better place.
We can use TinyML to take care wildlife.
We can use TinyML to improve the quality of our oceans or take care of the icecaps, things that are critical to our survival.
We might also be able to use TinyML to take care of our forests.
Now, over the course of the program, I'm going to introduce you to a variety of these different kinds of applications where the fundamental technology that sits in the palm of my hand, for instance, has the ability to actually have impact at a global level.
So let me be specific and give you an example.
Here, I'm talking about an ElephantEdge project.
There's a growing crisis in Africa right now about diminishing elephant population.
So organizations are coming together in order to protect the gentle beast.
Now, if we're not careful, these elephants can be wiped out in 10 years, so it's a very important thing.
So how do we help take care of the gentle giants?

Well, we need data.
We need lots of data.
So what if we could build some sort of intelligent device
like a smart collar that goes on these elephants that
is totally non-intrusive, does not harm them, but is there,
constantly sensing interesting signals so that we can help them have a better life?
Let me be specific.
For instance, if we had a smart collar, we might be able to do risk monitoring from poaching.
What if we could build machine learning models that could be put onto small, little devices where they were able to figure out if, for instance, the elephants were moving into a high-risk area and then be able to send real-time notification to park rangers to keep them safe?
The park rangers could then react to them and to make sure that the elephants go away from that danger zone.
Or perhaps, we can do human conflict monitoring.
We could build intelligent machine learning models that are able to sense and alert when elephants are heading into an area where, for instance, farmers live.
So that they can proactively get out of the way or do something in order to make sure that they're safe and the animals are safe.
We might be able to do activity monitoring.
We might be able to learn things about the animals.
When are they swimming, drinking water, or sleeping, and so forth?
And by learning all these things, we might be able to take better care of them.
We might also be able to do communication monitoring.
What we mean by that is that by sensing sound, we might be able to understand how these animals are actually
communicating together.
But in order to do all of that, we need TinyML.
We need to be right there where the data is, where that information, those sounds, those signals, where all of that is, so that we can do real-time machine learning.
These are all remarkable applications of TinyML.



So what I'm trying to say is that when you think of TinyML, don't just think about the future.


Think about the here and now, how this technology that's potentially in your hand, right, has the ability to transform the world.
Now, how are we going to achieve all of this, right?
Those are some grand applications.




How are we going to do program embedded microcontroller devices with machine learning capabilities?



\section{Who is this course aimed at?}





For starters, you might be a student.
You might be wanting to learn everything about ML.
You might be a student in high school, you might be a student in college, or you might just be a lifelong learner.
Either way, we want to help you expand your knowledge of ML.
Now, if you're new to machine learning, this is a great starting point for getting your feet wet with ML.
That said, if you know about ML, you will also learn something new.
Because we will be focused on tiny machine learning.
What are the aspects specifically important for deploying machine learning on tiny devices?
I do want you to understand one thing, is that this is an applied machine learning course.
In other words, we won't be getting caught up in theory and the math behind a lot of machine learning.
There are a lot of wonderful courses out there, and we will point you to them.
But in this course, we're going to focus very much on the applied side of deploying machine learning.
How do we engineer these systems?
Now, as a student, if you know about ML, you are definitely going to expand your knowledge.
Now, if you don't know about ML, you are definitely going to learn about ML.
Now, the second cohort is practitioners.
Now, you might be a practitioner who is interested in building a TinyML application, a product, a new innovation.
Well, if you want to be able to do that successfully, you're going to need to learn a lot about cross-layer design that's involved in TinyML.
Ranging from the applications to the run times to the hardware.
We're going to expose you to that vertical stack.
And by the end of this program, you will not only have designed tiny models, and learned the software stack that is needed for them, but you would also have effectively deployed it on a tiny little computer.
So that would give you full coverage.
And therefore, I believe you will be able to build a TinyML application by the end of this program.
Now, you might be an engineer working at a company.
And you might be thinking, hey, I want to hone my skills as the industry is evolving.
Well, this course is definitely meant for you in that case.
You could be an ML engineer wanting to deploy TinyML so that you can learn about new hardware that's coming out.
How do I take these big machine learning models and make them really small?
Well, you need to have a deep appreciation for the resource constraints that these small little embedded devices tend to impose.
Or it might be the other way around.
You might be a hardware engineer wanting to gain exposure to machine learning.
And you want to understand what are the concepts, and what does it mean to have a solution so that you can tackle the end to end system challenges for TinyML?



When we start off the course, we're not going to make any assumptions about what level of machine learning you have, or what level of embedded systems knowledge you have.
When we get into the fundamentals of TinyML, we're going to focus very much on just having the basic skills of scripting.
You've got to be able to program in Python.
Nothing too advanced, as programming is one of those things you can pick up on the fly, but you've got to know basic scripting. Where we will be focused on the applications of TinyML, we will be using a lot of what you know about Python and building applications in a Google Colab environment.
But for now, think of it as some sort of web page where you go to write your Python code.


Finally, this is where things change a little bit.
You will continue to program in Python, but you will also need a little bit of C and C++ programming background.
Nothing overly complicated like virtual functions and so forth.
Quite simple, you know, like C and C++.
And when things get tough, we'll be there to kind of help you nudge along the way, One other thing that we'll be introducing  is that, because we will physically be programming the microcontroller, you're going to need to use an IDE.
And so, that's an Interactive Developer Environment, and we're going to provide you an IDE that allows you
to program that microcontroller device.



\section{What will you learn?}


And across all threesteps, there are a certain set of fundamental principles we'll be going over.
Now, recall that step 1 focuses on the fundamentals of TinyML, where we're focusing primarily on the language of machine learning, with an interesting twist towards timing.
In step 2, we're going to be talking about different applications of TinyML, like keyword spotting, i.e., when you say, Alexa, or when you say, OK, Google.
There are many other interesting applications we'll touch on.
Step3, we're going to learn how to deploy that onto our tiny little microcontroller.

Now, across all three courses, there are a set of core areas that we'll be walking through.
And those are 

\begin{itemize}
  \item machine learning, 
  \item applications, and 
  \item embedded systems.
\end{itemize}

We will be touching on all three of them to the extent needed.
But the critical thing about this program is that we'll be focusing mostly on the interactions that are going to be happening.
For instance, you recall that we talked about ElephantEdge previously.
In ElephantEdge, what's the application?
The application is taking care or protecting the gentle giants.
Now, what's the ML aspect?
The ML aspect is actually doing risk monitoring, activity detection, and communication monitoring based on the sensors that are going to be on the caller.
So the critical nugget, the intersection, really is this, that you need to understand the fundamental needs of the application, and understand where, and why, and how the ML is going to serve a purpose.
In a similar way, when we're talking about TinyML applications, when we're looking at, hey, we've got this interesting thing that the application can benefit from ML, well, then the question becomes, how do I actually deploy this onto this tiny little microcontroller that has very limited capability in terms of computing and in terms of power consumption, and as long as -- they don't last forever if you keep running them with a lot of demand on the compute.
So how do you bridge these two worlds together?
Well, to do that, you need to understand the hardware's performance constraints and power constraints.
And at the same time, you need to figure out how you can take that machine learning and deploy it on the device in a way that it can run all the time.
So those interactions require you to understand not only the machine learning but also the hardware constraints, because it's that marriage that enables TinyML.
Now, there's also the aspect of knowing what the application needs and then intersecting that with what the hardware needs to be.
So you might have an interesting application, like the Smart Glasses that we spoke about earlier.
Well, what does this application really need in terms of the form factor?
Well, clearly, it's Smart Glasses.
So it has to be really thin.

But at the same time, that application demands you to have real-time characteristics.
For instance, if you're doing real-time speech translation, where someone is talking to you and you're listening to it, and it's translating the words, and they're going into your ears, well, you want that to happen in seemingly same time.
It cannot be like someone said something, and then 5 seconds later, it goes in your ear.
Well, that's not going to help with the conversation.
But all of that has to fit into the small little form factor, right?
And that's that intersection between the application needs these characteristics and the embedded system needs to be in this form factor and so forth.
You need to understand those interactions.
At the end of the day, it's all about putting each and every single one of these round circles together, because these domains need to intersect.
And that's when you actually enable TinyML.

And that's what we're going to focus on.


First and foremost, I want to say that in this class you're going to really focus on understanding the content.
But you're also going to have to practice it.
So in order to practice it, what we're going to be doing is a lot of hands-on learning.
So while you are reading about the concepts, you have a lot of assignments where you will actually be going through programming things in TensorFlow, which is the key thing.
But not only that, we will be talking about the different variants of TensorFlow.
And you need to understand this because these machine learning frameworks are specialized to do different tasks.
So if you want to enable TinyML, you need to understand how TensorFlow needs to be adapted in order to be able to run on a tiny little device.
And so we'll teach you about TensorFlow, TensorFlow Lite, and TensorFlow  Micro  and then deploying it onto the device.
Now, clearly you're going to have to solve certain tasks.
So you will have to do programming off that.
And you will be doing this in a Web IDE.
And I'll talk about Google Colab later if you're not familiar with it.
And you'll be deploying this onto a physical device by the end of the course.
So you're going to first train all the software.
And then finally, you will move it onto a physical hardware.
Now, this physical hardware is also going to have you interface with a couple of sensors, which is very exciting because not only are you just getting a small computer that you'll also be physically connecting of a different part.
So it'll give you a little bit of a taste for embedded systems.
All this is going to come packaged in a nice little box for you.
Or you will be able to order the parts for yourself.
When you open up this box, there are a whole bunch of different things that'll enable you to build a bunch of different TinyML applications.
Let's just take that small little microcontroller.
If you look at that microcontroller, this thing is jam-packed with sensors, because TinyML is all about sensors.
It's all about processing the data that's coming out of these sensors.
So you have a color, brightness, proximity sensor.
You have a digital microphone that can pick up sounds.
You've got a motion, vibration, orientation sensor.
You've got temperature, humidity, pressure sensor.
All of these different sensors can enable a variety of different application use cases.
Many of these sensors are quite critical to the kind of problems that people are trying to solve in the real world.
So you have the right building blocks here.
And all of this comes with an Arm Cortex-M4 microcontroller.
And the microcontroller is effectively the processing engine, the heart of the system.
And an Arm microcontroller is widely available in the real world.
It is deployed nearly everywhere.
So you will be dealing with a real microprocessor, where you could say that in the end class, that I've actually programmed this microcontroller to be able to do tiny machine learning, which is very exciting.
Trust me on this.
So you'll be not only dealing with the components of the board, but you'll also be packaging all of that excitement
into interesting applications.
For instance, we will be looking at something very similar to saying, OK, Google.
We will train on a small little model that will respond to your words.
And we'll use the vision sensor in order to be able to do interesting things from the camera input that's coming in, like detecting people.
And we'll use the accelerometer data and the gyroscope data in order to get it to respond to things when your hand is moving, for instance.
So we'll program some of these assignments in the software world.
And then in Course 3, we'll actually take those applications and learn how to deploy them.
Now, because it's an embedded system, you need to learn what the deployment means.
It's not as straightforward, and that's the beauty about it.
But once you learn the skill, you will have mastered a very, very powerful skill.




% >>>>>>>>>>>>>> Hier weiter


\section{What is tiny (ML) and why it's important for the future}


Now machine learning is a subfield of artificial intelligence that's focused on developing algorithms that can learn to solve problems by analyzing data for interesting patterns.

Here, we're going to specifically be looking at a very specific type of deep learning, which is a type of machine learning that leverages neural networks and big data to be able to make observations about interesting patterns in the data.


Well, there are other applications, like when you say, "OK Google," or, "Alexa," or when you talk to these systems.
Well, there's a lot of natural language processing going on in order to understand what you're saying.
Something that we already take for granted.
And then of course, there is the times when you go on social media and you give a thumbs up or you give a thumbs down.
All of that is helping drive our interest and our patterns of behavior on the network, because these machine learning systems are learning how to give us the information that we want to see more of.
For instance when I read the news, I tend to read a very particular type of news.
I'm not that much interested in all kinds of news.
But by giving a thumbs up or thumbs down, we're actually training the systems in order to help them serve us better.

Now these are great applications. Right?

Now, to be more specific, there are many specific subclasses of machine learning that are of interest.
For example, image classification, being able to tell whether the picture here has got a cat or dog.
If there's a cat a bunch of pixels are being fed in and ultimately the machine learning network says, hey, it's a cat.
Sometimes it might look at the same exact pictures and say it's a dog, which is completely incorrect.
So machine learning systems are statistical or stochastic in nature.
And so they're not always guaranteed to be correct.
This image classification, being able to say whether it's a cat or dog, is one particular type.
But of course, there are other interesting applications, such as object detection.
Object detection is basically taking another image, but in that image being able to precisely localize and say that it's actually an orange versus an apple.
Or you can even do instance level detections, where for every single orange that is in there, you're able to say that it is actually an orange, rather than just put a big box around all the oranges, saying that this whole collection is a set of oranges.
So that's another type.
Then there is segmentation, which is very heavily used, for instance, in autonomous cars, where the cars need to know where is exactly the human being in the picture.
Where are the lanes of the road?
Where are the shoulders in the road, for instance.
So segmentation is a way of kind of colorizing or pixelizing the entire image into different categories, where all the pixels for a certain thing are of the same color.
For instance, here, the purple color indicates that the side path is all tied together.
So that's a side path.
Then there is, of course, machine translation.
Here is an example that comes out of Google's pipeline, where we're looking at how machine translation is actually implemented.
Earlier, I was talking about machine translation on smart glasses, where you might have contextual hearing.
Well, if you want to do that, you've got to take a whole bunch of different words and then you've got to train a machine learning model.
Think of it, for now, as a black box.
And out of that comes this machine learning black box.
It's able to make predictions on a particular language, and then you're able to do natural language processing with it, yet another example.
Or recommendation systems, as I said earlier, you might be giving a thumbs up, thumbs down on different kinds of things, and that's actually driving your interactivity on the social networks.
Now, all of these capabilities require a remarkable amount of horsepower, remarkable amount of computing capabilities, so what companies are doing, they are taking all these computers and jam packing them into data centers, that are all just being dedicated in order to provide machine learning capabilities today.
These data centers are so big that they take up a big block.
So if you ever see a big cement block that's kind of running around that looks about the size of a football field, well that's probably a data center.
Now, what's really interesting is that in order to be able to provide that capability, companies like Google are building TPUs, tensor processing units.
Or companies like NVIDIA are building GPU, graphics processing units.
Both of these kinds of architectures, these computing systems, are capable of running machine learning extremely fast.
And that's needed, because at the rate you and I are using those machine learning services, when we give a thumbs down or a thumbs up, for instance, that's a remarkable amount of energy that's
being used.
So people are having to build custom designs in order to support all of the demanding needs of machining applications.
But building those big computing processors, packing them into big systems, and then putting them into a big giant data centers, well, big is not always better.

I'll tell you why. Think about it.


When you have a big data center like this, what can you do?
You can't have interactivity.
You can't carry a big data center in your pocket so you can take a picture.
Right?
Because when you take a picture, what do you want?
You want that system to be able to quickly respond back to you.
But if by the time you take a picture -- and let's say it has to go to some other computer somewhere else, well, that's going to be slow.
Right?
But imagine being able to put all that computer capability inside your pocket on your phone.
Right?
Why?
Because if you can do that, you can enable a variety of different applications that are much more responsive and interactive.
Like for example, when you say, "OK, Google," the machine can immediately wake up and you can say schedule a calendar appointment or something.
Or you might be able to do things like, "Hey, where in this text is," something interesting.
Q and A, as we call it, for instance.
Or when you take a picture, you want to enhance or remove something out of the picture.
Right?
So for example, in this picture, we're seeing that you're able to enhance the image quality.
There are all these things that you want to be able to do that are more close to you in a very responsive way.
Now, if you compare a big computing system to a small little device, right, what are the big differences?
If I pack a whole bunch of these big computers into a big cement block, well, that's a lot of power we're consuming.
And if I'm looking at a smartphone, well, there's very low power, because it's a small little battery-based device.
Now, if I look at it in terms of bandwidth, the amount of data that the data center can crunch, well, that's a lot of data that you can crunch, because all the machines are packed close to one another.
So as long as the data is nearby, you can move things really fast, and do a lot of high performance computing.
But if I'm on a smartphone, I don't have that much guarantee about the bandwidth, the ability to move the data off the phone or into the phone.
How many times have you had a poor network connection?
Where you're like, "Uh, why is this not loading?
Let me refresh."
Let me refresh that web page.
Or for instance, sometimes you might have great connectivity and other times you don't.
Right?
So there's no guarantee about what the connectivity in the phone looks like.
Or think about the use case, where for instance, you don't want to upload a lot of data or you don't want to download a lot of data because it costs you money.
Right?
So bandwidth, the amount of data you move and how quickly you can move it, is a critical concern.
So you want, typically on a smartphone, you have low bandwidth but it's responsive.
So in the data center, for instance, if I want to access it, it has a very high latency, because it's somewhere far away.
So I actually have to send the data to it.
So while I have a big amount of computational horsepower there, it's going to take me time to get there.
Versus is if I do that machine learning on my phone, well, it's right there on my hand.
So it can do it much faster than, imagine, moving all these bits over there and then bringing them all back.
But of course, if I have machine learning running on my phone all the time, right, and I wanted to be on all the time-- remember that's what I said.
The key thing is tiny ML.

This is about always on machine learning.
Well, what happens with your phone?
How many times have you had to plug your phone into the charger?
Sometimes I have to plug my phone in between half a day of my work, because I use my phone a lot.
Right?
And if it's constantly running something that is very demanding from the computing system, well, it's going to drain out of battery.
Now, this is where things get really interesting, because we want to move from data centers to smartphones.
But the reason we want to move that is because there's a lot of interactivity and a lot of interesting real time kind of behavioral patterns that we can exercise with machine learning, we can get even more as we make these devices completely pervasive in our homes.
Think about a phone.
A phone is only there wherever you are.
There are so many times I forget where I put my phone.
But if I have a smart device, in fact, I have about seven smart devices in my house, let alone my watches, these small little devices are increasingly becoming pervasive throughout us.
And we want all of these devices to be intelligent, because then there's a lot of information that they can sense and gather so that they can help us.
Certainly we have to be careful about that.
We'll talk about that more later.
For now I want you to understand the numerous possibilities, the kinds of devices that can all be intelligent.
And if they can all be intelligent that means they're all processing data on them real quickly, rather than getting the data, then sending it off somewhere, and then bringing it back, which costs a lot of power, which means they're going to drain out of the energy.
If you can process the data locally, then that's great.
Think about your smartwatch, it does machine learning on the device.
Think about my toothbrush.
My toothbrush is a smart toothbrush.
The Sonicare actually uses AI in order to help you brush better.
You can have smart earbuds.
Even a parking meter can be smart and tell you where to park or not park.
Or for instance, the Nest thermostat, which controls the temperature, is actually a smart device, because it learns based on how you change and when you change the thing.
So these devices are pervasive.
They are everywhere.
And imagine the possibilities if all these devices can actually be intelligent.
And that's what's tiny ML.



That's why people are so excited about tiny ML, because we're moving from thinking about all this amazing capability being some where in the cloud at some big corporation out there, versus you and I having that really rich experience right on our devices, being able to respond and interact with them in real time.
Now the question comes, where's the big data element?

Machine learning needs lots of data.
Well, where's all this data going to come from?
If you think about the big companies, like when you talk about Facebook or you talk about YouTube, well, yes, there are millions and billions of users all feeding the data up into them.
So yes, there's going to be big data and out of that big data, the genie of machine learning sprinkled some pixie dust, and out of that comes a neural network, this black box that's able to make predictions.
Well, when we go to small devices, where's the big data?
That's a good question.

If you look at the amount of data that's actually accessible at the end point, like at your watch or that toothbrush or that phone, there is an insane amount of data, because there's so many sensors that are out there at these end point devices.
The five quintillion bytes of data produced every day by IoT devices today, less than 1\% of all the data in these devices that we have today is actually being analyzed or used.
And if you think that today machine learning looks amazing, imagine what machine learning is going to do and look like if you analyzed that 99\% of that data that is just not being used today.
Can imagine how different a world we would be living in?
And that's what tiny amount is all about.
That's why we say that the future of machine learning is tiny and it's bright.
By being tiny, you can make it ubiquitous.
You can have it everywhere.
You can have it in every single endpoint device that is around you.
Think about it.


Sit down for a moment and think, just look around your room, find a device and think about what it would mean if that particular device was intelligent.
What if another device was intelligent?
What if another device was intelligent?
Can you imagine how-- not only is it about making an individual device interesting, but maybe these now intelligent devices might actually collectively be able to do something that each one of them was not able to do.
The possibilities are endless, and that's why the future of machine learning is tiny and it's bright.
And with that, I want to leave you with the key nugget, that there are many diverse applications of ML already in the real world, and it's increasingly moving from being up in the cloud, where there are big data centers, to the end point, where you and I are physically interacting with these devices.
And these end point devices, there are many, many of them around us.



\section{Learn about ML and tinyML applications}

\subsection{TinyML will soon be everywhere:}

TinyML will soon be everywhere, powering the next generation of smart embedded devices. These devices will be in our homes and in very remote locations, enabling remote monitoring for both industry and ecology. Today, in these remote monitoring settings, 99\% of raw sensor data is discarded, which is a wealth of data for machine learning! 

TinyML can provide a unique solution: by summarizing and analyzing data at the edge on low power embedded devices, TinyML can provide smart summary statistics that take these previously lost patterns, anomalies, and advanced analytics into account [1] [2].

In this reading, we survey a few emerging application areas that have great potential for TinyML. This list is a tiny (no pun intended) preview into the wealth of applications on the horizon. Later in this course we will do some basic review of machine learning, which some of you will need less than others but it is still a good review. In the next two courses we will be diving in much deeper around TinyML.

\subsection{Industrial Predictive Maintenance:}

In the industrial setting, TinyML is already being used to provide smarter sensing that enables advanced monitoring improving productivity and safety. For example, maintenance and monitoring of remote wind turbines can be quite challenging and time-consuming. However, if we could proactively predict that the machine will have trouble, we can predictively do maintenance ahead of any failures. Such “predictive maintenance” can lead to significant cost savings due to reduced downtimes, better availability of the systems for higher reliability in the product, which leads to overall higher quality of service for end-users/customers.

There are many TinyML applications for predictive maintenance. For instance, an Australian startup, Ping Services, has introduced a novel IoT device that continuously and autonomously inspects a turbine as it’s running. By magnetically attaching to the outside of any turbine (notice the small device in the image below) and analyzing detailed data at the edge and summary data in the cloud, the device can efficiently and effectively alert of any potential issues before a problem arises inside the turbine [3].

\begin{figure}
\GRAPHICSC{1.0}{1.0}{TinyML/1 Fundamentals/Windturbine}
  \caption{Ping monitor magnetically mounted on a turbine tower. The monitor is much smaller than the turbine}
\end{figure}  

\url{https://www.engineering.com/ElectronicsDesign/ElectronicsDesignArticles/ArticleID/20662/IoT-Device-Detects-Wind-Turbine-Faults-in-the-Field.aspx}

\subsection{Agriculture:}

Every day, the cassava crop provides food for more than 500 million African people. However, this vital stable is continuously under attack from a variety of diseases. The team at PlantVillage, led by Dr. Amanda Ramcharan, has developed the Nuru app to help farmers identify and treat these diseases. By running machine learning using TensorFlow Lite on mobile phones, the app enables real-time mitigation without the need for access to the internet -- a crucial requirement for many remote farmers (see the image below for the system in action). The next generation of this system will go farther -- leveraging tinyML and technologies like TensorFlow for Microcontrollers to deploy sensors across remote farms to enable better tracking and analysis [4].

\begin{figure}
    \GRAPHICSC{1.0}{1.0}{TinyML/1 Fundamentals/Agriculture}
    \caption{Two women using the PlantVillage app while in a field with plants}
\end{figure}  



\subsection{Healthcare:}

The Solar Scare Mosquito project deploys small smart Internet of Things (IoT) robotic platforms to help curb the spread of mosquito-borne epidemics such as Malaria, Dengue, and the Zika Virus. The system works by disrupting the mosquito breeding cycle by agitating water likely to contain mosquito larvae. The system uses rain and acoustic sensors to determine when it needs to agitate water to conserve battery and enable it to run on solar power indefinitely. It also sends smart summary statistics and alerts to warn of possible mosquito mass breeding events over lower power low-speed communication protocols. By making the system self-sufficient, small, and affordable, these devices can be deployed widely, preventing mosquitoes spread. All of the necessary components are included in a single component smaller than the size of a soccer ball [5].

\begin{figure}
   \GRAPHICSC{0.2}{1.0}{TinyML/1 Fundamentals/DeviceHealth}
    \caption{Diagram of the Solar Scare Mosquito 2.0 device. Shows the parts including the solar power, air pump, microphone for recording wingbeat frequency, water sensor, external connectivity.}
\end{figure}  




\subsection{Wildlife Conservation:}

\subsubsection{On the Land:}

TinyML is also already being used for ecological and environmental monitoring. For example, over the past 10 years, the Siliguri-Jalapaiguri railway line in India has had over 200 fatal collisions with elephants. Researchers from the Laboratory of Applied Bioacoustics at the Polytechnic University of Catalonia designed a smart acoustic and thermal sensor system using custom machine learning models running on solar power as an early warning system (see image below—all-in-one package with self-sustained energy source allows proximity to railway without added infrastructure, e.g., power lines) [6].

\begin{figure}
    \GRAPHICSC{1.0}{1.0}{TinyML/1 Fundamentals/Elephant}
    \caption{A team of researchers near train tracks testing an acoustic sensor so trains and elephants don't collide.}
\end{figure}  



\subsubsection{And In The Sea:}

Similar systems are also being deployed in the waterways around Seattle and Vancouver to prevent whale strikes in busy shipping lanes. These smart ML powered sensors enable constant real time monitoring and increased density of sensor deployments improving overall system efficiency and efficacy [7].

\begin{figure}
    \GRAPHICSC{1.0}{1.0}{TinyML/1 Fundamentals/Orca}
    \caption{Orca whale swimming in water}
\end{figure}  







\section{Applications that interest you?}


We hope these case studies got you excited and thinking about all of the ways that TinyML can be used to improve the world.  Here are three questions for starters: 

\begin{itemize}
  \item What kinds of new TinyML applications do you wish existed today?
  \item What existing applications do you think could be improved by TinyML?
  \item What applications that others have posted about are you excited about?
\end{itemize}




\section{How do we enable TinyML?}


We are  going to focus primarily on what it takes to enable TinyML.

 what does it take to make TinyML?
 
The two components are necessary:
There's an embedded systems component and then there's a machine learning component.
This combination is really what TinyML comes out of.
OK, well let's take an example of what that really means.
Now, let's go back to our favorite example about OK, Google, right?
Where we have a Google Assistant and then there's a physical device.
The assistant, being the virtual machine learning part, the device being the physical component.
Now, when you say "OK Google," the machine wakes up.
Let's break this process down into the three fundamental steps that are there.
The first step is that there is an input that's coming in.
This input is an audio input because you're saying something and then the machine is effectively picking up your audio signals.
Then there's a certain bit of processing that's actually happening.
And in doing that processing is where you're actually figuring out what are you trying to communicate.
Did you actually, in fact, mean to say, "OK Google," for instance, or "Alexa," or "hey, Siri?"
That processing happens.
And then there's a response that comes out where the machine physically generates some sort of output.
It might be lighting up the lights on the device to say that it's actually heard you.
For instance, on Alexa, you see a blue ring that goes around the top.
In Google devices, for instance, some sort of light shows up usually.
That's an actuation of the system responding back to you.
The three fundamental steps are input, process, and output.
This is fundamental to any embedded system.
And it's also true to any machine learning system.
That said, let's try and understand each one of these in greater detail.
So when we talk about input, the input that's going into this TinyML device for instance, what kind of input can we have?
Well, the critical thing about TinyML is that it's all about sensors.
There are many, many different kinds of sensors.
There are motion sensors.
There are acoustic sensors.
There are environmental sensors.
There are biometric sensors.
There are image sensors.
There are four sensors, lots and lots of different kinds of sensors.
I'm just showing you a general sprinkling of the possibilities here.
In the particular example, "OK, Google," for instance, it's an acoustic sensor because it's a microphone that is actually picking up the signal.
You might have biometric sensors for instance, things that detect a fingerprint, or your heart rate, and so forth.
Here are two examples I'm showing you here.
One is a glucose monitor from the Jacobs School at the University of California San Diego, where a simple little device just put on your skin can actually tell you what your glucose level is, which is very critical for certain patients.
Or you could have electrocardiography kind of devices, things where, effectively, they're able to, just from a single touch, be able to pick up the EKG signal and effectively plot the voltage over time and give you a representation of what the electrical activity really looks like.
Now, you can take that and you can couple that with much more intelligent activities like PPG.




Effectively, PPG is a means to actually analyze a signal at your blood conditions or changes in your blood conditions just by looking at light and how the light moves through the skin and hits the blood vessels.
It's a remarkable thing that you can do.
But if you had those kinds of fundamental sensors, right, which are inputs, then you can take that data And, you can do intelligent things with it.
Now, there are also other kinds of sensors, like image sensors, for instance.


The image sensor that's coming with a kit-- we have a small little camera.
It's a high resolution camera that actually interfaces with your microcontroller and you get vision, right?
There are many, many different kinds of sensors.
Now, you get all of that data that's coming in from a sensor, or sometimes multiple sensors, and you get to process it.
How do you process it?
Typically, in order to run machine learning, you need big processors.
Like, you need the big GPU processors, for instance.
Or you need something like the Google TPU that I showed you in the previous video.
Now, that's conventionally how people have been thinking about machine learning, big systems.
It's a complex workload and it means big computing horsepower.
OK, well, what does big really mean?
Let's use a proxy here for indicating what is big and what's small.
Let's look at the amount of area that it actually takes for the processor, the brains of the actual system.
Right?
If you look down here, the big processors are in the [INAUDIBLE]
for over 500 millimeters squared.
That is the effective area off the chip that we're talking about.
Then let's talk about something smaller, like your phone, for instance.
If you look at your phone, the A12 biometric processor from Apple, for instance, is much, much smaller.
Right?
It's only 83 millimeters squares, significantly smaller than the big GPU, for instance.
Then, of course, you have your smartwatch.
Smartwatch processor is even smaller, for instance.
Right?
If you look inside, this thing is only about 30 millimeters squared.
Oh, that's pretty tiny.
Now we're starting to get there from big, to small, to tiny.
But we're not there quite yet.
Let me show you something quite fascinating.
We're just getting started.
You ready for this?
Look at something really small, like this KL03 processor from Kinetis.
It's 3.2 millimeters squared.
I really want you to go out and get a ruler right now and try and draw an area of 3.2 millimeters squared.
It's a really small processor.
And it consumes very little amount of power.
In fact, it's a very small processor because it has little memory, little compute-- we'll be talking about these things at a level deeper in the next set of videos.
But before we get there, the context that I'm trying to get you to is that machine learning is typically thought of being run on big systems.
And now when we're talking about tiny, these are the kind of systems we're talking about.
Just to kind of put that into perspective here is that little processor sitting on a golf ball.
OK?
Do you know how big a golf ball is?
And then do you know how small each of the circles is on a golf ball?
It's smaller than the circle, one of those little dents that you have on a golf ball.
So it's really, really small is the point that I'm getting to.
And if it's that small, you can imagine putting it virtually anywhere.
You probably won't even notice it.




There are over 250 billion of these kinds of small, little processors everywhere in the world today.
These are called MCUs, or microcontroller units.
They are pervasive.
They are here already.
They're everywhere.
They're inside your automotive car.
They're inside your little iron.
They're inside your little coffeemaker, everywhere.
Now, the question is, can you make them smart by putting tiny machine learning on top of it?
So, let's talk about these small, little processors.
If there are 250 billion today, there are going to be many, many, many, many, many more in the coming few years.
Now, this plot on the x-axis is showing you years.
And the y-axis is showing you the millions of units that are going to get shipped, right?
Now, look at the projected forecast values.
These things are just going to continue growing.
The demand for microcontrollers in the world is insatiable because we are building so many new devices now.
So we're going to see this increasing trend.
Now, what's going to happen?
Is the price going to increase or decrease?
Well, as the volume of these microcontrollers that are needed increases, it's fundamental economics.
The price will drop because they are so pervasive.
So the manufacturing cost, effectively, goes down.
That's what this plot is trying to show you, that the price of a microcontroller is going to be below \$0.50.
It's almost going to be like \$0.50, which is really, really, really cheap.
Once you put this into the context of what these devices are capable of doing, especially when you enable TinyML on top of that.
These devices are pervasive, these devices are going to increase in quantity, these devices are very, very cheap.


Now, what's also interesting is that these devices consume very little power.
To put things into perspective, a big GPU like the one we showed you consumes 300 watts.
You need to plug that thing into a serious wall connector.
Now, as you go into something small, like the processor that's inside your phone, for instance, well, it consumes significantly less power.
It only consumes about three watts of power.
That's 100x smaller.
Well, when you go to microcontrollers like the Syntiant processor that we're looking at here, which we'll actually come back and learn a bit more about later down, their machine learning processor only consumes 140 microwatts. It's astoundingly small.



And this is what tiny machine learning computers are all about.
It's about consuming milliwatts or less.
And so these are very specific and concrete numbers
of what the future is holding.
So imagine if these processors consumed so little power, then we might actually be able to run them with real coin cell batteries.
Because the coin cell battery is going to look very, very big compared to the small, little processor.
And that's that potential we're talking about.
And then that's that processing part.
Now we come into the output.
So let's say you make some-- you take in the input, you process it, and now you generate some output.
What kind of output can a TinyML device do?
Well, you can have physical actuation by, like, activating servos, for instance.
Or you might be able to trigger speakers, generate some kind of signal.
In the case "OK, Google," that's effectively what we're doing.
The machine is waking up and it might say "hello" or something to that effect.
But it's not always about physical actuation of physical devices.
You can actually also have digital actuations.
In other words, you do some input process and you extrapolate some kind of interesting data.
And you send that digital signal out to a screen.
And then that's useful information.
So what are MCUs really enabling?
Well, fundamentally, it starts from the fact that-- MCUs are really small, so they can be completely pervasive.
Second is that they consume very little power, which makes them really practical to be able to deploy for a long time on small, little batteries, right?
So you don't need to tether them. You don't need to keep recharging them.
Then, they're really cheap, which means they can be everywhere.
And given that we're starting to build more number of devices, the demand for them is going to go up.
And if you can put machine learning on them, you're going to uncover a whole new set of applications.
And that's the beauty about it.
By building these small, little, tiny embedded systems and coupling them with machine learning-- which we'll talk about soon-- by building these two things together, we're going to unlock TinyML.
MCUs, or microcontroller units, are these small, little processors that are the fundamental building block of TinyML devices.

That's number one.
Number two, these devices are going to be pervasive or ubiquitous because they're cheap and they're going to be high in demand.
And number three, the question is, are they capable of running machinery models in them?





\section{What are the computing system challenges for TinyML?}

What are the challenges for TinyML?
Now, as always, TinyML is a composition of embedded systems and machine learning.
I'm going to continue to focus on the embedded system aspect before we transition over to the machine learning.
And specifically, I want to talk about what are the challenges in the embedded realm?
Because that sets the context, the scope, for how we deploy machine learning into that box.
So, as with any system, an embedded system is like a regular computer.
It's got components, like the hardware and the software.
Both of these are critical to understand at an embedded system level, in order to know how we can actually deploy TinyML on them, or what we need to do differently in order to deploy TinyML on them.
Hardware is the first area of focus for us.
Now typically, when you take any kind of computing system, the hardware is made up of three fundamental building blocks: the compute, the memory, and the storage.
The compute is the brain.
It's effectively where you do the processing.
The memory is where you temporarily store data.
It's like a short-term memory.
And your storage is your permanent hard disk.
When you take the power off, it still remembers your data.
So these are the three fundamental building blocks of any computing system.
Now when you put these three things together, there is a difference between what a microcontroller is and what a microprocessor is.
Now, if you're watching this video on a laptop or something like that, you've got a microprocessor that's doing a lot of the heavy lifting.
A microprocessor is a much bigger piece of computing engine, versus a microcontroller, the thing that we were talking about in the previous video, of which there over 250 billion of them out in the world.
Is a very different computing system as a whole.
Let me emphasize these differences for you.
When you take a laptop for instance, yes, there's a small part, which is called a microprocessor, but that microprocessor interfaces with a bunch of different system components.
It's only one part of a bigger puzzle that ties together.
So you have to buy the microprocessor.
And then you put all the other components together.
That's how you build a traditional desktop, for instance.
A microcontroller, on the other hand, is very different.
When you buy a microcontroller, it's already integrated.
For instance, the processor and the memory and the storage are all tightly, already coupled together.
It's one big monolithic block.
Big here, is relative, of course.
Here's a complete example for a microcontroller.
You've got your memory.
You've got your RAM.
You've got your processor.
Everything is fused into one single block that you buy, which is quite unlike when you buy your parts for a regular desktop computer, where you get to cherry pick and choose these different things.
And you can build them in any configuration you want.
For instance, my laptop has 32 gigabytes of memory, which is something I custom ordered from Apple.
Versus is if you buy the stock laptop, it only comes with 16 gigabytes of memory.
So I have this option to configure what kind of system I wanted.
But if I'm looking at a microcontroller, mm-hmm.
It's all packaged already into one single thing.
So it's either I buy this microcontroller, or that configuration of the controller.
That's one fundamental difference, major difference I want you to keep in mind.
If I was to kind of generalize these into critical bullet points, I would say the microprocessor is the heart of a general purpose computing system, where a microcontroller is the heart of an embedded system.

% >>>>>>>>>>>>>> Hier weiter

It's an embedded system, versus here, it's a general computer.
A microprocessor is just one part, and it interfaces with external components of memory and storage, versus, in the context of a microcontroller, these things are all tightly integrated already into the system.
Again, a microprocessor is more meant for general purpose systems like laptops, desktops, and servers, where you get to configure things.
Microcontrollers are typically designed for fixed function tasks, like be able to play music on an MP3 player, or be able to answer phone calls.
That's a very specific task.
So microcontrollers generally tend to be quite preset in their functionality.
And often, a big embedded system might include a phone, which is a very complex embedded system.
Might include a whole bunch of different microcontrollers that are all fused together, in order to be able to provide different kinds of functionality.
Now because of flexibility, a microprocessor ends up being very big, in fact, because your laptop is bigger.
You get to put all these parts together.
Those interfacing capabilities naturally lend themselves to being big.
A microcontroller is all about making them tiny, because this form factor is small, so you tend to squeeze things down.
Now, these key differences, if I was to step back up just kind of put some number ballpark numbers around, in order to kind of set the context of what a microcontroller really looks like, well let's take a look.
From a computing perspective, if you look at what's called the frequency or the clock rate, how fast the processor can run.
If you look at a regular microprocessor, it can run between 1 gigahertz and 4 gigahertz, versus a microcontroller can run between 1 Hertz and 4 Hertz.
From a memory standpoint, you can have megabytes to gigabytes of memory, versus you only tend to have small little kilobytes of memory on a microcontroller.
Storage tends to be in the order of gigabytes to terabytes, versus in a microcontroller, it tends to be kilobytes to megabytes.
From a power consumption it's in the order of 10's of watts, versus is in the microcontroller, as I showed you in the previous video, it's in the order of microwatts or milliwatts.
So it's a very, very big difference.
There is stark contrast.
It's in the order of magnitude difference, in terms of the compute, memory, storage, and the power consumption.




and so forth.
So the devices that we're talking about are extremely small.
That's the key takeaway down here.
Now what are the implications of all of this?
What does it mean to have less amount of memory?
What does it mean to have lower power consumption?
Well, that means you got to think about how complicated is the task that we're running?
We're going to learn about the interaction between the machine learning, the embattled system, and the actual application.
Well, if you want to understand the task, if I want to have smart glasses for instance, if it needs extremely high real-time processing, that is being able to have computing capability where everything runs extremely fast, well, microcontroller might not be a good fit.
That's not to say that it's not always a good fit.
It depends on what task you're trying to perform.
Secondly, it's like how much memory do I really need?
Microcontrollers come in all sorts of different forms and sizes.
For example here, we're looking at anything between two kilobytes and 512 kilobytes.
If you understand the task you want to run, and you understand the requirements, then you'll be able to pick the right microcontroller.
And the other question is, how long does a job have to perform?
Is it something that needs to be constantly running, or is it running intermittently?
Well that has implications in terms of the power consumption.

So, the key takeaway that I'm trying to emphasize down here, is that the microcontrollers that we're talking about, which are ubiquitous, the beauty about them is that they're ubiquitous.
But the downside is that they're very constrained, in terms of the capabilities.
But that's OK.
There are means and ways to which we can overcome those obstacles.
We'll get through those things.
But once you understand the confines, the scope, and the limits, then you will have a deeper appreciation for what it means to be able to deploy intelligence onto these heavily-resourced constrained devices.
And in the next video, we'll go a little bit further, talking more about the software side of the constraints. 


\section{What are more computing system challenges for TinyML?}


I'm going to pick up on where we left off-- what are the challenges for enabling TinyML?
Specifically, I'm going to continue to focus on the embedded systems side of the world.
Now, recall that I said that when we're talking about an embedded system, there is a hardware component and then there's a software component.
Previously we talked about the hardware.
Now we're going to shift our focus to talk about some sampling of the issues that we're going to face when we look at the software.
What is software?

Typically, what you find is that there are generally three levels of abstractions.
There is a high-level application, there are the libraries that provide support for those applications, and then there's the operating system that provides support for the libraries and the applications.
Now, I'm going to start off by talking about the operating system.
And specifically, what I want to help you understand here are the critical differences that you often find between what a general purpose computing system like a laptop or a desktop has versus what you will end up finding on an embedded system.
And this difference that we see is worlds apart.
And it's extremely important that you learn this difference, because it's through learning this difference that you gain deep appreciation for what a TinyML system really looks like, because that lays the foundation for building more complicated things on top of it.
So when I talk about an operating system, when I'm saying an operating system for a laptop, desktop, or a server, by and large you will probably pick one of three operating systems-- Microsoft Windows, Mac, or the Linux operating system.
And if I talk about things on the smartphone side of the world, there are two operating systems.
If I were to flip a coin, it would either be heads or tails.
In this case, it might be either the Apple iOS or it might be the Android operating system.
Now what's the beauty about having these widely deployed operating systems that run on most of the devices?
Let's take the Android ecosystem for instance.
An operating system and its associated parts are like the foundations for a house in my head.
If you have a solid foundation, you can build a ground floor.
On top of the ground floor, you can build a first floor.
You can build a second floor.
You have a lot of flexibility in what you can do.
But it all comes down to having the right foundation.
Now, in the context of the Android operating system, for instance, here I'm showing you the platform architecture of an Android OS.
Because it's got the right pieces in place, you can build high-level abstractions of libraries and stack on top.
Why is this important?


Because if you provide the right levels of abstraction, then you can enable end developers who are at the highest level, application developers, to be able to write very interesting pieces of code, like the applications that you and I have on our phones.
I have almost over 100 applications on my smartphone.
Each one does a specialized thing.
Now, that's the beauty if you have a right general-purpose operating system.
This is very much true even if you look at a Windows, Linux, or Mac operating system.
You have a nice base.
And then on top of that, you can build different kinds of applications, which is very interesting.
Now, when you go to an embedded system, what happens?
This is where I would say that the world kind of falls apart a little bit.
You see, in the embedded ecosystem, there isn't really that much of an operating system to say.
Why?
Because you take an embedded system, you typically want to specialize it to perform one task.
Think of your smartphone.
Do you only make phone calls on it?
Or do you listen to music?
Or are you surfing the web?
And are you on social media?
Or are you doing both at the same time nowadays?
You can do all those things.
You have all these flexible knobs and features that you have.
But an embedded system is not a general-purpose system.

An embedded system is typically designed to perform one task, like wake up when I say "OK, Google." Something specialized like that.
Or maybe it's designed to just do health monitoring and timing.
When you have something specialized like that, then you typically don't care so much for having a general-purpose platform.
Therefore, you don't see a lot of opportunity in terms of establishing an operating system.
This is not to say that embedded systems don't have operating systems.
I'm giving you here two examples-- FreeRTOS and Arm Mbed OS.
These are "operating systems" that exist, quote, unquote.
But typically, you don't tend to have these installed, because when you add any kind of extra stack on the embedded system, it takes away resources.
Remember, we're talking about memory and compute cycles that are very diminished and small, which means that if I put an Mbed OS or an RTOS, it might very well take up a little bit of those few kilobytes of memory I have, which means less left for the application.
That becomes a problem, don't you think?
So, therefore, you don't tend to see that.
If you don't see a good baseline, then building efficient systems actually becomes quite critical, because every one of these systems is specialized-- too much specialization.
So, let's go up one level and talk about libraries.
So we know at the embedded system level we don't have that sort of uniform, nice Windows, Linux kind of ecosystem.
What happens when you go to the library level?
Here's a simple piece of code that calls np.SaveTheWorld.
np is a library function call.
And it calls it 10 times.
We know it's a library function call, because we say import numpy as np.
This library is a Python library that's almost ubiquitously available on any given system.
The beauty of this library is that once you write the code in Python it can run on virtually any system.
It can be an Intel processor.
It can be an AMD processor.
It does not matter.
The code just runs uniformly.
You don't have to worry about it.
You just focus on running your code.
And then seamlessly the code will run on any hardware.
It creates a level of abstraction.
So you don't need to worry about the nitty gritty details of the lower level.
And that's the beauty about a general-purpose system.
Now this portability is something that we take for granted.
And we don't think about much.
But when we go over to an embedded system, this becomes a real problem.
Let me give you a very specific example by talking about the pi value.
What is pi?
22 divided by 7.
Very simple.



Well, that translates into a long floating-point number, or a number with a lot of values after the decimal place.
I just chose pi at 22 by 7 because it's an exciting value, nothing special.
But the key point here is that we have a floating-point value.
Now, in a machine, a floating-point value is typically expressed as three components here-- sign, exponent, and the mantissa, or fraction.
Don't worry about the details that's all available in the single-precision format for floating-point or double-precision format for floating-point numbers.
The key thing that I want to talk about is how a floating-point value is operated upon.
Typically, in any machine, if I have a value with decimal places, I might have two values, and I want to be able to add, multiply, divide, do whatever I want to do and generate some resulting value.
Now, I can choose to implement this functionality as a library, as a function call, just the way I did with np.SaveTheWorld.
And I'd say np dot multiply these two floating-point numbers.
Well, that tends to be slow.
So over the years, what's ended up happening is that we've gone from implementing these things as emulation mechanisms and software to actually physically putting data pack as we call it-- a floating-point data pack-- into the silicon itself, into the chip.
And if we do this, well, the chip ends up running super fast.
This was actually inside the Pentium processors for instance-- Intel i486.
If you open it up and you look inside, the die photo as we call it, there's a dedicated FPU that actually runs the floating-point arithmetic.
And therefore, it can run really fast.
That's great, because now, since then, it's proliferated, and virtually every chip has a floating-point unit inside it.
Therefore, you write the code once, and things run anywhere.
But that exposes a trade-off there.
If you want universal code portability across any system, you write code once.
Don't worry, it'll somehow magically run on any Intel processor or on any AMD processor.
That's great, but you pay in terms of the cost.
It costs money to actually physically implement it.
It takes extra power, because when you turn the unit on, it burns power, because it's actually triggering logic.
And it takes a lot of engineering effort to actually build it.
And sometimes those enduring efforts can actually be quite costly.



Well, go back and look up as the FDIV bug by Intel.
It's actually a real floating-point divide bug that Intel had, which got in mistakenly back in the '90s into one of its chips.
So it was literally generating incorrect results when you were doing a divide instruction.
So they had to go and pull all these things back-- all these chips back-- and then fix it.
Anyway, back onto the main line.
The point is that you want uniform code portability, well, you're going to have to pay a price in terms of these things.
Now, I could say forget the portability issues.
I'm an embedded system.
I want to be lean and mean and efficient.
So I don't want to worry about general-purpose portability.
If you run on my device, don't worry, you will run really well.
And by the way, my device is going to be very cheap.
It's going to consume low power.
And it's going to be much more simpler to engineer, because that's what I choose to do.
I choose to vertically integrate things and shave things off.
Now, that's beautiful from one system.
But what happens if I go to another system, and in another system, they actually have the floating-point capability?
So assume that I start off on the right-hand side with the checkbox where have this wonderful FPU capability,
and I write my code in my machine learning model accordingly.
And then I move it over to another device, and it doesn't run, because I made the assumption that the processor had a capability for the floating-point unit.

The code that I've specialized for one system is now not portable.
So this raises a fundamental question that you will now need to be thinking about.
How do we enable TinyML uniformly across these vastly different embedded systems if there is lower platform portability?
This is a critical challenge.


Embedded hardware is extremely limited in terms of performance, power consumption, and storage.

Now, I'm talking about embedded software.
And I said it's not as portable and as flexible as mainstream computing systems, because you don't tend to have these nice operating systems.
You don't only have these nice ubiquitous libraries that run arbitrarily on any given embedded device.
They tend to be very specialized.
And we need to understand these kinds of nuances.
We will be focusing on these kinds of things, because having this sort of knowledge is what is going to make you a much stronger machine learning engineer, for instance, who understands what it means to deploy TinyML.





\section{What are the ML algorithm challenges for TinyML?}


We are yet again on the topic of what are the challenges for TinyML.
While previously we were talking about embedded systems, now it's time to shift our gear a little bit over to machine learning.
Just the way an embedded system, both from a hardware and software perspective, were severely challenged, on a machine learning side, there are increasing trends that we're observing that might actually prove quite challenging to be able to unlock TinyML.
But again, that's not a challenge that cannot be overcome.
It's more of being able to critically think about what we need to do in order to be able to enable TinyML.
So let me start off by talking about what are some interesting trends that we're seeing in machine learning and where it's going.
Machine learning, think of it as a box.
And for now, think of it as how big or small the box is.
So what I'm showing you on the x-axis over time are effectively state-of-the-art models, machine learning models specifically in the context of NLP, natural language processing models and how big they've been getting over the recent years.
On the y-axis are what we call as the number of parameters.
Don't worry about it for now.
Once we get into the fundamentals of TinyML and when you start talking about the actual machine learning mechanics, the internals, these are all terminologies you will get quite familiar with in case you're not familiar right now.
Now, what we're seeing on the x-axis is time.
And what we are seeing is that as time has been going by, the complexity of the models is increasing quite steadily.
Now, if you look at the most recent model, which is called GPD3 from OpenAI, you will see that it's a remarkably complicated model as compared to all the other ones that
are down there.
Now, don't let the scale fool you.
This model is so big, or this box is so big that the rest of them look quite small when, in fact, just a few months, a few years ago, those other models looked like they were big differences in terms of their sizes.

The key point  is that machine learning models are growing fast in complexity.
They're growing extremely fast.
And moreover, in order to be able to keep up with them, the computing capability that is needed is growing remarkably fast.
So, for instance, from 2012 when the boom of machine learning really came about thanks to AlexNet and the use of big GPUs, over the recent few years what we have seen this is remarkable escalation in terms of the computing capability needed to be able to get those models out there.
So, for instance, between 2012 and roughly where we are now, the computing needs have grown by almost 300,000 times.
300,000 times!



That's an astounding amount of computational horsepower increase.
So at the same time, they were able to provide these black box machine learning models that we get are able to provide incredible capabilities to us.
But that said, it's requiring quite a bit of compute.
Now, if I was to go even further back before AlexNet in 2012 all the way back to the original source of history for machine running, back to the 1960s when all of this stuff was coming about, you'll see that back from 1960s to, say, maybe around 2010, there was roughly a steady, linear fast kind of pace.
And then once we got into the recent years, there's been a remarkable steep increase in terms of the computational capability requirements.
What is this saying?
That our interest in machine learning has grown so much and we are spending so many cycles on this that the pace of innovation needed for unlocking it now is quite remarkable.
Therefore companies like Google have been building these custom solutions, custom processors.
Just the way you have an Intel or an AMD processor that is able to do general purpose compute, companies like Google have been building custom processors just to run one task, just the machine learning task.
And they've been called TPUs.
These are called Cloud TPUs because they live in the cloud, not literally in the cloud, but in big data centers that
are remotely accessible to us.
If you look at these cloud TPUs, oh, they look like big skyscrapers.
They are physically about this much.
And the heat sink on them is about this tall.
The heat sink is there just to keep the processor cool
because it has to take out so much heat because the computing is so intensive.
And what are the companies doing?
They're taking all of these processors and jam packing them
into these big racks.
So there are multiple cloud TPUs all packed in, row after row,
column after column, right next to each other.
And all of these things are being sandwiched into big data centers that
are about the size of a football field.
And they consume so much power that they actually
have to be physically put right next to sources where there's water
or, like, for example, in this picture of the Google data
center in the Netherlands, you are seeing that it's actually
right next to renewable energy because it costs a lot of money
to power these things.
The fundamental message that I'm trying to get across
is that machine learning is becoming so computationally
hungry that we have to build big processors just dedicated for them,
then we have to sandwich and pack all of them together,
and then we have to incorporate these big data centers almost
near renewable resources and where water is
so we can actually keep these things cool and power them efficiently.
That's all wonderful.
Now, how does this compare to TinyML?
While things are getting bigger, I'm talking to you about the other side.
I'm talking to you about TinyML.
I'm talking to you about resource constraints
and how to pack these things that are tiny little devices.
So, the world is bifurcating in this.
But that's the beauty.
While there's a lot of interest in this, there's
a lot of interest in the capabilities of TinyML
because that's where the data lives.
You remember I said there's about less than 1%
of the data that's actually being used or analyzed for machine learning today.
So despite the fact that we're building all
of these remarkable machine learning processes and data centers,
that's only dealing with a small fraction of the data
that the world actually has access to where the sensors are
producing a remarkable amount of data.
And if you want to tap into that, we've got
to figure out how we're going to take that big model
and squeeze it right down into this small little form factor.
Now, that's quite the challenge.


So, how have things been evolving just generally in machine learning?
What this plot is specifically showing is
on the x-axis is the amount of computational power needed,
how many multiplies and accumulates and so forth need
to be done, all the arithmetic.
And the y-axis is the accuracy, how accurately in saying
it's a dog or a cat.
And the size of the circle is telling me how big
or how lean and skinny the model is.
So there are three pieces of information in this chart.
Now, what I want to do is walk you through a little bit
of that critical evolution and then set the stage
as to where things need to go.
So, the first and foremost model that people refer to
is AlexNet, which happened in 2012.
It basically was trying to predict a thousand
classes from ImageNet data set.
We'll talk about these things in much greater detail.
But the key thing is that it is able to get an accuracy of 57.1%.
And its model size was 61 megabytes in size.
Then we wanted better accuracy.
We want to be able to say a cat is a cat or dog is a dog.
And we want to do that with better accuracy than 50-something percent.
So VGGNet came along in 2014 as an example.
That boosted the accuracy to 71.5%.
But look at the size of the circle.
What happened?
We went from something that was smaller in the order of about 60 megabytes
to something that's almost 10 times as big at 528 megabytes.
That's enormous!
Well, we realized we can't just keep making things bigger,
so we tried to make them a bit more efficient
while also improving the accuracy.
So in 2015, Microsoft released ResNet, residual nets.
And this particular network improved the accuracy to 75.8\%
while shrinking the model size as it was getting better.

So we were doing great.
We were starting to improve the machine learning models
in order to make them both accurate and also be more cognizant of the size.
Then as smartphones became very prevalent
for machine learning deployments, that's when a major shift happened.
And we saw the evolution move towards mobile nets.
With mobile nets, we were saying that, no, the size is extremely critical.
We cannot just think about accuracy.
We have to make the thing small because it has to fit into your smartphone.
So what did we do?
We compromised on the accuracy, but we dramatically
shrunk the size of the network.
We made it only 16.9 megabytes.
What I'm trying to get across to you here
is that if you look at these couple of milestone markers
that I have cherry picked for you, what I'm trying to say
is that in our pursuit for better accuracy, better understanding
of the data that we're being given and being
able to see a cat is a cat or dog is a dog and being able to do so correctly,
pushed us towards naively making things bigger.
But slowly we understood the resource constraints of computing systems,
that we cannot just keep making things bigger.
Bigger's not always better.
So then we strive to think of ways to actually make things smaller.
As we go through this course, we will understand
the fundamental nuggets behind a way something
is much more efficient than another thing.
But the critical message here is to understand
the philosophical transitions that we're making as our journey continues
through machine learning.
Now, that's all great.
MobileNet, this is an awesome example.
Since then there have been other networks.
But the critical point is that the networks
that we have today, the models, they're still pretty big.
They're 17 megabytes or maybe they're a handful of megabytes.
It doesn't matter.
The key nugget is that our little embedded microcontrollers only
have a few kilobytes of memory.
With a few kilobytes of memory, it's an order
of magnitude difference between where usually state of the art sits
and how we need to cram things in.
So what did we do?
Well, this is where things get interesting for us.
Over the course of this program, we're going
to learn how to take these models that are performing certain tasks
and be able to compress them down.
We will learn about a whole slew of techniques.



\section{What are more ML algorithm challenges for TinyML?}


What are the challenges in terms of making machine learning models that are big-- how do we squeeze them down?
I'm going to give you a preview some of the things that we will be looking at in a lot more detail later on.
Now, from a model perspective-- one of the key things is when you have a machine learning model, it's all about how well can you shrink or compress the model down without losing its ability to fundamentally look for patterns in data?
Because that's what machine learning is, remember?
So, how can we go about doing this?
Let me give you a quick sampling for the kind of things that we will be getting into a lot more detail later on.
So, some examples and techniques that we'll be talking about are like pruning.
Pruning is this idea that if I have a network, then think of this network as having a bunch of connections.
And that's effectively a model.
When I prune, what I want to do is I want to be able to take some of the connections away.
And when I take some of the connections away, the question is, is this left hand side figure able to produce the same correct result as the right hand result?
Sometimes it can actually lose the accuracy-- as we call it-- or sometimes it might be able to maintain the accuracy.
We might also be able to remove certain neurons completely out of it.
Thereby, further reducing the amount of computation that's actually needed.
And in doing so, we shrink the size of the network and we also reduce the computational demand of the network.
And therefore, we can effectively squeeze this onto a tiny amount of ice.
And this is a very primitive technique that we will actually be using later on to be able to shrink our models down and deploy them efficiently.
Another thing that we can do is actually play around with the numerical representation.
The values that we actually compute on.
Because recall that embedded systems don't have fancy hardware or they tend to do very simple sort of arithmetic because they are supposed to be lean and mean as compared to general-purpose computing systems.
So a technique that we will be using quite heavily is called quantization.
Quantization is basically this notion that you have a floating point value.
Remember 22 by seven?
That's a floating point value.
You take that floating point value, which
is a long number of digital decimal places in place.
And you quantize it down.
Quantization means discretizing the values down
to a small subset of values.
So here, the end eight is basically going from negative 128 to 127,
which allows you to represent 256 values.
Now because you are now only having 256 values,
you have to shrink that whole range down,
which means you might lose some sort of values that would otherwise be
represented in a floating point value.
Why in the world would I ever want to do that?
Well, a floating 0.32 bit value takes a whole four bytes
to represent-- four bytes.
Versus an end eight value is simply one byte,
which means I get a 4x reductions by simply doing the quantization which is
a huge boom or boost in the model size.
So now I've shrunk it automatically down by 4x.
And there are a whole bunch of associated benefits
that come along with it from a system implementation sort of standpoint
that we'll talk about.
But quantization is extremely critical for us
because we are dealing with tinyML systems
where we don't have too much memory and we don't have to much compute.
Talking about compute, floating point values, as I said previously,
floating point values need dedicated hardware
often to be able to compute efficiently because floating point decimal
arithmetic is actually quite complex.
Versus int8, or just integer arithmetic, it is much more simpler
and it's usually found on every single embedded device.
So the key lesson here is that as we start thinking about embedded device,
you want to think about how big the networks are, what sort of precisions
they are being represented with, the numerical precision
and numerical formats-- there are a slew of things that we need to learn.
And I will teach you some of these things,
the critical ones especially, as we progress through.
There also some advanced techniques that are on the horizon
that are quite important for TinyML.
For example, things like knowledge distillation.
Knowledge distillation is a really nice fancy buzzword
for explaining how a teacher, who knows a lot from years of wisdom
and has a lot of information, is able to distill down
the critical information for the student without losing
the nugget given a particular task.
So the benefit of this, ultimately, is that a teacher network might be big.
A student network might be small.
Later on, we'll get into how a teacher network actually translates over
to a student network and how we can actually use this mechanism
to take something that a bigger machine learning task
knows how to do and be able to shrink things
down so that the tiny, small little machine learning model actually
knows how to do things better.
Again, this is just a sneak preview of what's coming down the pipe later.
Now, that's more from an algorithm and modeling sort of perspective,
what the algorithms look like.
But there is also stuff you need to know at the runtime level.
So, at the runtime level, you need to start thinking about the hardware
first.


When you have this big cloud TPU in its Google building,
where they want to be able to run many different kinds of machine learning
tasks, for instance, you need a framework, a software framework.
And so, typically, people deal with TensorFlow.
TensorFlow is a beautiful framework in which
you can write machine learning code.
Now if I want to do machine learning on a smaller device,
like a smartphone, which has got a much smaller computing capability, if you
recall from the previous videos--
the A12 processor, for instance-- is much smaller.
Well, if you do that, you realize that this particular system has
much less memory than a big cloud TPU normally would.
And then it would also have much less compute power,
and it's only focused on making inferences.
In other words, a cloud TPU might actually
have to learn how to detect patterns in the data.
The smartphone is always going to be just basically looking
for patterns in the data.
There's a big difference where you're learning how to do things
and where you only execute.
Learning something takes a lot of time.
For instance, I bet you right now you're learning about tinyMLs, so your brain is exploding with information that it has to process and figure out all the connections. 




Training versus inference.
Your smartphone is only looking for interesting patterns, for instance.
So if that's the case, then you don't want to necessarily use TensorFlow
because it's a rather big framework.
What you'd rather have is a smaller, leaner and meaner software framework
that is dedicated to just meeting the specific requirements that
are needed on a smartphone.
For instance, something that takes less memory.
Something that has less computational power requirements.
Something that only focuses on looking for patterns of data, quote,
unquote, "inference."



So if I were to put these two systems next to one another head to head, I'm going to start at the ultimate bottom line and work my way upwards since this is the most important point here.
Who are these frameworks targeted at?
For instance, like, if you look at TensorFlow versus TensorFlow Lite, these are two different software framework.
TensorFlow is really targeted at a machine learning researcher who is trying to figure out how to actually build the machine learning algorithm.
TensorFlow Lite is really focused on an application developer who is really interested in using the machine learning algorithm that comes out and integrating that into an application that provides end user service to you and me.
That's the fundamental difference.
And because of this, the frameworks have fundamentally different philosophies.
For example, TensorFlow has to do distributed compute.
It has to use many cloud TPUs, for instance, or many GPUs all at the same time to try and get the machine learning algorithm to learn the patterns that it finds in the data.
Well, we're not learning things on the mobile device.
We're just exercising what we already learned, so you don't need to communicate with a whole bunch of different smartphones.
That's a simplification in the software stack. The binary size, when you're running things in big data centers with those big processors that are all coupled together, well, you're not really worried about how much memory the software framework is taking.
But if you're talking about putting that little TensorFlow framework onto a smartphone, you're really concerned about how much memory it's consuming because you don't want it to use up all the space.
And once we get into the machine learning model itself because we're dealing with-- we're training.


In other words, we're learning.
There are things that are variable, just like our neurons are
variable in our head.
We're constantly adjusting as we are trying to learn things.
But once we learn it, we don't even think about it.
We just execute things really fast.
It's almost like a muscle memory or gut reaction.
Same thing for how the topology of the network looks like.
Topology is nothing but a fancy buzzword for saying,
what do your connections look like.
And in TensorFlow, because you're still learning to find patterns in data,
you're constantly adjusting things, and you need to be flexible.
But again, on TensorFlow Lite, when I already
know the task that I'm going to perform, I don't need any flexibility.
All of these things, these are just a sampling of the things.
But the key thing I'm trying to get to is
that there are big differences based on the system you're targeting
and what the software is going to look like,
and you need to understand the software to be
able to deploy things efficiently.
The critical difference here between TensorFlow and TensorFlow Lite
is the difference as to where you're deploying it.
So what's a typical pipeline look like?
The first time when you're learning a model,
if you want to be able to enhance a picture in your camera--
and what you have to do is you have to learn how to enhance that.
You have to find interesting patterns and data
and learn how to make that happen.
After that, once you've already figured that out,
you just want to make it like a function call.
You just call it.
So we have something.
Let's call it TensorFlow Lite Converter that
converts a big beefy model into something lean and mean.
That gives you a file that's called .tflite that's nice and thin.
And you take that, and then you can deploy this
into different kinds of smartphone devices.
So this is your general pipe.
You started with a lot of flexibility, and then you shrink things down,
and you make them lean and mean, and then
you target them for the specific device.
This is effectively what we need to enable the next major step going
from big processors over to smartphone processors
into the next smallest little thing.
The question is, how do we do it?
These devices have even less memory.
They have even less power as we've looked at before.
And they're only focused on inference, very much like a smartphone.
So for this, we're going to need to look at a software stack, which
is a TensorFlow Lite software stack.
Specifically, if you look at the training pipeline, for instance,
or how things happen, you use TensorFlow right at the beginning,
then after that were the conversion, the optimization, the deployment,
and the usage is all kind of tied down with TensorFlow Lite.
And when you deploy the model, you're typically
going to be deploying it on any given kind of system.
On a smartphone, it's in iOS or Android or maybe a Linux
kind of platform, which you might target a bigger embedded kind of device.
Or, in the case of microcontrollers, you want
to deploy it what you need is a TF Lite engine,
and then you can finally start using it for doing tasks.
So I'm just showing a smartphone here as an example,
but the general idea is that once you manage
to squeeze things out and put them on the device,
then you can start repeatedly using it.
For all of that, we'll be focusing on TensorFlow Lite.
So we'll walk you through from TensorFlow to TensorFlow Lite
to, in fact, what's called TensorFlow Micro.
That's at the runtime level.
So we talked about the models, we talked about the runtimes,
and at the hardware as well there's a lot of innovation going on specifically
to make tinyML very efficient.
Now, in this series of courses, I'm not going
to emphasize that much on the hardware.
I'll give you inclinations about the kind of things that are coming down.
You will be knowledgeable by the time you finish this program in terms
of what kind of trends and patterns we're
seeing in terms of the industry, where the industry is going.
But, that said, this is going to require you
to take a dedicated course that's specifically focused
on building machinery and hardware.
Nonetheless I will give you enough clues and tips all through the thing in order
to cover all three layers.
Now, with that said, what I'm trying to wrap up here
is that in the last four videos, including
this one, what we saw was specifically some of the challenges
that we're going to have in terms of the embedded system side
and the machine learning side in order bring them together
to fit inside a tiny little device.



\section{Why the Future of ML is Tiny}

\subsection{Why the Future of Machine Learning is Tiny and Bright}

\subsubsection{Tiny Computers Are Ubiquitous}

It is estimated there are over 250 tiny microcontrollers in the world and that over 50 billion will be sold this year. Microcontrollers (or MCUs) are packages containing a small CPU with possibly just a few kilobytes of RAM and are embedded in consumer, medical, automotive, and industrial devices. To put MCUs into perspective, there are only about 10 million servers in use across the planet and 4 billion smartphone users. Microcontrollers typically do not get much attention because they are often used to replace functionality that older electro-mechanical systems could do, in cars, washing machines, or remote controls, but they are ubiquitous.

Driven by new applications in wearables, automotive use cases, smart appliances, etc., the demand for MCU is anticipated to increase steadily. Figure 1 shows the data from IC Insights. The number of units sold in 2023 is expected to exceed 35 Billion units per year!


\begin{figure}
    \GRAPHICSC{0.4}{1.0}{TinyML/1 Fundamentals/DemandForecast}
    \caption{Microcontroller (MCU) Demand Forecast showing actual sales (in millions of units) for 2016-2018 and forecast for 2019-2023. Showing a linear expected increase from 20,000 million units in 2016 to almost 40,000 million units in 2023. Source IC insights.}
\end{figure}  


\subsection{Microcontrollers are Cheap}

MCUs are designed to be cheap enough to include in almost any object that’s sold. They are cheap because they aren’t designed to be general-purpose computational workhorses that run complex workloads. Instead, they are often designed to perform specific tasks, slowly, albeit steadily. Therefore, they usually cost much less than a typical processor. The magnitude difference in price between an average MCU and an Intel or AMD processor found in your laptop or computer might be an order or two orders of magnitude difference in price. 

The average price of an MCU is already less than USD \$1. As seen in Figure 2, the average sales price (ASP) for a microcontroller is hovering close to 55 cents, and as the demand surges, it is expected that this will fall below 50 cents in the future. On a global scale, the microcontroller market (\$M) is poised to grow by USD \$6.74 Billion between 2020-2024.

\begin{figure}
    \GRAPHICSC{0.4}{1.0}{TinyML/1 Fundamentals/PricingForecast}
    \caption{Microcontroller (MCU) Pricing Forecast showing actual average selling price (in US dollars) for 2016-2018 and forecast for 2019-2023. Showing a decrease from \$0.75 in 2016 to $\approx$ \$0.63 in 2018 and a linear expected decrease to almost \$0.55 in 2023. Source IC insights.}
\end{figure}  







\subsection{MCUs are Resource-Constrained, Ultra-low Power Systems}

The holy grail for almost any embedded device is for it to be deployable anywhere and require no maintenance like docking or battery replacement. Any device that requires tethered electricity faces a lot of deployment barriers. It can be restricted to only places with electrical wiring. Even where electrical wiring capability is available, it may be challenging for practical reasons to plug something new in, for example, on a factory floor or in an operating theatre. Putting something high up in the corner of a room means running a cord or figuring out alternatives like power-over-ethernet. The electronics required to step-down or convert the main-line high voltage to low voltage are expensive and waste energy.

But perhaps the most significant barrier to achieving ubiquitous deployment computers everywhere is how much energy an electronic system uses. Here are some rough numbers for standard components based on figures from \href{https://www.cambridge.org/us/academic/subjects/engineering/wireless-communications/smartphone-energy-consumption-modeling-and-optimization}{Smartphone Energy Consumption}, and so it is no wonder that smartphones need to be tethered to the wall and recharged every night:

\begin{itemize}
  \item A display might use 400 milliwatts.
  \item Active cell radio might use 800 milliwatts.
  \item Bluetooth might use 100 milliwatts.
  \item Accelerometer is 21 milliwatts.
  \item Gyroscope is 130 milliwatts.
  \item GPS is 176 milliwatts.
\end{itemize}

A key opportunity with using MCUs is they require a very minimal amount of energy. A microcontroller itself might only use a milliwatt or even less. Still, you can see that peripherals (accelerometer, gyroscope, GPU, etc.) require much more energy to stay powered on.

A coin battery might have 2,500 Joules of energy to offer, so even something drawing only one milliwatt will have severe consequences. Of course, most current products use “duty cycling” and power naps to avoid being always on, but we see how tight the budget is even then. The overall thing to take away from these figures is that while processors and sensors can scale their power usage down to microwatt ranges, displays and especially radios are constrained to much higher consumption, with even low-power wifi and bluetooth using tens of milliwatts when active. The physics of moving data around is generally well-known to require a lot of energy.

The general wisdom is that the energy an operation takes is proportional to how far you have to send the bits. CPUs and sensors send bits a few millimeters and are cheap. Radio sends them meters or more and is expensive. We don’t see this relationship fundamentally changing, even as technology improves overall. We expect the relative gap between computing and radio costs to get more expansive because we see more opportunities to reduce computing power usage.

\subsection{Tiny Machine Learning (TinyML) on MCUs}

Can you imagine a future where every one of the 250B tiny MCUs runs intelligent machine learning algorithms that can sense their surroundings, predict events in real-time, and even make nudges or recommendations based on the sensors' activity? We can transform the world.

For instance, smart consumer devices like smart toothbrushes will have MCUs that are tightly coupled with sensors that dynamically adjust to your brushing intensity based on pressure. Medical devices in the future may use microcontrollers with biomedical sensors to control drug delivery as and when needed. Microcontrollers in automotive applications can aid functional safety on the road, detecting engine conditions, etc. to ensure our safety. Finally, MCUs will likely have widespread applications in industrial devices for continuous process monitoring and anomaly detection for applications such as predictive maintenance that can save millions of dollars in productivity loss and downtime by forewarning maintenance engineers.

TinyML is a relatively new machine learning paradigm, yet it is producing remarkably astounding results. Several recent examples include audio, visual, and sensor fusion applications, such as voice and facial recognition, voice commands, and natural language processing.

Looking more further into the future, we imagine a world where we have a tiny battery-powered image sensor that I could program to look out for things like particular crop pests or weeds and send an alert when one was spotted. These could be scattered around fields and guide interventions like weeding or pesticides in a much more environmentally friendly way.

\subsection{TinyML at the Edge/Endpoint}

Computing data locally, rather than streaming all the data to the cloud or edge servers, which is costly, is a feasible solution for intelligently processing the data available at the sensors. This will not only conserve energy, but it will also help unlock new applications and use cases. Consequently, given the existing widespread deployment of MCUs, and the minimal energy consumption of MCUs, there is growing interest in providing AI functionality at the edge/endpoint. This idea isn’t particularly new. Both Apple and Google run always-on machine learning (ML) based neural networks for voice recognition on these kinds of power-efficient chips.

There are many more examples of AI capabilities at the endpoint, and people are just starting to realize that there is a unique match between deep learning and MCUs. If you are with us so far, then it’s obvious there's a massive untapped market waiting to be unlocked with the right technology to enable AI on these cheap and ubiquitous devices. We need ML solutions that can work on cheap microcontrollers, which use very little energy, that rely on computing capability, not radio, and can turn all our continuously streaming sensor data into something useful. This is the gap that machine learning, and specifically deep learning, fills.

It has suddenly become possible to take noisy signals like images, audio, or accelerometers and extract meaning from them in the last few years by using neural networks. Because we can run these networks on microcontrollers and sensors themselves use little power, it becomes possible to interpret much more of the sensor data we’re currently ignoring. For example, imagine we want to see that almost every device has a simple voice interface. By understanding a small vocabulary, and maybe using an image sensor to do gaze detection, we should control nearly anything in our environment without needing to reach it to press a button or use a phone app. We want to see a voice interface component that’s less than 50 cents that runs on a coin battery for a year, and we believe it’s possible with the technology we have right now.

\subsection{The Future of TinyML is Bright}

We can conjure up a thousand other products. But to get there, first and foremost, Pete, Laurence, and I are most passionate about ensuring that the technological imperative behind them is so compelling that we will all be able to build whole new applications that we can’t even imagine today. For all of us, it feels a lot like being a kid in the Eighties when the first home computers emerged. None of us had any idea what they would become, and most people at the time used them for games or storing address books, but there were so many possibilities. A new world emerged out of those burgeoning systems. So we challenge you to invent the future. Come on and embrace the future with us by learning all we have to teach you about tinyML!





\section{Responsible AI/ML}

Here we are, with the countless number of possibilities for TinyML despite the fact that we know that there are some hard problems to be solved both on the hardware side from the embedded systems perspective as well as from the software and algorithmic perspective on the machine learning side.
That said, we are excited about TinyML, aren't we?
Because I've shown you that there are a large number of possibilities about what you can do.
In fact, there are probably possibilities that I'm not even imagining that you are already thinking about.
And that's what's exciting about TinyML.
Now, technology is powerful.
It can bring about a lot of change.
It can bring about both good and bad change.
So one of the key things that I want to make sure you learn as we go through this program is that you understand what it means to deploy machine learning out in the field and to do that in a responsible manner.
And that's what responsible AI is all about.
AI offers us great benefits.
But at the same time, it introduces new risks and ethical challenges that we must be very conscious about.
We cannot build AI into all these tiny, little devices and scatter them out everywhere thinking that only positive things are going to come out of it.
That would be rather naive. 
Think, for instance, that just by having a simple little phone on you, this little device-- if you look at all the gyroscope data just out of that one thing, you can actually figure out the unique individual.
How so?



Well, because the way that you and I walk, there is a certain gait or cadence or rhythm to it.
Believe it or not, that's true.
And if you track that data, which these devices are totally capable of doing today, then people would be able to identify you precisely without ever asking for your unique name or your ID or anything like that.
And that's just one dangerous example, because that infringes upon our privacy and security aspects.
So I'm not saying this in order to scare you.
Rather, I want to help you understand that there is a big opportunity here.
If we think about AI being everywhere-- I mean, think about the new risks and ethical challenges proactively-- we can actually use that to uncover new opportunities.
And that's what I'm going to hope to convey throughout this video.
As new technologies come about, they often tend to be disruptive.
They're not always neutral.
And in other words, you cannot assume that they're not going to cause any
damage.
Why?


Because AI can bring about changes in current practices and the way we do things, because now machines are starting to do some of the things for us, they can influence human decisions.
Think about it.
For instance, when my phone says "breathe" or, for instance, when a sensor gives me a nudge to do something because it's picking up some kind of interesting pattern, it's influencing my behavior.
Now, imagine if these sensors are all around us.
Would you even know if they're influencing our behavior?
Or they could, worse, be regulating our behavior.
So consciously or unconsciously, we might be falling susceptible to the [? own ?] technologies that we're building.
And it's extremely crucial that as we engineer these systems, we open our mind to looking at technology a bit differently and more holistically.
Not just looking at the use cases, but understanding the potential ramifications and implications of those use cases.
Let me tell you, for instance.
There was a research that was done that shows that 65\% of Americans believe that companies often fail to anticipate how their products and services impact society.
Most people are wigged out by this.
I would like to think that any piece of product or any technology that I'm bringing into my house has been carefully thought about and that it's not so much that I have a piece of technology and, oops, company says, I didn't know that.
Right?
So there's a little bit of a fear.
Now, this is just in America.
I'm sure the stats are different in other countries.
You should go try and find out.
So what does it mean?
It basically means that we need to be responsible when we're building this AI technology.
While AI can actually mean better markets, responsible AI can actually mean that it actually increases the marketability.
Because now people feel safer rather than being worried that, hey, this piece of technology might do something dangerous that I don't know or understand.
It might actually increase the product adoption.
People might actually want to buy that product, because they feel safe with it.
They actually trust that particular device.



But you and I, as engineers and scientists and researchers, we need to be doing this proactively in order to convince the people that are out there doing their job that we have done our job.
So this is not a new realization.
Companies have been coming together-- major organizations like Amazon, IBM, Facebook, DeepMind, Google, Microsoft.
Everybody's been coming together in order to join forces, in order to build a partnership on the AI where the key principle is to understand what the risks are proactively and then be able to build intelligent products that are safe for all of us so that we can all benefit and society as a whole can benefit from this technology.
The reason I'm putting these different companies up here is because I want you to understand that when you have these different organizations coming together, they are communicating a critical message to you and me.
They are saying that they value these.
These are their first-order principles.
In doing so, they are trying to say that they want the next generation of leaders-- you-- to actually know these kinds of principles.
Because when you go over to work with them and build the next TinyML product, they're going to expect that you have learned the value of responsible deployments.
So let's go further.
If these companies all come together-- by these companies, I don't mean just this short list that I've listed here.
Just about any company that's using machine learning and AI in the field in order to enable products, if they come together, they can do well.
And by doing good, you can effectively improve your capabilities.
Now, there are a lot of unique features that TinyML obviously brings.
These are unique opportunities, but they also present unique challenges.
Some of the challenges we already talked about, like battery-powered, on-device ML running constantly and so forth.
Now, there are issues that are not purely technical in this space that are just as important.
And I want to talk about these things from the perspective of three important use cases or deployment scenarios that we're going to see for TinyML.
We're going to understand, obviously, the technological benefit.
And then we're going to see what the cons of deploying TinyML could be so that you and I can proactively design systems where we can enable the adoption of TinyML much more robustly and safely.




Let's say, for instance, we focus on industry.
Or we can talk about the environment as I'm giving you some examples.
Or we can talk about how TinyML actually influences our human behavior.
Now, if we look at the industry landscape, can we question, is responsible AI actually applicable?
Well, the truth is that it's actually applicable in industry environment and humans.
Specifically, if you look at industry, for instance, there are a lot of benefits of having TinyML being deployed in the industry, because it reduces downtime for repairs.
How so?
Well, there are things like anomaly detection mechanisms that we will learn about where, for instance, by detecting sounds or behavioral patterns of a machine, if something is about to go wrong, these machines give us little hints and cues in the vibrations they make or the sounds that they're emitting.
And by observing these things in real time right at the end point through TinyML, we can figure out if this motor or if this pump is actually going to break down.
And by knowing this proactively, what we can do in the industry is go ahead and get someone out to replace it.
Because imagine what the other consequence is.
The other consequence would be that that machine breaks down.
Then someone picks up a phone and says, hey, this machine is broken down.
Would you please come fix it?
All that time, the entire manufacturing pipeline might be shut down waiting for this particular product to be repaired.
Instead, if someone proactively knew that it needed repair, they could go ahead and fix it.
And the downtime would be much shorter.
You can also improve the efficiency effectively of that, right?
And it's also much more cost effective.


Why?
Because you're not suffering from a long downtime.
Your downtime shrinks, which means you're staying much more productive.
These all look great.
But what about the negative pitfalls?
The pitfalls could be, well, how accurate are your predictions that this system is actually going to fail?
What happens if it's a jet engine that's going to fail that costs millions and millions of dollars?
Are you just going to arbitrarily replace it just because your jet engine sensors seem to indicate that something was wrong?
Well, if so, then you have to worry about the accuracy.
What about explainability?
Who's responsible if it's a false alarm?
Is the operator responsible?
Or is the machine learning smarts that we package together actually responsible?
And overall, having this sort of predictive maintenance thing, how is it going to affect the workforce?
For instance, are people going to lose their jobs?
And is that necessarily a good thing for society?
Or are we actually going to enable new jobs?
So these are the kinds of questions that we need to think about in industry, for instance.
Let's talk about the environment.
I gave you example with ElephantEdge, where the goal was to kind of save these elephants from extinction.
Well, what are the benefits of having TinyML, the smart collars around them?
Well, you get these very detailed insights, right?
Often, when you want to learn something about these creatures that are out in the wild, it's like finding a needle in a haystack.
OK, well, if you have TinyML going around with these animals, then you're going to get very precise, detailed insights.
There's less wasted data, because it's not like you're collecting all this data, writing out all these logbooks,
and then analyzing them while something might have already changed in their environment.
So there's a long lead time before you figure something out and go back.
Well, with TinyML, you're talking about real-time responses.
So you actually waste less data.
And it's cost effective, because you don't need someone constantly walking behind the elephants in order to figure out what's going on with them.
You could effectively be deploying these collars out onto the field.
And then humans are safe, and they don't need to be venturing out into dangerous places.
And you can also overcome the limitations of human labor, just the capability that humans have in terms of going out and working year-round in forests and so forth.
So there are a lot of perks about doing it in the environment.
But what are the downsides?
Is the data that you're getting actually reliable?
Are you sure?
As we will look at this very closely later on in the Fundamentals of TinyML course, we will be looking very closely at, how good is the data?
How sure are you about the data you're collecting?
Because machine learning, the notion of influencing patterns from data that you're looking at, is a stochastic thing.
There's no guarantee that it's always right or it's always wrong.
So how can you be 100\% reliable?


Are you sure for the question you're asking that the data that the machine learning algorithms are producing is actually correct?
Who has access to that data?
If we're not careful about securing the data, someone illegal might actually get the data.
Worse, imagine if we're trying to avoid the high-risk conditions that we talked about previously, where we don't want to get these animals into some dangerous poaching areas, right?
And what happens if the data that is getting collected that they're reaching the poaching areas actually lands in the hands of the poachers?
Well, then they get smarter.
That would not be helping us while we're deploying all those technologies.
So we got to think about what the security implications are.
And that goes to the last point that I'm talking about.
Can these devices actually be hacked by poachers?
Even if we put mechanisms in, if they are hacked, then what ends up happening?
So we could actually make things really worse for these animals if we're not careful, proactively thinking about these things.
Let's talk about humans and how TinyML has great effect on us.
Well, TinyML makes things more accessible.
There are many people who are suffering from disabilities who can't touch type and do all those things, many of which you and I might be taking for granted.
Well, if TinyML, for instance, is all about speech-- like when we say, OK, Google, or, hey, Siri, these machines wake up.
Well, I don't need to touch anything.
That's great.


That makes it much more accessible for everybody.
Everybody should have access to technology.
And then it also makes the design of these products much more intuitive.
We don't need buttons to push here and there, because that has implications on how you design.
If it's all integrated and smooth, then it leads for a much better UI.
We often communicate as humans.
We don't necessarily go push someone's shoulder or press someone's head.
We don't do that.
We do that because we're interacting with machines.
But if TinyML can become pervasive, then it can enable us to build more intuitive, human-like experience, because that's what, ultimately, we're all trying to build technology for, helping us have a better experience in life.
And effectively, that would streamline our design.
We would end up with a much smoother sort of experience.
Bottom line, if TinyML is great, machines would be supplementing and working for us rather than us trying to understand how the machine works.
So those are just some of the examples.
But what are some of the downsides?
Well, you've got to worry about, is this going to work?
Is this technology going to work uniformly across all different kinds of populations?
One of the big issues that you will find in machine learning that we will touch upon quite a bit is actually diversity issues.
Machine learning models rely on data.
So if we're not getting data from different countries, then what are we going to do?
Because if the data is not there, then these algorithms are not going to learn on that data.
If they don't see it, they don't know it.
So what you don't know is what you don't know.
And that's going to become a serious problem, especially
if you're using machine learning to make intelligent decisions for you.
And does it threaten the user's privacy?
If this TinyML is always on, which I've kind of emphasized, then
what is it learning about us?
What is it learning all the time, right?
What if something is being captured that you don't want to be captured?
But the machine thinks it's doing good by trying to help you.
But if someone gets into that data, then you've got a serious problem.
So how can you protect that user's data?
These might all seem like simple questions,
but they're actually fundamentally deep.
And they need serious technical chops in order to be able to solve.
And I'm talking about these things with you
because they are extremely important.
The consequences of not being careful can be disastrous.
Why do I say this?
Here are just some high-level examples of things
that I want you to just Google and try and find on the web.
Microsoft's disastrous Tay experiment shows the hidden dangers of AI.
This was a bot that was basically trained on data.
I'm not going to say what it did.
I want you to go and do your homework and look up what it actually did.
There was a microphone that was hidden inside Nest.
I shouldn't say "hidden."
Rather, Google says that it was there as future-proofing the device.
But no one mentioned it in the manuals that there was
a microphone inside the Nest device.
So this raised a big question about, are they listening to us?
They never told us that there's a microphone in there.
So this leads to a trust issue, which really
can backfire on the organization.
And your product falls apart because of that sometimes.
And so you need to be careful.
And algorithms can also be unfortunately rather discriminative because
of the data that they're seeing.
If the data that's going in is biased, then they're
naturally going to be biased because of the way they're being taught.
It's no different than you and me learning about something
from a textbook.
Because if the textbook is only giving us one viewpoint, then that's all we're
going to learn, right?
Unless we get exposure to other textbooks that
give us a balanced portfolio, we would not
have a broader, wiser wisdom sort of experience.
And these are very similar things for these AI or machine learning systems.
And we've got to think about all these things.
Today, what we do is largely we find these problems,
and then we put a giant Band-Aid on top of them.
And then we say, OK, great, we fixed it.
Not a problem.

% >>>>>> hier weiter

But this is not how we can go about it, especially when
we're going to be deploying TinyML into billions of devices all around us.
There are going to be too many mistakes that we're going to make.
So what's the point here?
The point is that we need to proactively think about things.
We need human-centered AI or human-centered machine learning.
We need to proactively, systematically be thinking about these things
while we're building the systems.
We need to think about the design of these things.
And what are the consequences of the design?
When we are developing the algorithms, how
should we be thinking about the development of the algorithm
so that we're conscious of the potential issues that we might run into?
And once we throw these devices out into the field
and when they're operating out there, how
do we consciously think about protecting the data that they're sensing?
And what do we do with the mechanisms that
need to be in place before they're put out there?
So design, development, and deployment-- these are three critical stages.
Each one of these, we will actually be going through in a lot more detail.
And the core nugget of them is that we want to keep human values in the loop
at all stages of developing machine learning for these systems.
So we will be talking about design, development,
and deployment.

It's going to be centered around when we design the machine learning models.
What are the things that we need to think about?
What are the characteristics of the data, for instance?
When we are starting to train the models for very specific tasks as we develop them, what are the nuances that you need to be worried about so that you're making sure you're teaching them the right things?
Third is when you deploy them.

You need to be careful about exactly what it means to put them out
in the field and be responsible for the kind of data and how the data is secure
and so forth.
And we'll be touching upon all these things.
So design is going to talk about, what am I building?
Who am I building it for?
And what are the consequences if it fails?
It's philosophically understanding the issues.
Development is, what's the data?
Is the data set biased?
How is the model going to react?
Is it going to be fair for everybody?
Deployment-- when I put this out there, is the model
going to drift if the data is kind of shifting,
if the characteristics of your ecosystem are changing?
Is it actually going to adjust to it, or is just
going to keep doing the same thing that it was taught to do?
And how do you deal with security, which is a very big critical piece?
And how do you preserve a user's privacy?




Don't forget that when we build technology, especially advanced technology that is very, very pervasive and going to be ubiquitous, that we do so in a highly responsible manner.



\section{Forum: Discuss AI/ML failures}



The following three case studies are all real-world examples of when AI has failed in some way. First, read through the following descriptions of each case:

  \begin{enumerate}
    \item  Winterlight Labs auditory detection of Alzheimer’s disease

In 2016, Winterlight labs designed an AI-powered auditory test for Alzheimer’s disease, where users’ speech would be recorded and AI would be used to detect signs of Alzheimer’s such as vocabulary richness, pauses in speech, and syntactic complexity. However, the initial research findings revealed a serious problem; non-native English speakers were being inaccurately flagged as having Alzheimer’s disease. Since the data that was used to train the model had been collected from native English speakers from Ontario, Canada, this technology was unable to work reliably across different populations. 

\item  Wireless baby monitors hacked

In 2018, there were several instances where wireless baby monitors were hacked which ultimately made national news headlines. In one case, a hacker used his newfound access to the baby monitor device to broadcast threats and shout sexual expletives. In another case, a more benevolent hacker used his newfound access to warn parents about the susceptibility of their device, in hopes that the parents would be able to address the situation before being targeted by nefarious hackers. 

\item  Hidden microphones in Nest devices

In 2019, users of Nest Guard devices were shocked to discover hidden microphones inside the device. Users were concerned about the invasion of privacy as well as the breach of trust that resulted from not being properly informed about the specs of the device. From Google’s perspective, the on-device microphone was simply a form of future-proofing that would allow the device to be compatible with updates that supported new functions later down the line.

\end{enumerate}

Now, reflect on the following questions and join the discussion going on in the forum!

\begin{itemize}
  \item [A.] Which one of these cases do you find most concerning? Which concerns you the least?
  \item [B.] What do you consider to be the relevant ethical challenges? 
  \item [C.] What do you think the designers of this technology could have done differently?
  \item [D.] How can you apply learnings from these examples to your own job? Your personal life?
  \item [E.] Do you agree or disagree with what others have already posted in the forum? Why?
\end{itemize}

%<<<<<<<<<<<<<<<< bis hier

% 1.4 Getting Started

% ab hier in die Datei GoogleColab überführt
%>>>>>>>>>>>>>>>>>>>>>>>>>>>>>
\section{What resources will be needed for each course?}

VIJAY JANAPA REDDI: Hi In this video, I want
to talk a little bit about the resources you need.
So you're well prepared for the entire program.
Starting with course 1, where we will be focused
on the fundamentals of machine learning, what you really need are quite simple.
You just need an internet connection in order to be able to watch your videos,
and you need a Chrome browser.
The Chrome browser is a specific requirement primarily because we will
be using a programming environment called Google Colab which is going
to rely on working well with Chrome.
So if you use Chrome, you will likely not have too many issues.
And so we strongly encourage you to use that browser.
Then, as we move on into the applications of TinyML,
we will need a little bit more than what we needed in course 1.
Specifically, we will need your laptop or your desktop
to have a microphone and a laptop camera.
Why do we need that?
Don't worry.
We're not trying to listen in on you, like we did
when we talked about responsible AI.
Why we need those devices are primarily because we want to engage with you.
We want to be able to get you to actually train
you own special word that wakes the machine up, for instance.
And in order to be able to do that, you're
going to want to collect a bit of data from the microphone
by saying different words, understanding how to say it in different ways,
so that your model is general and flexible enough,
or capture different kinds of images.
So we want you to be able to have access to a camera and a microphone.
Then, as we go into course 3, things will get even more fun,
because we'll be giving you a special kit.
And the kit needs you to physically interface the microcontroller
board with your laptop.
And so in addition to the things you need from course 2,
in course 3 we will be adding a couple of extra parts.
And together, you will be effectively building
on all the things you've learned as you step through the courses.
And with that, I want you to have a really good time,
and I want you to be prepared for this course.
So please make sure you have all the resources you need,
and I will see you in the next video.

\section{Video Introduction to Colab}

Continuing our preparation for the course, here is an introduction to Google Colab. Google Colab is an online integrated developer environment to design, train, and test our machine learning models. \href{https://www.youtube.com/watch?v=inN8seMm7UI}{Watch Jake VanderPlas from Google give a wonderful intro to Colab.}

\section{Colab in this course}

VIJAY JANAPA REDDI: Welcome back.
I'm glad you had an introduction to Colab.
Now in this video, what I want to talk a little bit about
is how we will be using Colab in order for learning tiny machine learning.
Well, first and foremost, we will be doing
a code walkthrough of algorithms and different kinds of concepts
using Colab.
For example, we will be looking at a small snippet
of algorithmic code for the machine learning model,
and we will learn how to tweak it in Colab.
And then, there will be quizzes and assessments
that will help you kind of internalize those critical concepts by having
you fill in some of the blanks that are out there.
In addition, because we will be using the TensorFlow framework through Colab,
one of the critical things that we will be doing
is we will be using the power of the cloud
in order to train our neural network sufficiently.
If you tried to teach the machine learning models
to look for interesting data patterns manually on your laptop is what exactly
I mean.
If you try and do that, there will be cases where you
might have to wait for a month or more.
Instead, if we use the capabilities that Google Colab gives us in order
to train our networks much faster, then we
will be able to accelerate the whole learning
process by using GPUs for instance.
And that's going to save you a lot of time.
Trust me, there are assignments in this program that will sometimes need
about 13 hours or so in order to train.
But don't worry about that.
We won't have you sitting in front of the screen for 13 hours.
We have found ways to kind of improve that experience.
Another way that we will be using Google Colab
is to help you analyze the way that the algorithms are actually learning.
So we will be visualizing interesting characteristics
about the machine learning training that is actually happening.
And in doing so, you will understand what
are some of the critical implications of the way you set some hyper
parameters, like learning rate, for instance,
or what's the impact of the loss function?
These are things that we will be getting into in the next chapter.
So with that said, I hope you got a little bit of information
about Google Colab.
And once you start using it, I promise you,
you will feel much more comfortable.
And I'll look forward to seeing you soon.

\section{Learning Colab}

As mentioned earlier, we’ll be using the Colab environment for much of our programming in this specialization. In order to ensure that everyone has some hands on experience with Colab, please follow the link below to Google’s Introductory Colab. In it they provide a handful of quick tips and hands on exercises to get you started using the Colab environment.

\url{https://colab.research.google.com/notebooks/intro.ipynb}


\section{Tips for using Colab in this course}


Now that you’ve heard about Colab and gotten a chance to explore a simple Colab we thought we’d give you a couple of tips and tricks for using Colab in this course:

\begin{enumerate}
  \item Each Colab instance will stay alive as long as you interact with it once every 90 minutes.
  \item Each instance will only last for a maximum of 12 hours so make sure to download all of the files you need ASAP to avoid losing them. We will discuss particular files of note as we start to use longer Colabs later in the specialization.
  \item You can access the files in the the screen by clicking the folder icon: 

\begin{figure}
    \GRAPHICSC{1.0}{1.0}{TinyML/1 Fundamentals/Colab1}
    \caption{This image is a screenshot of the Colab UI showing one portion of the interface.  In this portion, there is a folder icon highlighted in yellow, which you can use to get to the files you need for this course.}
\end{figure}  




  \item By default all of your files are saved under /content and by default you start only seeing that folder. 
  \item Remember, while the code cells in Colab default to python, any bash command can be executed by prepending the command with \PYTHON{!}. For example you can unzip a file with \SHELL{!unzip file.zip -d destination\_folder}
  \item To download a file from Colab simply click the three dots next to the file's name and then select download
  \item To upload a file to Colab simply click the icon with the piece of paper with a folded corner and an arrow on it below the word Files
  \item To make permanent changes to a Colab you need to select file->Save a copy in Drive which will make a copy of the Colab in your drive with your changes saved.
  \item You can drastically speed up Neural Network training by activating a GPU instance. This can be done through runtime->Change runtime type and then selecting GPU under Hardware accelerator
  
\begin{figure}
    \GRAPHICSC{1.0}{1.0}{TinyML/1 Fundamentals/Colab2}
    \caption{selecting GPU under Hardware accelerator}
\end{figure}  
  
\end{enumerate}



% bis hier in die Datei GoogleColab überführt
%<<<<<<<<<<<<<<<<<<<<<<<<<<<<


% ab hier in die Datei TensorFlow überführt
%>>>>>>>>>>>>>>>>>>>>>>>>>>>>>
\section{Introduction to TensorFlow}

For machine learning in this course we are going to use Google’s TensorFlow Machine Learning framework. TensorFlow is widely used across industry and academia. \href{https://www.youtube.com/watch?v=yjprpOoH5c8}{To get you started with TensorFlow here is Google’s wonderful intro to TensorFlow video}.

\section{TensorFlow for this course}

VIJAY JANAPA REDDI: Welcome back from the video on TensorFlow.
Hopefully, what you have learned from the video
is that you have understood that TensorFlow is a critical software
piece of infrastructure that is needed, in order
to be able to train a machine learning model,
in order to be able to teach an algorithm,
to look for interesting patterns in data, and how we go about doing that.
We will be training models, specifically in the context of TinyML.
But as you might recall, from the earlier videos,
I told you the TensorFlow is big, bulky, and beefy.
Why?
Because it was originally designed for big computing systems.
But as we move into deploying things on small, little embedded
microcontrollers, what we need to do is to think about how
we make the system lean and mean.
And that's why TensorFlow Lite was developed.
TensorFlow Lite goes through a series of steps in taking that big model
and converting it down into a format that
is much smaller, much more compact, and much more
suitable for the endpoint devices.
In addition, it does optimizations like the ones
that I hinted, talking about pruning, and quantization, and so forth.
And then, finally, sets up the model, so that it's effectively
deployable on the edge, where you can actually do inferences.
Now, all of that is extremely important.
In fact, I would argue that if you want to learn about TinyML,
you should learn about how to train the model.
But of course, you also got to learn how to deploy
the model under the specific constraints of an embedded microcontroller.
So we will be stepping through how to use this TensorFlow stack properly,
the architecture, in order to train a model
and target a very specific microcontroller.
But one of the key things I'm going to make sure that I enforce to you
is that the lessons that you learn, as we
step through this process [INAUDIBLE] generalizable
so that you can actually change the microcontroller that we
might be looking at.
And then, still, all the things you have learned, in going from training,
to deployment, onto the microcontrollers is totally valid.
With that said, I hope you have gotten all the tools
that you need to know about.
And if there's any questions or any doubts,
feel free to just post them on the forum.
And I'm sure some one of us will be able to reply.
Or the vast community of your colleagues and students,
who are taking this class, will actually be able to help you out.
I'll see you soon.

\section{Sample TensorFlow code}

While TensorFlow is written with fast custom C++ code under the hood, it has a high level Python interface that we will use throughout this course. As such, you’ll need to be comfortable with basic coding in Python for this course.

For example we might create a custom Neural Network model using code that looks something like the below. Don’t worry if you don’t know what all of the words mean, we’ll walk you through all of these topics later in the course. For now just make sure you know enough Python to understand the following high level takeaways:


\begin{enumerate}
  \item This is a custom class that extends the \PYTHON{tf.keras.Model} base class
  \item It defines four specific kinds of \PYTHON{tf.keras.layers}
  \item It has a \PYTHON{call} function which will run an input through those layers (functions)
\end{enumerate}


If that resonates with you then you probably know enough Python to do well in this course!

If you don't feel like you know enough Python yet, don't worry. Both \url{Python.org} and \url{LearnPython.org} have great resources you can use to get you up to speed. We're sure that with a little work, very soon you’ll be ready to proceed with this course.

\begin{code}
    \begin{lstlisting}[language=Python]
 # Load in the TensorFlow library
 import tensorflow as tf
        
 # Define my custom Neural Network model
 class MyModel(tf.keras.Model):
     def __init__(self):
         super(MyModel, self).__init__()
         # define Neural Network layer types
         self.conv = tf.keras.layers.Conv2D(32, 3, activation='relu')
         self.flatten = tf.keras.layers.Flatten()
         self.dense1 = tf.keras.layers.Dense(128, activation='relu')
         self.dense2 = tf.keras.layers.Dense(10)
        
     # run my Neural Network model by evaluating
     # each layer on my input data
     def call(self, x):
         x = self.conv(x)
         x = self.flatten(x)
         x = self.dense1(x)
         x = self.dense2(x)
         return x
        
        
# Create an instance of the model
model = MyModel()
    \end{lstlisting}
    \caption{Simple Example for TensorFlow}
\end{code}

\section{Objectives}

The goal of this section of the course is to introduce students to “ML language” that we will be using throughout the TinyML course. More specifically, this course focuses on the basics of machine learning and embedded systems.

Some of you might have a deep machine learning background. In this case, parts of this section will be more of a review for you. We encourage you to still go through things, take the tests, and get excited for the next course which will dive deeper into TinyML, covering frameworks and applications that you may be less familiar with.

\section{Prerequisites}

\begin{itemize}
  \item Basic Scripting in Python
  \item Basic Usage of Colab (covered in Chapter 1)
\end{itemize}


% bis hier in die Datei TensorFlow überführt
%<<<<<<<<<<<<<<<<<<<<<<<<<<

% 2.1 The Machine Learning Paradigm

\section{The Machine Learning Paradigm}

LAURENCE MORONEY: Hi, I'm Laurence Moroney,
and I lead AI advocacy at Google, working
on a product called TensorFlow that you'll hear all about in this course.
But before you get into the nuts and bolts of using machine learning
and programming with it, I just want to explore what machine learning actually
is on a fundamental level.
I personally like to see it as a new paradigm for programming computers,
where instead of just coding a solution, you can teach
a computer how to learn a solution.
This can be incredibly useful, and it can open up new scenarios for you
to solve problems that can't be solved with programming alone.
Let's dive in.
Explicit coding is the bread and butter of software developers
all over the world.
You think about the scenario that you're working on
or the problem that you're trying to solve,
and you express the solution to that by defining the rules that
determine the behavior of a program.
For example, in this bricks game, the ball moves along a path, which can get
changed when it hits a brick.
As a programmer, you have to figure out the behavior.
How does the ball bounce?
Which brick do I remove?
What happens when the ball hits a wall?
Or what happens when the ball misses the bat?
Everything needs to be figured out in advance.
It's predetermined by the programmer, and it's then coded and tested.
The more complex the scenario, the more code you'll have to figure out.
And with the law of diminishing returns, often you'll
encounter scenarios that it just isn't feasible to figure out
all the code for.
We'll explore one of those in a moment.
But first, this diagram shows coding scenarios like the one
we just explored.
Traditional programming is solving a problem by figuring out and expressing
roles using a programming language like Java, C++, or Python.
You then have these act on data, like the location of the ball
and the bricks in the previous game, and then you get answers from that.
Is it game over?
Does the score go up?
Where should the ball be plotted next?
It's really powerful, but it can be limited.
So let's explore one scenario where this doesn't work well at all.
Consider activity detection.
Say you want to write an app that uses sensors
on a phone or a watch or something else to determine a person's activity.
You could, for example, use the data about their speed
and write a rule that determines if the speed is below a certain amount,
then they're probably walking.
You have data.
You have a rule.
And then you can get an answer.
And then say you extend this to determine if they're
running by building on that rule.
If, for example, it's below 4 miles an hour, then they're walking.
Otherwise, they're running.
OK, great.
That still works.
And then maybe you can extend that even further to see if they're biking.
If the speed is below 4, they're walking.
Otherwise, if it's below 12, they're running.
Otherwise, we can say they're biking.
But then how would you handle golfing?
What rule could you write that determines
that they're actually playing golf?
Also, by now, you've probably realized that the other rules are also
a little bit naive.
You can't just go by speed.
You might run downhill faster than your bike uphill, for example.
So let's go back to our diagram.
Recall that we mentioned that traditional programming is
when you create rules in a programming language and these rules act on data
and provide you with answers.
With activity detection, we had created rules
to figure out the user's activity, be it walking, biking, or running,
but then we hit a wall with golfing.
Not only that, we can also see that our algorithm was that little bit naive.
Our rules don't really work that well.
For example, you might run downhill faster than you bike uphill.
And as a result, your app might fail when written with rules like this.
Machine learning can help solve this problem.
And it can be represented with a simple rearrangement of the diagram.
So instead of you trying to figure out the rules that act on the data
to give you an answer, what if you go the other way around
and you provide the answers with the data
and you have a computer that can figure out the rules that
will match them together?
Once it's done this, it can then apply those rules to future data
in order to figure out the answers about it.
So in our activity detection scenario, we could gather a lot of data
from sensors and label them with the activity that the user is following.
So now, by matching parts of the data with the label that
describes the activity, the computer might
be able to figure out the rules for what makes an activity walking or running
or biking or, yes, even golfing.
For example, I've highlighted four bits within each of the binary blocks,
and you might notice that these four bits are
different between the four activities.
Now, that could be a signal for which activity is which.
With machine learning, we have some APIs that
allow us to write code to potentially identify such distinguishing factors.
And you'll see that a little later in this course.
But first, let's consider at a very high level how this works.
First, the computer will make a guess as to the relationship
between the data and its labels.
It does this by randomly initializing a neural network.
And you'll delve into that a little later.
But beyond the specifics, it's simply guessing as to the relationship.
It then measures how good or how bad that guess is.
The terminology often used here is loss.
Higher loss is roughly analogous to lower accuracy.
So then you can measure the results of your guess,
and you can then use the data from the accuracy measurement
to figure out your next guess, optimizing
based on what you already know.
And if you then repeat the process with the logic being
that each subsequent guess gets better than the previous one,
then your model becomes more and more accurate.
It's learning from experience what the best guess might be.
So let's return to the diagram describing
machine learning from a high level.
I've changed the word "answers" to "labels" here,
and that's generally used for the term that describes your data.
But the rest still stands.
You start with label data and go through the process I previously mentioned.
And you end up with a set of rules that matches that data to those labels.
This constructs a model, and from this model, you can then give it new data,
and it will figure out how closely that data fits the set of labels.
And it will return you something called a set of inferences,
and these are the probabilities that the data matches each specific label.
OK, so now it's your turn.
In the next section, I'm going to give you a set of numbers.
And there's a relationship that matches the x to the y, a rule
that you can figure out.
Give it a try, and think about how you figured it out.

\section{Finding patterns}

In the previous video you learned about how traditional programming is where you explicitly figure out the rules that act on some data to give an answer like this:

\begin{figure}
  \includegraphics[width=\textwidth]{TinyML/1 Fundamentals/TraditionalProgramming}
    \caption{Rules and data feed into traditional programming which then outputs answers.}
    \Mynote{Besser mit tikz}
\end{figure}  



And then you saw that Machine Learning changes this, for scenarios where you may not be able to figure out the rules feasibly, and instead have a computer figure out what they are. That made the diagram look like this:

\begin{figure}
  \includegraphics[width=\textwidth]{TinyML/1 Fundamentals/MachineLearning}
    \caption{Answers and data feed into machine learning and from this rules emerge.}
    \Mynote{besser mi9t tikz}
\end{figure}  



Then you read about the steps that a computer takes -- where it makes a guess, then looks at the data to figure out how accurate the guess was, and then makes another guess and so on.

So, consider if I give you a set of numbers like this:

$$X: -1, 0, 1, 2, 3, 4$$

And then I give you another set of numbers like this:

$$Y: -3, -1, 1, 3, 5, 7$$

Can you figure out the relationship between the two sets? There’s a function that converts -1 to -3, 0 to -1, 1 to 1, 2 to 3, 3 to 5 and 4 to 7. Can you figure out the relationship. Think about it for a moment.

Often when I ask people about it, they see that the 0 is matched to -1, so Y is (something) times X - 1. Maybe they’ll take a guess at the something, and come up with 3.

Then fill in the gaps, if $Y=3X-1$, then

$$X: -1, 0, 1, 2, 3, 4$$

Becomes

$$Y: -4, -1, 2, 5, 8, 11$$

Other than working for 0, it fails for everything else. In ML terms, you can define this as your loss is ``high''.  With what you learned from that, you might think, what if it's $Y=2X-1$?

Then, when you fill in the results for $Y=2X-1$, you'll get:

$$X: -1, 0, 1, 2, 3, 4$$

Becomes

$$Y: -3, -1, 1, 3, 5, 7$$

\ldots which matches your original data perfectly. Your loss is zero.

You've just gone through this process:

\begin{figure}
  \includegraphics[width=\textwidth]{TinyML/1 Fundamentals/MachineLearningLoop}
    \caption{Making a guess leads to measuring your accuracy which leads to optimizing your guess. This then gets repeated where you make a guess again and repeat the process. Arrows connect the steps.
    }
\Mynote{Besser mit tikz}
\end{figure}  


In the next video you'll explore this in a little more depth.

\section{Fitting lines - Thinking about loss\ldots}

LAURENCE MORONEY: Earlier, you looked at a process
that defines machine learning, where you make a guess.
You then measure how good or how bad that guess is.
And then, based on that, you can make another guess.
And keep repeating the process.
Measuring how effective your guess is is often called measuring your loss.
Minimizing your loss means you increase your accuracy.
So you'll learn different strategies for minimizing the loss soon.
But first, let's take a look at it and a way that you can measure it.
So for example, you can have the values for x and y that you can see here.
And you want to figure out a way to match x to y.
You know that y is a function of x.
You just don't know what that function is.
So you make a guess.
Say you think it's y equals 3x minus 1.
It's as good a guess as any.
But how do you measure how good or how bad it is?
Well, first, we can calculate what y would look like.
And you can see that here.
Now that we know what y is supposed to look like, I can see my y.
And I can see the real y.
They're not the same.
In some cases, they're pretty good.
The second entry, for example, looks OK to me.
And in some cases, they're really bad.
And the last entry, where I guessed it would end up 11.
But the real value is 7.
Now, is there a way for me to formulize a calculating how good
or how bad these are?
Well, maybe I could look for it by plotting them.
And here, you can see the real values of y.
And you can see my guesses.
And by drawing a line between them, I can, by the length of that line,
figure out which ones are good and which ones are not so good.
I can measure the length of these lines.
So it's minus 1, 0, 1, 2, 3 and 4 for my six points.
So if I now remove the chart, I can see these clearly.
So maybe, I could add them up to get the overall loss from my guess.
But that seems wrong.
Because even though both of these are wrong, they would actually cancel out.
And it would look like I actually have three of my six correct,
when in fact, I only have one.
So what if I square them to remove any negatives?
And now, with my losses squared, it looks like this.
They're all positive.
So I won't have any canceling each other out.
And my correct one is still 0.
So now, if I total up all of these, I'll get 31.
Because I squared them, I can take the square root
to get a fair measurement of my error.
And as you can see, it's 5.57.
So now, let me repeat that with another guess.
What if I guessed that y equals 2x minus 2?
Now that I can see that my y is still off the real y,
but when I measure the difference and I square that, I get five ones.
If I add them up, I'll get 5.
And the square root of that as 2.23.
So my loss is now smaller than it was.
And I know I'm heading in the right direction.
And if I now guess y equals 2x minus 1, I can calculate my loss as 0.
And now, I know I have the right formula to get the right answer.
This video just explained loss and how to measure it.
Now, of course, the next step should be, how can you
use that information to generate the next guess,
knowing that the next guess should actually
be a better one than the current one?
That process, called optimization, will be covered a little later.
But first, let's take a look at some code that uses this process.

\section{Coding exercise}

In this next Colab you will get to explore linear regression and loss functions for yourself. We have provided a cell that will compute the predicted Y values as well as the loss function for your guess of w and b. Simply change their values and explore how the output and loss changes.

Note: You may receive a warning that the Colab was not authored by Google, and this may occur on future Colabs as well. That is entirely normal as we authored many of these Colabs ourselves for this course

\url{https://colab.research.google.com/github/tinyMLx/colabs/blob/master/2-1-4-ExploringLoss.ipynb}


\section{Gradient descent}


LAURENCE MORONEY: Earlier we saw how the training loop works.
You make a guess.
You measure how accurate that guess is by calculating its loss.
And then you use that data to make another guess
with the next guess usually being a little bit better than the current one.
The idea is that the smaller the loss, the more accurate your guess is.
So in this video, we'll explore how you can minimize your loss.
So let's first look at how gradients and derivatives can help us
understand how to minimize loss and, in other words, how
an optimizer function works.
Let's plot our loss function.
You'll remember that we were squaring values to remove the negatives,
so our function is a square of something.
Square functions have a parabolic shape like this one.
If we want to know the minimum of the loss function,
we just look to the bottom of the parabola.
It doesn't matter whatever the parameters that make up the function
are and wherever this would be plotted on a chart.
We will still know that the minimum is at the bottom.
So if we made a guess as to what the parameters of our equation would be--
say we had started with y equals 10x plus 10, for example--
and we calculated the loss for the data in our equation,
we know that it would be very high and so very far away from the minimum.
Maybe it would look a little bit like this.
Now it gets interesting.
And here's why you generally see a lot of math and calculus
in machine learning.
If you differentiate the values with respect to the loss function,
you'll get a gradient.
The nice thing about having a gradient is
that you can use it to determine the direction towards the bottom.
It's actually the negative of the gradient points that way.
You have no idea how far it is to the bottom,
but at least you know the correct direction.
It's like a ball on a slope.
Gravity will give you the direction of the bottom of the slope,
but you don't know how far it is to go.
So if you want to go towards the bottom, you can take a step in that direction.
You know what the direction is, and you can pick a step size.
The step size is often called the learning rate.
So given a direction from the gradient and a step size from the learning rate,
you can now safely take a step towards the minimum.
And you'll end up here.
You're closer to the bottom.
You still don't know how far it is, but you're closer.
So you can repeat the process.
By differentiating again, you can get the gradient.
And from the gradient, you can understand
the direction towards the bottom.
So you can move again in the direction of the gradient,
and you'll end up here.
Repeat the process again, and you could end up here.
You're getting closer.
Another step and you'll be here, closer still.
The next step could have us overshoot the bottom
and end up here because of our step size.
But that's OK, because when we differentiate from this spot,
we'll have a new gradient pointing back towards the bottom.
So we might end up here, overshooting again.
And after a few steps, we could find ourselves oscillating over and back.
But we're generally approaching the minimum.
It's important to choose a good learning rate or we could oscillate forever.
Or indeed, an advanced technique is to adjust the learning rate on the fly.
It might be good to have it bigger when we're further up the curve
and gradually reduce it step by step so that we can get to the bottom
more quickly.
To illustrate, choosing the learning rate wisely is important.
If it's too large, we could overshoot like this
and then overshoot again on the way back like this
and then overshoot again like this.
We're never actually reaching the minimum.
Or if the learning rate is too small, it might take a long time
to reach the minimum, a long, long, long time.
Tiny increments will get you there.
But as you can see, it's not the most efficient way.
OK, so now it's your turn.
You'll try a co-lab first that does this gradient descent using TensorFlow,
and then you'll see how to figure out the minimum loss.

\section{Coding exercise: Gradient descent}

In this next Colab you will get to minimize the loss function for our linear regression model using TensorFlow’s built-in GradientTape class. You will also get to visualize the training process, seeing how both the bias and weight converge on the optimal solution across training epochs.

\url{https://colab.research.google.com/github/tinyMLx/colabs/blob/master/2-1-6-MimimizingLoss.ipynb}


\section{Introductory neural networks - First Neural Network}

PROFESSOR: Previously, you saw how you can start fitting parameters
to an equation by making a guess about those parameters,
measuring your loss on that guess, and then
using what you knew to make a better guess by plotting
on a chart of the loss using calculus to get the gradient,
and then using that to figure out your way towards minimum loss.
This is ultimately how neural networks work under the hood.
So in this video you're going to switch gears and see
how to code a neural network that achieves the same goal.
Recall, here's the set of numbers.
We have a bunch of x's and a bunch of y's.
And we want to figure out the relationship between them.
Earlier, we used the process of machine learning-- guessing, measuring,
optimizing by hand, and we got the result y equals 2x minus 1.
Let's now see how a neural network would do this.
Here's the code.
There's a lot going on here that you may not be familiar with.
So I'll break it down line by line.
This first line defines our model, and it's the simplest neural network
that I could think of.
You've probably seen neural networks that look a little bit like this.
So that we can understand the terminology in the code,
let's explore some of that here.
Each of these circles will be considered a neuron.
Each of these sets of neurons is called a layer.
So for example, this would be one layer, this would be another,
and this would be another.
These layers are also in sequence.
The first layer feeds the second which feeds the third.
As a result, we'll use the term sequential to define this network.
You might also have heard the term hidden layers
when dealing with neural networks.
The hidden layers are the middle ones that are neither input nor output.
So that middle one, you can see here, is a hidden layer.
The neurons are all connected to each other, making a dense connection.
And we'll use that term to define the name and the type of the layers.
So now if we go back to our code, we can see the term sequential.
We're defining a sequential neural network.
Within the brackets, we would then list the layers in this neural network.
We only have one entry, so we only have one layer.
This layer is a dense, so that we know it's a densely connected network.
The units parameter tells us how many neurons will be in the layer.
And we can see that we only have one.
So this code defines the simplest possible neural network.
There's one layer and it has one neuron.
There's one other thing to call out when defining
the first, and in this case the only, layer in a network,
you have to tell at the input shape.
Now here our input shape has only one value.
We're training a neural network on single x's to predict single y's.
So again our input shape really couldn't be any simpler.
Visually, our neural network looks like this.
We only have one layer, and it has one neuron.
It takes an input, which is a single value we'll call x,
and it will learn the parameters to produce y based on our training data.
So now hopefully you can see that we're defining a new model, which
is a neural net comprised of a single layer with a single neuron
and that neural net will take a single input value.
When we compile the model, we'll define the loss function and an optimizer.
The loss function is mean squared error, as we used previously.
But instead of calculating everything ourselves,
we can just tell TensorFlow to use this last function for us.
The optimizer is SGD, which stands for Stochastic Gradient Descent.
And this will follow the process we saw in the previous video of using
a gradient to descend down the loss curve to try and reach a minimum.
We're not doing any of the math or the calculus ourselves
we'll just let TensorFlow do it for us.
Now we specify the x's and the y's.
TensorFlow expects these to be in a numpy array,
and that's what that np stands for.
And you can see that the array of y's has the corresponding y
for a value in the x.
So y equals minus 3 when x equals minus 1, y equals minus 1 when x equals 0,
and so on.
The process of figuring out this relationship or training
is called fitting.
So we do model.fit.
We're fitting the x's to the y's.
We'll do this for 500 epochs, where an epoch is each of the steps
that we had discussed previously-- make a guess, measure the loss of the guess,
optimize, and continue.
So this line of code will do that 500 times.
Then when we're done, we can use the model to predict the y for a given x.
So if x is 10.0, what do you think y will be?
Well, you can try it out for yourself now.
Maybe you think if x is 10, then y would be 19.
And the answer to that surprisingly is, it's not.
And can you figure out the reason?

\section{Coding exercise: neural networks}

In this next Colab you will get to explore training your first neural network—see just how easy that can be with TensorFlow. You’ll also see the result of the question posed at the end of the last video—will the predicted Y be 19 for X = 10? Launch the Colab to find out!

\url{https://colab.research.google.com/github/tinyMLx/colabs/blob/master/2-1-9-FirstNeuralNetwork.ipynb}


\section{More neural networks}

Up to now you’ve been looking at matching X values to Y values when there’s a linear relationship between them. So, for example, you matched the X values in this set [-1, 0 , 1, 2, 3, 4] to the Y values in this set [-3, -1, 1, 3, 5, 7] by figuring out the equation Y=2X-1.

You then saw how a very simple neural network with a single neuron within it could be used for this.

\begin{figure}
  \includegraphics[width=0.5\textwidth]{TinyML/1 Fundamentals/Neuron}
    \caption{An input leads into a neuron as an arrow. Out of this an output arrow is generated.}
    \Mynote{besser mit tikz}
\end{figure} 


This worked very well, because, in reality, what is referred to as a ‘neuron’ here is simply a function that has two learnable parameters, called a ‘weight’ and a ‘bias’, where, the output of the neuron will be:

\bigskip

Output = (Weight * Input) + Bias

\bigskip

So, for learning the linear relationship between our Xs and Ys, this maps perfectly, where we want the weight to be learned as ‘2’, and the bias as ‘-1’. In the code you saw this happening.

When multiple neurons work together in layers, the learned weights and biases across these layers can then have the effect of letting the neural network learn more complex patterns. You’ll learn more about how this works later in the course.

In your first Neural Network you saw neurons that were densely connected to each other, so you saw the Dense layer type. As well as neurons like this, there are also additional layer types in TensorFlow that you’ll encounter. Here’s just a few of them:

\begin{description}
  \item[Convolutional] layers contain filters that can be used to transform data. The values of these filters will be learned in the same way as the parameters in the Dense neuron you saw here. Thus, a network containing them can learn how to transform data effectively. This is especially useful in Computer Vision, which you’ll see later in this course. We’ll even use these convolutional layers that are typically used for vision models to do speech detection! Are you wondering how or why? Stay tuned!
  \item [Recurrent] layers learn about the relationships between pieces of data in a sequence. There are many types of recurrent layer, with a popular one called LSTM (Long, Short Term Memory), being particularly effective. Recurrent layers are useful for predicting sequence data (like the weather), or understanding text.
\end{description}

You'll also encounter layer types that don’t learn parameters themselves, but which can affect the other layers. These include layers like dropouts, which are used to reduce the density of connection between dense layers to make them more efficient, pooling which can be used to reduce the amount of data flowing through the network to remove unnecessary information, and lambda lambda layers that allow you to execute arbitrary code.

Your journey over the next few videos will primarily deal with the Dense network type, and you’ll start to explore how multiple layers can work together to infer the rules that match your data to your labels.

\section{Neural networks in action}

You've now seen a very simple example for how computers can learn. There's no great mystery to it -- it's a simple algorithm of making a guess, measuring how good that guess is (aka the loss), and then using this information to optimize the guess, and continually repeating this process to improve the guess.

What you've seen -- fitting numbers in an equation -- might seem trivial, but the methodology that you used to do this is the same as is used in far more sophisticated scenarios.

To understand just how powerful this simple method - Machine Learning - can be, lets now explore a couple of new and exciting case studies:

The first is the story of a young woman, Nazrini Siraji who used Machine Learning to detect diseases in crops, helping to stem the destruction of crops in her home country of Uganda: \url{https://www.youtube.com/watch?v=23Q7HciuVyM}

Next, is air cognizer, built by undergraduate students in India, who realized that pictures of the sky, when matched to labels on air pollution could be used to build a new type of air quality sensor, using just the cameras on their phones: \url{https://www.youtube.com/watch?v=9r2VVM4nfk8}

Finally, here’s a talk from a Google engineer about how Google used images of retinas to build a diabetic retinopathy detector using TensorFlow that performs state of the art diagnosis of this disease: \url{https://www.youtube.com/watch?v=oOeZ7IgEN4o}


\section{Coding assignment}

Welcome to your first assignment. Most assignments will be in Colab and will ask you to fill in some parts of the code to complete the Colab. Remember, if you are taking the course for a certificate during the end of section test you will be asked questions about the assignments.

This first assignment is a bit of a warm-up and simply has you fill in two variables to complete a simple linear regression example in TensorFlow. 

Note: Remember, you may receive a warning that the Colab was not authored by Google, and this may occur on future Colabs as well. That is entirely normal as we authored many of these Colabs ourselves for this course.

\url{https://colab.research.google.com/github/tinyMLx/colabs/blob/master/2-1-13-AssignmentQuestion.ipynb}

\section{Coding assignment solution}

We hope you enjoyed your first warm up assignment. We’ll be following a similar format throughout the course and will always also provide the solution following the assignment.

For this assignment we expected you to provide the following:

\begin{lstlisting}
SHAPE = [1]

LOSS = "mean_squared_error"
\end{lstlisting}

That said, for the \PYTHON{LOSS}, as we will explore later in the course, there are many loss \href{https://www.tensorflow.org/api_docs/python/tf/keras/losses}{functions included in TensorFlow}, and for this example many of them would have also resulted in a good solution.

If you’d like to re-try the assignment, now that you know the solution, follow the link below: \url{https://colab.research.google.com/github/tinyMLx/colabs/blob/master/2-1-13-AssignmentQuestion.ipynb}


\section{Initialization and learning}

Recall that Machine Learning, the core process powering TinyML, is the process of having Data and Answers, and from them attempting to infer the rules that link them:

\begin{figure}
  \includegraphics[width=\textwidth]{TinyML/1 Fundamentals/MachineLearning}
    \caption{Answers and data feed into machine learning and from this rules emerge.}
    \Mynote{besser mit tikz}
\end{figure}  



So, if your data is the numbers in this set: $[-1, 0 , 1, 2, 3, 4]$, and your corresponding answers are the numbers in this set: $[-3, -1, 1, 3, 5, 7]$, one way to begin to infer the relationship between them (assuming it's linear) is to have a function, where we say Y = wX+b, and we have to try to figure out the values of w and b.

A process to do this is shown below:

\begin{figure}
  \includegraphics[width=\textwidth]{TinyML/1 Fundamentals/MachineLearningLoop}
    \caption{Making a guess leads to measuring your accuracy which leads to optimizing your guess. This then gets repeated where you make a guess again and repeat the process. Arrows connect the steps.}
    \Mynote{besser mit tikz}
\end{figure}  


We could make a guess as to the values of $w$ and $b$, and measure how accurate that guess is.

So, for example, if we guess that $w=10$, and $b = -10$, we would get the set of answers for our data as $[-20, -10, 0, 10, 20, 30]$, where our correct answers are $[-3, -1, 1, 3, 5, 7]$.

From this we can measure our loss, by plotting our predicted answer against our actual answer.

\begin{figure}
    \GRAPHICSC{1.0}{1.0}{TinyML/1 Fundamentals/Scatterplot}
    \caption{A scatter plot showing the correct answer as compared to the predicted answer which is plotted for the given guess. The correct answer dots start at $y = -3$ at $x = -1$ and increase linearly to $y = 7$ at $x = 4$. This is compared to predicted values of $y = -20$ at $x = -1$ increasing linearly to $y = 30$ at $x = 4$.}
\end{figure}  



So, for example, for $x=-1$, we predicted $-20$, when the answer is actually $-3$, so we are off by $-17$


\begin{figure}
    \GRAPHICSC{1.0}{1.0}{TinyML/1 Fundamentals/Scatterplot2}
    \caption{The same scatter plot from above is shown highlighting the error at $x = -1$ showing an error of $-17$ (predicted value of $-20$ vs. actual value of $-3$).}
\end{figure}  



And when x=3, we predicted the answer to be 20, when in fact it is supposed to be 5, so we have an error of 15.

\begin{figure}
    \GRAPHICSC{1.0}{1.0}{TinyML/1 Fundamentals/Scatterplot3}
    \caption{Again the same scatter plot from above is shown now highlighting the error at $x = 3$ of $15$ (predicted value of $20$ vs. actual value of $5$).}
\end{figure}  

If we want to find out our error, we could average each of these.

However, by doing this our negative (-17) and positive (15) errors would mostly cancel out, so the smart thing to do to ensure that we don’t have them do this is to square the error amount, average out the squares of the errors, and then get the square root of that.

This gives us a loss function of root mean squared error.

When we now want to optimize our guess, and do better than these results, we can look at the loss function, and figure out a way to minimize our loss. As the loss function uses a square, the curve of the function is parabolic.

\begin{figure}
  \includegraphics[width=\textwidth]{TinyML/1 Fundamentals/Lossfunction}
    \caption{A graph of a parabolic loss funciton curve showing its u-shape.}
\end{figure}  


And the value of our loss with our current parameters can be plotted on this. Our loss was high, so we can say we’re pretty far up the curve.

\begin{figure}
  \includegraphics[width=\textwidth]{TinyML/1 Fundamentals/Lossfunction2}
    \caption{The same u-shaped loss graph as above showing a ``high'' loss plotted with a circle placed relatively high up on the left side of the curve.}
\end{figure}  


When we differentiate our values against the loss function, we’ll get a gradient, and we can use the gradient to figure out which direction we have to go in to move down the slope!

\begin{figure}
  \includegraphics[width=\textwidth]{TinyML/1 Fundamentals/Lossfunction3}
    \caption{The same u-shaped loss graph as above showing a ``high'' loss plotted with a circle placed relatively high up on the left side of the curve showing the gradient of the value as an arrow pointing down toward the bottom of the graph.}
\end{figure}  

And we can then move down the slope in the direction derived from gradient. We’ll ‘jump’ by a certain amount, and we can call that amount the ‘Learning Rate’

\begin{figure}
  \includegraphics[width=\textwidth]{TinyML/1 Fundamentals/Lossfunction4}
    \caption{The same u-shaped loss graph as above showing a ``high'' loss plotted with a circle placed relatively high up on the left side of the curve showing the gradient of the value as an arrow pointing down toward the bottom of the graph AND showing the actual step you take depending on the learning rate as a smaller arrow pointing in the same direction as the gradient down toward the bottom of the graph.}
\end{figure}  

After that, we are now closer to the bottom of the curve. The parameters that give us our location on the curve can then be the next ‘guess’, and by definition, these will have a lower loss, and we can repeat the loop. This is called back propagation.

\begin{figure}
    \GRAPHICSC{1.0}{1.0}{TinyML/1 Fundamentals/Lossfunction5}
    \caption{The same u-shaped loss graph as above showing that now our dot has moved lower on the left side of the u-shaped graph by following that gradient and taking that step mentioned above.}
\end{figure}  



The process can be repeated, and over time the value will get closer and closer to the bottom of the parabola, which is the minimum value of the loss function.

This process of using the gradient of the value to figure out the direction of the minimum, and then jumping down the curve towards the minimum is called gradient descent.


\section{Neurons in code - Understanding Neurons}

PROFESSOR: Previously, we looked at how you could fit parameters
to a linear equation and using this you could have a computer
learn over time the best set of parameters to solve that equation.
This really is the fundamental underpinning of machine learning.
And in this video, we'll see how that principle forms
the concept of a computational neuron, which is
the building block of a neural network.
Earlier we saw the code for what I like to call the Hello World of neural
networks, where you'll use a neural network to learn the relationship
between a set of x's and a set of y's.
The architecture there was a neural network with a single layer.
And this contained a single neuron indicated by this line of code
where we have a sequential defined and it only has one entry for the one layer
and that entry has only one unit for the single neuron.
So our neural network will really look like this, a single neuron with x in
and y out.
After training the neuron to learn the relationship between x and y,
we can then ask it to predict the y for a new x.
So for example, what is the y when x equals 10?
We know that the relationship between x and y is y equals 2x minus 1,
and the neuron learns something close to that.
When you print the results of model.predict(10.0),
you'll get something close to 19, but not exactly 19.
It'll be more like 18.998.
To see why, let's look under the hood a little.
Revisiting our neuron with an x in and a y out,
we can say that y is a function of x.
And that function is y equals wx plus b, the linear equation.
So we could redraw our architecture a little bit,
thinking in terms of parameters.
If y equals wx plus b, we need to know the values of w and b
in order to get y from x.
So how can a computer figure these out?
Well, we can start by emulating the neuron.
And here we can do so with a class called Model.
This has internal values w and B, and we can initialize them
to anything we want.
Say, we set them to be 10.0 each.
The class will also have a call function that will return wx plus
b for a given x.
So for our known set of x's, let's see what this model will give.
As we initialize w and b to be 10, y equals wx plus b will give us this set.
Where x is minus 1, y be 0.
Where x is 0, y will be 10, and so on.
Now this, of course, is very different for the set of y's that we already
know.
When x is minus 1, y is supposed to be minus 3.
But we got 0.
When x is 0, y is supposed to be minus 1, but we got 10.
Indeed, when x is 4, y is supposed to be 7, but we got 50.
We're clearly way off.
So recall the paradigm diagram where we make a guess, measure our accuracy,
optimize based on this information, and then repeat.
You can see at this point with w equals 10 and b equals 10,
we've made our first guess.
And by comparing our y's with the expected ones,
we can see that our accuracy was way off.
If we want to measure that accuracy, we can create a loss function in code
where we give it the predicted y, which is the set we calculated,
and the target y, which is the set that we were supplied.
And from there, we can calculate the root mean square of the difference
as explained in an earlier lesson.
Then this code is our underlying training code.
We'll use a gradient tape to keep track of our differential variables.
And within it, we can calculate our loss using the loss function
that we just saw.
And you can see that here.
Our current loss is the result of the loss function if we pass it the values
returned by the model and our real y's.
The next step is to optimize our guess.
And as we're using a gradient tape which tracks our variables,
we can differentiate our model's w's and our model's
b's against the loss function to get the gradients for w and b.
These gradients then help us get the direction
towards the bottom of the loss curve, so we can generate our next guesses
by reassigning w and b within our model to a new value in the direction
that the gradient gave us.
So on our machine learning paradigm diagram, the next step was to repeat.
We now have a new w and a new b, so we can make a guess with them
and repeat the whole process.
We can see that in simple code like this,
where we loop through the training function
that we just looked at 50 times.
So 50 times, we'll make a guess, measure the loss,
differentiate to get the gradient of the variables with respect to the loss,
and then use that to tweak their values in the direction of loss and so on.
By the time we've done this, our initial w
equals 10 and b equals 10 will have changed to something like this--
w equals 1.98 and b equals negative 0.94.
Now you and I both know what the desired values are based on the x and y
that we were given and they're supposed to be 2 and minus 1, respectively.
So this process has gotten us pretty close.
And now if we want to use this neuron to give us the y for any future x,
it will calculate y equals 1.9 x minus 0.94.
OK, so now it's your turn to try this out.
I've given you the code in the colab, so try it out
and you can see it all in action.

\section{Coding exercise: neurons in action}

Now that we’ve seen that we can use a single neuron to solve the linear regression problem, go back to the minimizing loss Colab and see if you can spot how we were using neurons all along. Do you find this surprising? Did you figure it out the first time you saw the Colab? If so, what gave it away?

\url{https://colab.research.google.com/github/tinyMLx/colabs/blob/master/2-2-3-MimimizingLoss.ipynb}

\section{Coding stepback}

In the previous video we created a function that had two parameters -- w and b, and returned a value f(x) = wx+b. You then saw how to use the machine learning training loop to adjust these parameters so that the correct values could be ‘learned’ over time.

You also, earlier, saw how this works in TensorFlow with machine learning.

\begin{lstlisting}
my_layer = keras.layers.Dense(units=1, input_shape=[1])
model = tf.keras.Sequential([my_layer])
model.compile(optimizer='sgd', loss='mean_squared_error')

xs = np.array([-1.0,  0.0, 1.0, 2.0, 3.0, 4.0], dtype=float)
ys = np.array([-3.0, -1.0, 1.0, 3.0, 5.0, 7.0], dtype=float)

model.fit(xs, ys, epochs=500)
\end{lstlisting}

I've changed the code slightly so that instead of having the layer code within the \PYTHON{tf.keras.Sequential()}, I declared it as \PYTHON{my\_layer}. Now, if you train this network, and try to predict a value for a given X like this:

\PYTHON{print(model.predict([10.0]))}

You'll see a value that's close to \PYTHON{19}. It did this by learning the internal parameters of the neuron, and you can inspect the neuron by looking at \PYTHON{my\_layer} using it's \PYTHON{get\_weights()} parameter:


\begin{lstlisting}
print(my_layer.get_weights())
[array([[1.9980532]], dtype=float32), array([-0.9939642], dtype=float32)]
\end{lstlisting}

You'll see that you get back two arrays. The first contains the \PYTHON{w} value -- which after running for 500 epochs gives you a value that's very close to 2! Similarly the second contains the b value, which got learned to be very close to -1.

But what would it look like if you used more than just a single neuron? So, for example, if you used a multiple neural network that looks like this?

\begin{figure}
    \GRAPHICSC{1.0}{1.0}{TinyML/1 Fundamentals/Neurons}
    \caption{Two inputs from $x$ into a neuron. Each neuron creates an output that is combined leading to the output $y$.}
\end{figure} 


To implement this in code, you'd use 2 layers, the first with 2 neurons, and the second with 1 neuron. It would look like this:


\begin{lstlisting}
my_layer_1 = keras.layers.Dense(units=2, input_shape=[1])
my_layer_2 = keras.layers.Dense(units=1)
model = tf.keras.Sequential([my_layer_1, my_layer_2])
model.compile(optimizer='sgd', loss='mean_squared_error')

xs = np.array([-1.0,  0.0, 1.0, 2.0, 3.0, 4.0], dtype=float)
ys = np.array([-3.0, -1.0, 1.0, 3.0, 5.0, 7.0], dtype=float)

model.fit(xs, ys, epochs=500)
\end{lstlisting}


So, if you look at the parameters in this case, there’ll be a little bit of a difference in the second layer. Typically we have said that our neuron has an input, x, and an output, y, and y will equal wx+b, where w and b are learned parameters. However, when we have 2 inputs, as you can see above, what will happen is that the formula will change where there is a separate w for each input. 



The neuron in the second layer has 2 inputs, so it will, instead of having \PYTHON{y = w*x+b}, it will have \PYTHON{y = w1*x1+w2*x2+b}, where \PYTHON{x1} is the output of the first neuron in the previous layer, and \PYTHON{x2} is the output of the second neuron in the previous layer. Naturally, if there are more than 2 neurons in the previous layer, then that number of weights will be learned.

If you run the above code to learn the parameters, then inspect the parameters:

\PYTHON{print(my\_layer\_1.get\_weights())}

You'll see something like this:

\PYTHON{[array([[ 1.4040651, -0.7996106]], dtype=float32), array([-0.50982034,  0.20248567], dtype=float32)]}

This will give you the weights and biases for the neurons in the layer. Note that it isn't a list of weight and bias for the first neuron followed by weight and bias for the second. In the above example 1.4040651 is the learned weight for the first neuron and -0.7996106  is the learned weight for the second. Similarly -0.50982034,0.20248567 are the learned biases for the first and second neurons respectively.



Similarly:

\PYTHON{print(my\_layer\_2.get\_weights())}

Will give you something like:

\begin{Python}
[array([[ 1.2653419 ],
        [-0.27935725]], dtype=float32), array([-0.29833543], 
        dtype=float32)]
\end{Python}

As mentioned earlier -- you can see that there are 2 weight values in this array, and a single bias. These weights are applied to the output of the previous neuron, and then they are summed and added to the bias.

You can inspect them manually, and apply the sum yourself like this:

\begin{lstlisting}[language=Python]
value_to_predict = 10.0

layer1_w1 = (my_layer_1.get_weights()[0][0][0])
layer1_w2 = (my_layer_1.get_weights()[0][0][1])
layer1_b1 = (my_layer_1.get_weights()[1][0])
layer1_b2 = (my_layer_1.get_weights()[1][1])

layer2_w1 = (my_layer_2.get_weights()[0][0])
layer2_w2 = (my_layer_2.get_weights()[0][1])
layer2_b = (my_layer_2.get_weights()[1][0])

neuron1_output = (layer1_w1 * value_to_predict) + layer1_b1
neuron2_output = (layer1_w2 * value_to_predict) + layer1_b2

neuron3_output = (layer2_w1 * neuron1_output) + (layer2_w2 * neuron2_output) + layer2_b

print(neuron3_output)
\end{lstlisting}

Will output:

\PYTHON{[18.999996]}

Practice this code for yourself in the colab!

This same process will apply for bigger and more dense neural networks, and will allow you to build models that learn more sophisticated patterns. You’ll see this shortly with an introduction to Computer Vision, where the model can learn the patterns in pixels in an image and classify them against labels. Then, later in the course, we’ll explore how these models need to be adapted in the context of TinyML.

\section{Coding exercise: multi-layer neural network}

Now let's explore the code from the previous reading through a Colab. You’ll be able to train both the single layer and multi-layer model described in the previous reading and explore their predictions and trained weights and biases. Which model learns faster? Which ends up making better predictions?

\url{https://colab.research.google.com/github/tinyMLx/colabs/blob/master/2-2-5-FirstNeuralNetworkRevisited.ipynb}


\subsection{Introduction to classification - Going Further}

SPEAKER 1: Up to this point, you've been exploring
how computers learn with machine learning
through fitting internal parameters of a function, or a neuron,
to match x and y values together, when there
was a linear relationship between them.
This is often called regression, where you're
using a neural network to predict a single value from one or more inputs.
Another common use for neural network is classification, where the network
can tell you something about the data.
For example, if it's an image, it can tell you the contents of that image.
Or if it's a set of values describing an object,
then it can tell you what that object is.
We'll now look at how we can move from our single neuron example
to a multi neuron, multilayer network.
And then how that can work to classify content.
To this point, we've been looking at a single neuron
and using it to get the y value for an x, when there's
a linear relationship between y and x.
But neural networks can be more sophisticated than this.
For example, we can have a network with two layers--
the first with two neurons and the second with one.
The neurons in the first layer each have an internal w and B that they can learn
and they will then output that to the second layer, which
will calculate based on its inputs, and not the x's directly.
With a single output layer like this, we can do a regression
where we calculate a value, like we've been using all along.
We can also have multiple neurons in the output layer instead.
This methodology is often used in classification.
So if we want to design a neural network that
can tell the difference between a cat and a dog,
then we could have one output neuron representing cats
and another representing dogs.
So then, if an image of a cat gets fed into the neural network, the neuron
representing cats, will have a high probability value,
and the one representing dogs will have a low probability value.
And similarly, if we feed in an image of a dog,
then the neuron representing a dog should have a high probability,
and the one representing a cat should have a low probability.
And this gives us a way of labeling cats and dogs
by having our label contain the desired output of the set of output neurons.
So if our first neuron represents a cat and our second represents a dog,
then we can say that our label for cat images is 1, 0.
And our label for dog images is 0, 1.
Or if we want to create a neural network that recognizes the handwriting
digits 0 through 9 with an output neuron representing each,
we could use labels like this.
Every output neuron is 0, except the one representing the digit.
Let's now look at some code that will use a neural network
to recognize the digits 0 through 9.
We'll need data-- images that are labeled to train with.
And fortunately, there's a data set called MNIST
that's already built into TensorFlow.
With these two lines of code we can load the data.
The data set is split into training images and labels, as well as
validation images and labels.
The data set looks like this.
There are 60,000 labeled images that you can use for training and 10,000
that you can use for testing or validation.
And we'll get all of that simply by asking
TenserFlow for the MNIST training and validation sets.
The handwriting digits are 28 by 28 monochrome, which means
each pixel value is from 0 to 255.
We want to normalize these to use them in a neural network.
And it behaves much better if the values are between 0 and 1.
So we can just divide the image values by 255.
Then next, we're going to define a neural network.
Earlier, you learned about sequential and how
sequential can be used to define a neural network by listing
the layers of that neural network.
And that's what you're seeing here.
It's a little more complex than the single neuron that you saw previously.
Our images are 28 by 28 in dimension.
So the first thing that we can do in our sequential
is to define an operation to flatten these out,
so that we can feed them into the neural network.
The next layer has 20 neurons, so we need our images
to be in the same dimension as that layer, considered
to be 20 by 1, so our images will also need to be something by 1,
in order to fit.
The easiest way to do this is to flatten the 28 by 28, which is 784 pixels,
into a 784 by 1.
And that's what a flattened layer does.
We can tell it to expect the input shape to be 28 by 28.
If it gets anything else, it'll throw an error.
Next, we'll define a layer of neurons.
And we have 20 neurons in this layer.
This number is purely arbitrary for now.
And later you can explore methodologies to pick the optimum number.
Too few, and you may not have enough to learn the contents of the image.
Too many, and you could over specialize or your network
could be very slow to learn.
The activation function is a new concept that we'll explore shortly.
The final layer will have 10 neurons in it because we have 10 classes.
And these are the digits 0 through 9.
It also has an activation function, but we'll explore that shortly too.
So our network with 20 neurons in one layer and 10 in the other,
will look like this, where every neuron is connected
to every neuron in the other layer.
All of these connections inspire the word "dense" for this layer type.
We'll want to feed the pixels of an image into the neural network.
So for example, these pixels are from one of the training images,
and we can see it's the number 7.
So we would desire that the neuron representing the number 7
is the one that would light up.
And we'll train it with this set of 0s and a 1 as the label representing 7.
So remember, that our image is 28 by 28.
So instead of the square image like this being fed in,
it'll be flattened into 784 by 1.
And this will be fed into the neural network instead.
So if we return to our neural architecture,
we can see this definition--
the flattened layer to flatten the images, the dense layer of 20 neurons,
and the output dense layer of 10 neurons for our 10 classes.
Our layer of 20 neurons has something called an activation function on it.
It's called ReLU, which stands for rectified linear unit.
The important thing to know is this--
every neuron will call its activation function when that layer is used.
The ReLU activation function changes any output that is less than 0 to 0.
So when we evaluate y equals mx plus b on any of these neurons,
we'll throw away any values that are less than 0.
It's really commonly used in dense layers,
so that neurons don't cancel each other out.
Imagine if one of these neurons return the value 10 and the other minus 10.
Without ReLU, we'd get a net of 0 from that pair of neurons.
And in general, you don't really want this.
For the math-inclined, you might also notice
that this can introduce a non-linear relationship between the layers.
And that can help you capture complex relationships beyond just
what the neuron alone can do.
Our output layer has another activation function.
This time it's called softmax.
What this will do is help us find the neuron from amongst the 10 that
has the highest value.
It saves us from writing a lot of code to loop through the outputs
and pick the biggest one.
Our process of compiling and training is pretty much the same
as we had for a single neuron example.
The key difference, of course, is the functions
that we use for optimization and loss.
In this case, we'll use an optimizer called Adam.
TenserFlow has many built-in optimizers, but instead
of worrying about the math for optimization
or the gradient calculation and descent stuff that we saw earlier,
we can just pick a function and use it, making it easy for us to experiment.
Adam is particularly useful because, if you
recall when we spoke about gradient descent and the jump size
from the learning rate, Adam can actually vary its learning rate
helping us to converge more quickly.
And we can say the same for loss functions.
There's a menu of them that we can choose from,
but often our loss function has to map closely to our scenario.
For example, when we were doing regression,
mean squared error made sense.
But it doesn't make sense for classification,
as we have a number of categories instead of a single value.
And we need the loss across all the categories.
This means that you'll generally use a categorical loss function.
And you can see that in the name, like I'm using here.
Model.fit runs the training loop.
And in this case, we can train for just 20 epochs.
And then we can test how well our network did by calling model.predict
on one or more images.
Here it's running out on the entire validation set of 10,000,
but I'm just printing out the classification
it gave me for the first one, and comparing it against the actual label.
This set of 10 numbers is the set of 10 probabilities for each of the 10
values as output by the 10 neurons.
And the 7 is the label for the first image, so we know it's a seven.
If we inspect closely, we can see that the neuron representing the number 7,
has by far the highest value.
It's telling us we have a 99.89% chance that we're looking at a 7.
If you look closely, it's the eighth neuron
because we started counting in 0 i.e. the first neuron is
the probability we're looking at 0.
The second neuron the probability we're looking at 1, and so on.
So next you can try this out for yourself,
and you can inspect the neurons in each of the layers
to see what they've learned.
And you can see also how this maps to classifying an image.

\section{Coding exercise: DNN}

Now let's explore the MNIST classification example through a Colab. You'll get to train the model just described in the previous reading, learn how to inspect and analyze the various trained layers in the model, and test the final model on additional data sets!

\url{https://colab.research.google.com/github/tinyMLx/colabs/blob/master/2-2-7-ExploringCategorical.ipynb}

\section{Training data, validation data, test data -- Learning from zero\ldots}


LAURENCE MORONEY: We've been looking at classification
and, in particular, classifying images.
When we trained, we looked at the accuracy of the training
and after some basic testing, we were able to see that our neural network was
able to recognize things pretty well.
But that can lead you into a false sense of security.
And I'd like to explore some of the issues around this
and look at some methods that you can use to avoid errors when
you train your neural network naively.
Imagine a scenario where you want to train
a neural network to recognize shoes.
It's like teaching somebody who had never seen a shoe before about what
a shoe actually is so that in the future when they see an object,
they can decide whether what they're seeing is a shoe or not.
Now, we know there's a huge variety of shoes,
and there's no hard and fast rule about what makes a shoe a shoe.
Typically, the steps that you'd follow are
to get as many examples of shoes as possible,
train a neural network using those examples, and then profit.
Earlier, you trained a neural network to recognize handwriting images,
and it needed 60,000 examples.
That will be a lot of shoes if you were going
to try to train it with shoes instead.
But you'd still have to gather as many shoes as possible
and train by looking at them.
And you might end up with something like this.
An increasing accuracy number as you train.
By the time it finishes training, it has a very high accuracy, maybe even 100%.
Now that might mean you think you've written
an amazing AI that can recognize shoes, and it's time to profit.
But then you show it a shoe like this, and it fails.
You thought you were 100% accurate in recognizing shoes.
But the reality is that you are 100% accurate at recognizing
the types of shoe that you trained the neural network on
and that 100% figure led you to a false sense of security.
Your perfect neural network architecture may not be so perfect after all.
There's a number of ways to help prevent this, and we'll look at the simplest
first.
Consider our shoe scenario where we had a lot of shoes,
and we can represent them using the yellow rectangle.
We used all of our data to train our neural networks, and why not?
The more we train, the better, right?
In principle this sounds fine, but we've no way
of testing it against previously unseen data
because we used all of our data to train it.
So what if we hold back some of the data?
Don't train with all of the data.
Save some for validating that the neural network is training well.
And by having this set, then we can use that to measure
how it does with unseen data.
And we can represent this with the red rectangle.
The full data is still the same, but the training data is now smaller.
It's the full set minus the validation set.
And then, let's take this a step further.
How about holding back a little more that we
can use to test our neural network after it's done training.
Our first subset, the yellow rectangle, is the data that we'll train with.
Our second subset, the red rectangle, is the data
that we'll use to validate the training by testing after each step
how the network is doing.
But we don't use this data to train the network.
Our third and final subset is the blue rectangle.
And this can be used as a clean set to test the network
after we're done training it.
Following this methodology, we could, for example,
pick a neural network architecture and train it.
On the training data, it's 99.9% accurate.
But it does worse on the other two sets.
For example, say it's only 80% on the test set.
This can be like our shoe scenario earlier.
We designed a neural network that's great on the training data
but not so great on the other data.
Our 99.9% makes us think we have a much better network than we actually do.
So why don't we redesign our neural network architecture and try again?
Our accuracy on training might go down but what's most important
is keeping the accuracy of the network on the training data
as close to the validation and test accuracy as possible.
And that will give us a true signal of the real accuracy of the network.
So returning to the code that we've been using to train endless handwriting
digits, recall that we had training images and labels
and validation images and labels in the data set.
We could load them with one line of code in TensorFlow.
But when we trained, we just used the training images and labels.
And we could see our training accuracy epoch by epoch.
It got to 97 and 1/2% and just 20 epochs and that's pretty good.
But that was just on the training data.
We have no context for how good or how bad it is with previously unseen data.
So we can update our model.fit to use the validation data like this.
When you train, you then get reports on the accuracy
of the network on the training data step-by-step
as well as the validation data.
I've highlighted them in this diagram.
And while the accuracy on the training data gets up to 97.59%,
the accuracy on the validation data is at 96.26%.
They'll almost never be the same.
And your goal is to strive to get them as close to each other as possible.
This is a sign that your neural network is good at generalizing and not
over specializing to the training data.
If you have a test set, MNIST doesn't include one
but other data sets do, you can then perform model.evaluate on it
after training is complete to see the results.
And you'll get an accuracy value for that.
You'd expect it to be similar to the validation one
but maybe a little bit lower if your network is well-architected.
Your confidence and your network's ability to classify data
should be based on this value and not the training one.
Now that you've looked into classification and understanding
how to train with split data sets, the next thing
is a quick quiz to test you on what you've learned so far.

\section{Realities of coding}

When writing code for neural networks, often the bulk of your attention will be on the model architecture -- because you want to find the optimum architecture for accuracy on your training set, and that accuracy also carries over to validation and testing.

This is good practice, but do keep in mind that often you will have many more lines of code to handle everything else in your system.

With MNIST, your code looked like this:

\begin{lstlisting}
import tensorflow as tf
data = tf.keras.datasets.mnist
(training_images, training_labels), (val_images, val_labels) = data.load_data()
training_images  = training_images / 255.0
val_images = val_images / 255.0

model = tf.keras.models.Sequential([tf.keras.layers.Flatten(input_shape=(28,28)),
                                    tf.keras.layers.Dense(20, activation=tf.nn.relu),
                                    tf.keras.layers.Dense(10, activation=tf.nn.softmax)])

model.compile(optimizer='adam',
              loss='sparse_categorical_crossentropy',
              metrics=['accuracy'])

model.fit(training_images, training_labels, epochs=20, 
          validation_data=(val_images, val_labels))
\end{lstlisting}


Note that it only took 3 lines of code to load and prepare your data -- this is because the MNIST dataset was already available to you, so you could get it with 1 line, and then it was just another 2 lines to normalize the images. As you build real systems, getting the data may not always be this easy -- and it might require significant coding to get it into a place where you can train a neural network with it. Going into the specifics of that is beyond the scope of this course, but as you continue to learn neural networks, it’s good to practice with different types of data, so you can prepare them for ingestion into your neural network architecture.

Indeed, often you'll find about 80\% of your code might be in the preparation, and only 10\% (or less) in your neural network architecture, leaving the other 10\% for testing and other tidy up.

\section{Coding assignment}

In this assignment we'll dive a little deeper with a series of hands on exercises to better understand DNN learning with Tensorflow.

Remember that if you are taking the class for a certificate we will be asking you questions about the assignment in the test!

\url{https://colab.research.google.com/github/tinyMLx/colabs/blob/master/2-2-12-AssignmentQuestion.ipynb}


\section{Coding assignment solution}

We hope you enjoyed diving deeper into DNNs and TensorFlow.

For the data normalizing / rescaling we expected you to use:

\begin{lstlisting}
training_images  = training_images / 255.0

test_images = test_images / 255.0
\end{lstlisting}

This is because the RGB data format uses values in the range [0,255] to represent each color and so dividing the color values by 255 will result in the data falling in the range [0,1]. Note that we need to make sure to use floating point division or the number will always be either 0 or 1 instead of in the range [0,1].

Later running the following:

\begin{lstlisting}
classifications = model.predict(test_images)

print(classifications[0])
\end{lstlisting}

Should result in an output that looks something like:

\begin{lstlisting}
[1.6892707e-05 1.0681998e-06 1.3354841e-06 5.3753224e-06 2.1088663e-06

5.2157331e-02 9.2272086e-07 6.5788865e-02 2.6547234e-03 8.7937140e-01]
\end{lstlisting}

These numbers are the probability that the value being classified is the corresponding value, i.e. the first value in the list is the probability that the handwriting is of a '0', the next is a '1' etc. Notice that they are all VERY LOW probabilities except the last value which is about 88%. That is the neural network telling us that it is fairly confident that it is seeing an example of the 10th class.

We then asked you to double the number of neurons in the model this resulted in a code input that should have looked like:

\PYTHON{NUMBER\_OF\_NEURONS = 1024}

This should have resulted in a more accurate model that took longer to train as it required more computation but was able to better fit the data as our original model was too small to adequately learn this particular dataset. However, that doesn't mean it's always a case of 'more is better', you can hit the law of diminishing returns very quickly!

You then added a layer into your model which there are many possible answers for how you could have done it (different numbers of Neurons and/or different activation functions). One possible suggestions is as follows:

\PYTHON{YOUR\_NEW\_LAYER = tf.keras.layers.Dense(512, activation=tf.nn.relu)}

When training your new model if you also used a large number of Neurons in your new layer, e.g., 512, you should find that training accuracy improves in a similar manner to doubling the number of Neurons in the single layer. If you are lucky it might even be better! If you chose a small number of Neurons you may have seen no effect. This is simply an artifact of the particular model structure and dataset we are using. Again, remember there is no one size fits all solution in Neural Network model design.

We then un-normalized the data by inverting whatever you did at the start of the assignment. For us that would be:

\begin{lstlisting}
training_images_non = training_images * 255

test_images_non = test_images * 255
\end{lstlisting}

Which resulted in the model doing a worse job in training due to the poor numerics of the un-normalized data.

Finally we added a custom callback function to help us exit training early if we reached some goal. In this case to exit early but only upon reaching at least the 2nd epoch one thing we could do is exit once the training accuracy is over 86% (as in all of our training runs we got as high as an 83.5% accuracy on the first epoch). One way to do that would be:

\PYTHON{logs.get('accuracy')>0.86}

Of course if you model does better than that then you’d have to adjust that threshold.

As always you’ll find the link to the assignment again below in case you want to explore it a bit more now that you have a solution.

\url{https://colab.research.google.com/github/tinyMLx/colabs/blob/master/2-2-12-AssignmentQuestion.ipynb}


\section{Recap}

Up to this point you have learned the Machine Learning paradigm, and how it works to allow a computer to figure out the parameters that fit a function. When that function gives you the rules to map data to answers, you’ve changed from a programming paradigm where you have to figure out the rules that act on the data to give you answers:


\begin{figure}
  \includegraphics[width=\textwidth]{TinyML/1 Fundamentals/TraditionalProgramming}
    \caption{Rules and data feed into traditional programming which then outputs answers.}
\end{figure}  

To one that looks like this, where you provide the answers with the data, and the computer figures out the rules that maps them to each other:


\begin{figure}
  \includegraphics[width=\textwidth]{TinyML/1 Fundamentals/MachineLearning}
    \caption{Answers and data feed into machine learning and from this rules emerge.}
\end{figure}  

To achieve this, of course, you still need to write code, but instead of writing the code for the rules, you write the code that figures out what the rules are. The algorithm to do this looks like the following:


\begin{figure}
  \includegraphics[width=\textwidth]{TinyML/1 Fundamentals/MachineLearningLoop}
    \caption{Making a guess leads to measuring your accuracy which leads to optimizing your guess. This then gets repeated where you make a guess again and repeat the process. Arrows connect the steps.
    }
\end{figure}  

As an example, you used two sets of numbers, for X and Y and wrote code that used the above algorithm to figure out the relationship between them where Y = wX+b, and the above algorithm could figure out the values of w and b that would fit.

This was the basis of a mathematical neuron, where the neuron could store the values of w and b, and again, using the above algorithm, those values could be figured out to give the desired Y for an input X.

\begin{figure}
  \includegraphics[width=\textwidth]{TinyML/1 Fundamentals/Neuron2}
    \caption{An input of X is fed into a neuron and an out Y is created.}
\end{figure}  


When these neurons are used together, in layers, more complex relationships could be figured out, and you saw how having multiple output neurons could take you into the territory of classification, where, each neuron could have a role in the final output, which of course is perfect for computer vision. Below is a simplified neural network diagram illustrating this point where input data can be classified -- i.e. does the network ``see'' a dog or a cat.

\begin{figure}
  \includegraphics[width=\textwidth]{TinyML/1 Fundamentals/CatDog}
    \caption{An image of a cat’s face is fed into multiple neurons. Each of these neurons outputs a decision of whether the cat’s face represents a cat or a dog.}
\end{figure}  

Basic neurons, stacked together like this form what is called a Dense layer, based on the fact that they’re densely connected to each other. When there are multiple layers between the input and the output, you have a Deep neural network, which gives us the term Deep Learning. Additionally, the layers between the input and the output are often called >hidden layers.

As you continue your Machine Learning journey, you’ll discover that there are many other types of layer that learn different parameters to just the w and b that a neuron learns.

For example, there are Recurrent layers, that are beyond the scope of this course, but they are able to learn values in a sequence where each neuron learns values and passes those values to another neuron in the same layer. The basic idea can be represented like this:

\begin{figure}
  \includegraphics[width=\textwidth]{TinyML/1 Fundamentals/Recurrent}
    \caption{A graphical depiction of a recurrent layer showing the input values x0, x1, x2 flowing into the three neurons and outputting values y0, y1, y2. Importantly the inputs go to and come from each neuron AND there are arrows denoting information flowing from the left to the right ACROSS the three neurons. This is the WITHIN layer flowing of information.}
\end{figure}  

Where, if we have a sequence of values, X0, X1, X2 that we want to learn a corresponding sequence Y0, Y1, Y2 for, then the idea is the the neurons (F) can not only try to match X0->Y0, but also pass values along the sequence, indicated by the right pointing arrows, so the input to the second neuron is the output from the first and X1, from which it can learn the parameters for Y1 and so on.

There are many variations on this idea, but a very powerful one is called Long Short Term Memory where not only are values passed from neuron to neuron, so that X0 can impact Y1, X1 can impact Y2, but values from further back can also have an impact -- so that X0 could impact Y99 for example. This is achieved using a data structure called a Cell State where context can be preserved across multiple neurons.

\begin{figure}
  \includegraphics[width=\textwidth]{TinyML/1 Fundamentals/Recurrent2}
    \caption{The same graphical depiction of a recurrent layer as above now overlayerd with a Cell State bar that connects all three neurons. This is the ´´memory'' that is shared across the layer on top of the information shared within the layer.}
\end{figure}  



Another type of neural network layer is called a Convolution, where a filter that can transform data, particularly image data, can be learned. You’ll explore those next!

I hope this helps you understand that Machine Learning goes beyond simple neurons that learn a w and b in the $Y=wX+b$ scenario.


\section{Introducing convolutions}


LAURENCE MORONEY: As you just saw, dense and flatten
aren't the only layer types supported in TensorFlow.
Another very important one, particularly for computer vision, is convolutions.
The main difference with a layer like this
is that you move beyond just learning weights and biases in your network.
In this video, you will explore how to learn filter values or convolutions
and why they're important.
And because convolutions are so efficient at helping
a computer understand the contents of an image,
you'll find they're extremely important in tiny ML too.
To recap, a single neuron can learn the relationship between y
and x, where y equals wx plus b.
And layers of these functions can be used to take in more complex values,
like the pixels of an image, and be trained
to classify images like a cat or a dog.
We then explored using the MNIST data set of handwritten digits
that are 28 by 28 and how a couple of dense layers
could be used to classify them.
Under the hood, every neuron was learning linear relationships,
and a complex set of these trained on pixel values
could be used to classify the simple images.
But what happens when the images get more complex,
like these images of a human and a horse on a beach?
With MNIST we had simple handwritten digits
that were centered within a 28-by-28 frame.
But with images like these, it can get much more complex.
We're zoomed in on the horse, for example.
The image may not take up the full frame,
either, or it could be off center.
Using convolutions or filters can help with this problem,
and we'll look at how they work next.
So for example, if you take a look at this image from Fashion-MNIST
it's a 28-by-28 grayscale image of an ankle boot.
We can use this to demonstrate a very simple filter.
For every pixel in the image, we'll take its immediate neighbors.
So for example, if the current pixel has value 192,
we can see that its neighbors above are 0, 64, 128.
Its neighbors below are 142, 226, and 168, for example.
We can then define a filter where we multiply each pixel
by the respective filter value.
So the pixel above and to the left, which is zero,
is multiplied by the top left filter value, which is minus one.
That gives us zero.
We do the same for each pixel, and we sum up the values.
This will become the new value for our current pixel instead of 192.
This might seem really unusual, but it's actually
a common tool in image processing.
And if you've ever done any kind of image processing
such as what you might do on Instagram or Photoshop,
that's pretty much how it works, by having a filter
and applying that filter to the image.
Here's a simple example.
Consider the image on the left.
If I apply the filter shown to it, I'll get the results on the right.
It greatly enhances the vertical lines, and it dims everything else.
So you could consider this to be a vertical line detector.
And similarly, this filter can spot horizontal lines,
darkening almost everything else in the image that isn't a horizontal line.
By applying filters like these, you can remove almost everything
but a distinguishable feature.
And this process is called feature extraction.
When combined with a process called pooling,
we can also do something really powerful.
Pooling is simply the process of removing pixels
while maintaining important information.
So for example, if you look at the four-by-four square of values
on the left of the screen, these emulate what pixels in an image
might look like.
We can break this down into four two-by-two blocks of pixels,
and then we can see that in the center of the screen.
So our 0, 65, 48, 192 from the top left becomes the first block.
The 128, 128, 144, 144 from the top right
becomes the second block and so on.
From each of these blocks, we pick the biggest value
and we throw the rest away.
We then reassemble these, and we have a new set of four pixels, which
was pooled from the original set of 16.
Now, where this gets really interesting and powerful
is if we apply it to an image after filtering the image.
It can have the effect of enhancing the features that we extracted.
And not only that, if you're applying many filters to your image in a layer,
you're in effect making many copies of your image.
So if it's a large image, you end up with a lot
of data flowing through your network.
So it's good to have a way to compress that data without losing
the important features.
And this is just one layer.
When you apply multiple layers of filters,
then really complex features, such as faces or hands
instead of the vertical or horizontal lines I've shown here,
could be spotted and extracted.
Recall the image earlier where we had that filter that
extracted vertical lines.
After a pooling, the image will look like the one on the right.
We can see that those lines pop a lot more.
And importantly, the image, of course, is now one quarter
the size of the original.
How does this impact computer vision?
Well, let's consider a scenario.
Here's a picture, and you don't know its content.
I've simulated that by making the image really blurry.
So say there's a filter or a set of filters
that can be passed over the image that extracts these pixels.
It's two vertical shapes, and they look a little bit like human legs.
And then maybe other filters can extract features like these.
They're blobs with cylinders that stick out of them.
Now our human brain sees these instantly as hands,
but an untrained network doesn't know that yet.
It just knows that a filter can extract something that looks like this.
And yet another filter extracts this, roughly a circle
with another two circles on it as well as some parallel features.
Our human brain recognizes this as a face with eyes and a mouth.
But, again, the computer has no context for what these are.
It only knows that a particular filter could extract this.
So then the computer could match the features that
were extracted by the first filter--
in other words, the legs--
with those from the second filter, also known
as the hands, and those from the third filter, also known as the face.
And if all of these features are present in a picture that is labeled human,
then when training a neural network, we could
have it learn the filters that extract information
that's consistently present in images that are labeled human.
So then these could be used to spot future pictures
to see if they contain a human.
And perhaps the network could spot different features
that are present in images of horses.
By learning sets of filters that can then spot human or horse,
we now have the beginnings of a computer vision model
that can handle sophisticated images and predict what's within them.
So to recap, filters or stacks of filters
can extract features like hands or ears.
And then during training, they can be combined
with the labels of the known images to help us train a network that
can predict image contents.
So our model, instead of just weights and biases, can now also learn filters.
So when we pass in data like an image, we can predict its contents.
Remember, the network learned filter values
that it has a high confidence that if they can extract features
and those features are present in the image then
it can match that to a label.
So now our inferences are going way beyond just y equals mx plus b.
So in this case, we know we have a human.
If we return to the machine learning paradigm diagram,
we remember it's looping through, making a guess, measuring our accuracy,
and then optimizing our guess repeatedly.
In this case, instead of weights and biases,
we initialize the filter values randomly,
and applying these to the images will give us our guesses.
By measuring our guesses as to the class of the image
based on the extracted features against the ground truth,
we can now measure our accuracy.
And as before, we can use an optimizer to tweak our filter values,
knowing that we'll be stepping in the right direction
before we repeat the process.
So if my filters end up turning my image into these three images
where I have extracted features that look like legs,
hands, and a head and this image is labeled human,
then these filters will be associated with human.
And then they can be used to predict pictures.
In this case, we can see that he's human,
but the computer could only recognize that by looking at the filter values,
extracting features from them and if they extract the features as expected.
Do note that I'm using legs and hands and faces here
as examples of things that were extracted that match a filter.
Now, these are things that our human brains will recognize as features,
so I used those to demonstrate the concept.
When the computer is matching features to an image,
it might see things in the image that a human doesn't.
There may be patterns of pixels that consistently
match on a labeled image that mean nothing to us as humans.
And these became the basis of DeepDream, where
parts of an image that match features in a data set
called Inception with millions of images of thousands of classes are extracted,
and then the original image that matches that feature is partially overlaid.
So for example, the patterns in the sand around the woman
closely match features that were extracted
from what looks like animal skins in the middle image
and are then overlaid and then later match the animals themselves,
which you can see overlaid in the third image.
And the effect ends up being trippy and dreamlike.
I thought it might be just fun to see this.

\section{Coding exercise: filters}

In this Colab you’ll explore how filters and convolutions can be used to extract features from images. You’ll see firsthand how different filters can extract different features. You’ll also learn about pooling and explore how it can help reduce the amount of information while retaining important features. Click the link to get started!

\URL{https://colab.research.google.com/github/tinyMLx/colabs/blob/master/2-3-3-ExploringConvolutions.ipynb}


\section{Converting DNN's to CNN's}

SPEAKER: So far, we've seen how to build a neural network that
uses dense layers or a DNN to learn how to classify images,
as well as the principles behind convolutions
that will learn filters that can be matched to features and images.
Let's now take a look at enhancing our existing DNN's
to add convolutions to them.
So we can see the impacts.
For a recap, here's our DNN encode that we've
been using to classify mnist or fashion mnist images.
For this example, I'll use the fashion mnist
version, which is a little bit more sophisticated than the handwriting
digits.
And it makes it harder for us to get a higher level of accuracy.
The model architecture is here.
And it's a simple sequential, which flattens the input and it
through dense layers.
When we train a model with it, we'll be able to see the accuracy and validation
accuracy reported back to us.
And in this case, after 20 epochs, I got an accuracy of 89.53%
and a validation accuracy of 86.77%.
These results are great.
But we might be able to improve on them using convolutions.
So let's take a look.
Now, here's the complete code for an updated neural architecture
to use convolutional neural networks, instead of a straight up DNN.
First, we can see that on our input layer, we need to use an input shape.
But it's a little different.
It's now 28 by 28 by 1, instead of 28 by 28.
And this is because a convolution of layer
expects the image dimensions to be fully defined.
So a color image will have three channels for red, green, and blue.
So it will be something by something by 3.
But a monochrome image only has one channel.
So it's specified here.
Our image is 28 by 28 with one channel.
So it's 28 by 28 by 1.
Our output layer will be 10 neurons as before.
It's because we have 10 classes.
We'll also have a flatten before we get into the DNN.
This will occur after the filters have done their job.
The filters will act on the image as we saw previously.
So they will still have a rectangular image.
And once they've given updated pixel values,
then we'll flatten that out before we get into the DNN.
And then, the network can behave as before.
So let's now look at how to define the convolutional neural network.
First, will specify the convolutional layers using Conv2D.
It's 2D, because the images have a width and height.
And we'll pass filters across them to change the image to extract features.
Remember, these filter values will be learned over time.
The Pooling that we described previously is implemented
by specifying max pooling as a layer.
There are different types of Pooling.
There's max, min and average.
But for this, we'll use max pooling, where we
take the largest value from the pool.
Looking at the convolutional layers, we can explore the parameters.
First is this number, which is the first parameter in the convolutional layer.
And this is the number of filters to apply to the image at this layer.
So in this case, we have 64 filters that will be applied to the image.
Remember that the filter values are going to be randomly initialized as our
Guess.
And then, they will be learned over time.
The 3 comma 3 here indicates the size of the filter.
Recall that when we were talking about filters that we had defined a 3
by 3 filter that multiplied across the current pixel and each
of its neighbors.
This could be a different value.
For example, 5 by 5 to look at neighboring pixels up to two pixels
away, and so on.
But it will always be odd numbers--
3 by 3, 5 by 5, 7 by 7.
And it's similar for the Pools, where the size of the Pool
is specified as a parameter, which in this case is a 2 by 2 pool.
So we'll be picking one pixel out of four.
If we want different size pools, we can define them here.
Remember that by Pooling, we're also reducing the size of our image,
helping us to emphasize and summarize the features that we extracted.
The convolutional layer had 64 filters in it,
which means 64 copies of the image are going to be made.
So reducing the size of them will reduce the amount
of data flowing through our network.
OK, now that you've seen how to update your code to allow for convolutions
and saw how you can improve the accuracy of a fashion mnist classifier using
them, next you're going to try a Colab
where you'll try it out for yourself, and measure the accuracy.

\section{Coding exercise: CNN}

In this Colab you’ll explore the power of Convolutional Neural Networks (CNNs). You’ll train both a traditional DNN and a CNN and see how CNNs can far outperform standard networks on computer vision tasks. You’ll then dive into both the convolutional and max pooling layers that power the CNN.

\URL{https://colab.research.google.com/github/tinyMLx/colabs/blob/master/2-3-5-FashionMNISTConvolutions.ipynb}

\section{Mapping features to labels}

Imagine you are teaching a computer how to ‘see’ something. By that we mean not just be able to process the pixels in an image for color or shade, or anything like that, but to be able to understand the contents of an image.

To simulate this, you could say that the understanding of the contents of the image are effectively blurred\ldots like this:




There's information in there, but I don't know what it represents, so I simulate that by blurring the image.

Now, next, imagine that there's a convolution (filter) that can extract something from the image consistently, and the something that it extracts is always present when the image is labelled with a particular class. In other words, if there’s a filter that always produces something like this, when the image is labelled human.


\begin{figure}
  \includegraphics[width=\textwidth]{TinyML/1 Fundamentals/Features2}
    \caption{Blurry image ranging from some tan to black colors. In the bottom two legs seem to be visible.}
\end{figure}  


You and I both know that these clearly seem to be human legs, but the computer doesn’t know that. It just knows that two cylinder-like objects like these tend to show up for a particular filter, and only on images that are labelled human.

Then, similarly, there’s another filter that produces this, and only for human images:

\begin{figure}
  \includegraphics[width=\textwidth]{TinyML/1 Fundamentals/Features3}
    \caption{Blurry image ranging from some tan to black colors. In the middle two human like hands appear to be visible but are still blurry.}
\end{figure}  



\ldots and then another that produces this, or something similar, again, only for human images:


\begin{figure}
  \includegraphics[width=\textwidth]{TinyML/1 Fundamentals/Features4}
    \caption{Blurry image ranging from some tan to black colors. At the top something that looks like a human face with some hair.}
\end{figure}  


You and I both recognize this as a human face. But the computer doesn’t. It only knows that something clear is regularly extracted by that filter, when the image is labelled as human.

Thus, when a set of filters is learned that consistently extracts content like this when the image is labelled as human, we could say that the following ‘equation’ holds:


\begin{figure}
  \includegraphics[width=\textwidth]{TinyML/1 Fundamentals/Features5}
    \caption{Collection of the previous three images, one with hands, one with legs and one with a face. The addition of these images gives an output of ``Human''.}
\end{figure}  




If there are 64 filters, for example, in the final layer, they may all return ‘nothing’ and it’s just these three end up having significance for this class.

Similarly, a different set of filters could return values for the label HORSE, and everything else (including the three filters that gave us human hands, legs and face) would return nothing, so we’d get:

\begin{figure}
  \includegraphics[width=\textwidth]{TinyML/1 Fundamentals/FeaturesHorse}
    \caption{Similar collection of images but with a a horse, with legs, ears and face of a horse.  The addition of these images gives an output of ``Horse''.}
\end{figure}  



Now we have a set of filters that a model has learned that can extract the features that indicate what is a horse and what is a human!

Do note that for this example I used features that you and I recognize, like hands, legs and feet as distinguishing between the two, but the computer is NOT limited to that. It might be able to consistently ‘see’ patterns in images that you do not, and that might be a more accurate determinant of the class of the image. The field of convolutional visualization studies this, and it’s fascinating to learn the interpretability of images that are classified using this method!

\section{Coding exercise: computer vision}

In this Colab you’ll build off of the prior Colab exercise using CNNs to learn the Fashion MNIST dataset, and you’ll start to visualize the convolution and pooling layers to better understand how the model “sees” the world. Hopefully this will give you more insight into how neural networks see the world as compared to how you see the world!

\URL{https://colab.research.google.com/github/tinyMLx/colabs/blob/master/2-3-7-FashionMNISTConvolutionsVisualizations.ipynb}

\section{Coding assignment: CNN}

In this assignment we are going to take your learnings from class about building CNN models and apply them to a new dataset, CIFAR-10. You’ll have to design your own model and then select your optimizer and loss function. In the end you should end up with a model that can relatively accurately identify objects from this much harder dataset!

\URL{https://colab.research.google.com/github/tinyMLx/colabs/blob/master/2-3-9-AssignmentQuestion.ipynb}


\section{Assignment solution}

We hope you enjoyed developing a model that can solve the CIFAR-10 dataset!

There are many possible solutions to this assignment. One good approach (as suggested by the hints), is to use convolutional and pooling layers to find high level features and then dense layers to classify those features. A possible solution to the model design is shown below:


\begin{lstlisting}
FIRST_LAYER = layers.Conv2D(32, (3, 3), activation='relu', input_shape=(32, 32, 3))
HIDDEN_LAYER_TYPE_1 = layers.MaxPooling2D((2, 2))
HIDDEN_LAYER_TYPE_2 = layers.Conv2D(64, (3, 3), activation='relu')

HIDDEN_LAYER_TYPE_3 = layers.MaxPooling2D((2, 2))

HIDDEN_LAYER_TYPE_4 = layers.Conv2D(64, (3, 3), activation='relu')

HIDDEN_LAYER_TYPE_5 = layers.Dense(64, activation='relu')

LAST_LAYER = layers.Dense(10)
\end{lstlisting}

For this assignment an optimizer with an adaptive learning rate often performs quite well and for the loss we want something that can help optimize for differences across categories. Again there are many possible options that will work. A possible solution is:

\begin{lstlisting}
OPTIMIZER = 'adam'

LOSS = tf.keras.losses.SparseCategoricalCrossentropy(from_logits=True)
\end{lstlisting}

As always you'll find the link to the assignment again below in case you want to explore it a bit more now that you have a solution.

\url{https://colab.research.google.com/github/tinyMLx/colabs/blob/master/2-3-9-AssignmentQuestion.ipynb}




\section{Recap}


Over the last few lessons you've gone from a single neuron to classify $Y=wX+b$ to having layers of these that could classify images, giving you basic computer vision.

You then learned about filters, and about how filters work to amend images:

\begin{figure}
  \includegraphics[width=\textwidth]{TinyML/1 Fundamentals/tempsnip_3.png}
    \caption{A screenshot of the Colab used in the last section showing how a convolutional filter applies a pre-defined set of filter values to a spatial patch of an image outputting a single value. An example is show for an image of a shoe and with pixel values and filter values and the mathematical computation down to compute the result. That math is shown to be the sum of the multiplication of filter and image values as discussed in the previous section.}
\end{figure}  



Filters could then be used to extract features from images. Here you can see a basic filter that removes most of the information from the image except for vertical lines:

\begin{figure}
  \includegraphics[width=\textwidth]{TinyML/1 Fundamentals/tempsnip_4.png}
    \caption{Another screenshot from the Colab used in the previous section showing an image of two people walking up some stairs that have a banister with vertical bars and then the result of applying a vertical line detetcting filter to that image. This makes the image highlight vertical lines in white and the rest turns black.}
\end{figure}  


When combined with pooling, an image like the one above could be compressed and could have the features enhanced.

\begin{figure}
  \includegraphics[width=\textwidth]{TinyML/1 Fundamentals/tempsnip_5.png}
    \caption{Another screenshot from the Colab used in the previous section showing the result of pooling applied to the vertical line highlighted image which compresses the image and highlights the lines more and reduces the total number of pixels.}
\end{figure}  



You explored adding convolutional layers with Fashion MNIST, and by extracting features, and then classifying those, instead of attempting to classify raw pixels, you achieved a better result. But these images were very simple -- 28x28 monochrome images containing a single subject, and it was centered. Can convolutions help with understanding the contents of more complex images?

The short answer is ``yes'', and you'll explore that with images that are more real-world in the coming lessons. As a taste of what’s possible there's a great paper here -- \url{https://arxiv.org/pdf/1311.2901v3.pdf} -- where the authors explore how a Convolutional Neural Network can ``see' features. So, on the left diagram here, you can see the results of a filter when applied to what are obviously dogs, as well as curved items like clocks. The images on the right are actually reconstructions of images using these filters -- where these images were generated by the neural network by going backwards through filters!

\begin{figure}
  \includegraphics[width=\textwidth]{TinyML/1 Fundamentals/tempsnip_6.png}
    \caption{A final screenshot from the above linked paper showing the results of visualizing features. This results in the images on the right of dogs and clocks and cans having the dog facses and curved sections of the cans highlighted}
\end{figure}  

.

That's obviously beyond the scope of what we’re doing here -- but it's interesting to see where the journey might take you! For some code on how to create visualizations of your filters check out: 

\URL{https://keras.io/examples/vision/visualizing_what_convnets_learn/lu')}


\section{Preparing image data}

Now that you've looked at using convolutions
to improve your classification of images or computer vision,
let's take a look at going beyond the simple MNIST images
to use more realistic content, larger color images where
the subject isn't necessarily centered.
Previously when using MNIST or Fashion-MNIST,
it was very simple for you to get the data.
You simply load it from the data sets and it gave you
a set of training images and labels and the same for validation images
and labels.
You did this just by saying \PYTHON{mnist=tf.keras.datasets.fashion\_mnist}.
And then you would load the data from it.
It was a nice luxury to have two lines of code to give you your data.
But it's not always that simple.
So let's take a look at another example.
One useful data set for learning is the Horses or Humans data
set, which has over 1,000 images that are 300 by 300
in full color of horses and humans in various poses.
Note that unlike MNIST and Fashion-MNIST,
they're not always centered in the frame and sometimes they're
not the only content.
So for example, there's background content like beaches and trees
and sky and stuff like that.
A methodology in TensorFlow to handle these
is to create directories containing the images.
And then, they'll be labeled based on their parent directory.
So your training set can be in one directory
with subdirectories for the images.
And the labels will be the names of those folders.
So we can have a training images for humans in a Humans folder.
And we can have training images for horses in Horses folder.
And similarly for validation, we can have folders for Horses and Humans.
And within them will have the validation images.
In TensorFlow, you can use the Image Data Generator to turn directories
like this into a trainable pipeline.
It's imported with code like this from tensorflow.keras.preprocessing.
The Image Data Generator can perform processing on the images.
So for example, we need to normalize our images, which we had done previously
by dividing by 255.
And we can now do that as an initialization on the Image Data
Generator.
Then, we can create a training generator by flowing
from the directory of images.
And you can tell it the path of the directory containing the training
images.
You can also specify the target size of the image.
So if they're of different sizes, the generator will resize them to this.
We'll design a neural network to take in 300 by 300 images.
And in the case of Horses or Humans, this is the default size.
But you'll see with other data sets like Cats versus Dogs,
the images are on all different shapes and sizes,
and you'll need a standard size before you can feed them into the network.
Training can be done in batches.
So we'll take batches of 128 images at a time.
Then, this final parameter is the class mode.
And there can be two different types--
binary, when there are two classes, and categorical,
when there are more than two classes.
This is Horses or Humans that has two classes,
so we'll set class mode to be binary.
Similarly, the validation data can have a generator made for it.
We point this at the directory containing the validation images.
Now we can explore our model architecture.
We have large images containing complex features.
So this network will certainly look much more sophisticated
than what you've seen already.
And we'll try it with four layers of convolutions,
and each has an associated pooling.
Our input shape is 300 by 300 by 3 because our images are 300 by 300
and they're in color, which has three channels.
And then finally, here at the bottom, we have our output layer.
We have two classes, which would normally have two output neurons.
But one other option they can do is to use a single neuron
and set the activation function to sigmoid.
This will pull the output value towards 0 for one class and towards 1
for the other.
It's another option that you can use when you have only two classes.
As always, we'll compile our model specifying
a loss function and an optimizer.
In this case, we can try a different optimizer, RMSprop,
which stands for Root Mean Square Propagation.
And it's always good to experiment with different ones.
Now when we do our model.fit, we can inform
it to use the generators for the training and validation data.
For the training data, we can specify the training generator
that we set up earlier.
We don't need to specify a parameter name.
And as we batched our data, we can specify
how many steps we want to take with each step loading a single batch.
We can pick the number of epochs to train across or just try things out.
We can say train for a few epochs, say 15.
For validation, we just passed the generator for the validation data
and it will perform the validation for us.
The validation data is also in batches, so we
can specify how many steps per epoch we want to take with the validation data.
After training for just 15 epochs, we can reach a pretty good accuracy
on both the training and validation sets.
This is a pretty small data set, so don't
be lulled into a false sense of security that it's 100% accurate
all Horse or Human images.
But it looks pretty good.
So next up, you can try the code for yourself.
Play with the parameters, add or remove layers,
try different values for neurons or different numbers of filters.
You might actually even be able to improve on the model
or maybe you can speed up the training.


\section{Coding exercise: complex images}

In this Colab you'll get to explore hands on the concepts discussed in the previous video. You’ll start by downloading and organizing the horses v. humans image database using Generators. You'll then train a model using your organized data and then get to upload your own horse or human image and see if the model can detect it correctly. Finally you'll get to visualize the model layers.

Note: as many learners have pointed out in the discussion below, the final model will only work well on images that are similar to those in our limited example dataset. As such, do not be surprised if you can trick the model!

\URL{https://colab.research.google.com/github/tinyMLx/colabs/blob/master/2-4-3-HorsesOrHumans.ipynb}



\section{TensorFlow Datasets (TFDS)  for image data}


In the previous examples you saw how to download a ZIP file containing images, and how, if the images are sorted into subdirectories, you could use a generator to automatically label the images based on their directories.

So, for example, in the horses-or-humans dataset, if you downloaded both the training and validation zip files, you’d end up with a structure like this:


\begin{figure}
  \includegraphics[width=\textwidth]{TinyML/1 Fundamentals/TDFS1}
    \caption{A diagram of the data that has been loaded into memory showing how the images are split into training and validation folders. In side of each of those folders is a horses folder and a human folder. In each of those folders are a series of images (e.g., 1.jpg, 2.jpg, 3.jpg).}
\end{figure}  



Within your \FILE{Images} directory, if you unzip the training files into a \FILE{training} folder, it would have subdirectories for \FILE{Horses} and \FILE{Humans}.

\begin{figure}
  \includegraphics[width=\textwidth]{TinyML/1 Fundamentals/TDFS2}
    \caption{The same diagram above highliting the training row of the depiction again showing the training has both horses and humans folders each with various images in them.}
\end{figure}  



By pointing an ImageDataGenerator at this subdirectory, it would automatically label images in the ‘Horses’ directory as horses, and similarly the images in the Humans directory would be labelled as such.

All of this requires the images to be sorted into named subdirectories, downloaded, and unzipped into the same structure. It works well for training, but it isn’t the way of flowing data into a neural network. And this is just one way of representing image data -- there are many -- and often you have to have some domain-specific knowledge of the data (at the very least the structure of a zip like this, but it’s often more complicated) and code around that to get the data into a way that you can train a neural network with.

\section{TensorFlow Datasets}

A common way to flow data into your network for training is TensorFlow Datasets.

The goal behind TensorFlow Datasets (TFDS) is to expose datasets in a way that’s easy to consume, where all the preprocessing steps of acquiring the data and getting it into TensorFlow-friendly APIs is done for you.

You’ve already seen a little of this idea with how Keras handled Fashion MNIST.

As a recap, all you had to do to get the data was this:

\begin{lstlisting}
data = tf.keras.datasets.fashion_mnist

(training_images, training_labels), (test_images, test_labels) = data.load_data()
\end{lstlisting}

TensorFlow Datasets builds on this idea, but greatly expands not only the number of datasets available but the diversity of dataset types. The \href{https://www.tensorflow.org/datasets/catalog/overview}{list of available datasets} is growing all the time including pictures, text, audio, video and more. Check out the link to see some of the datasets.

TensorFlow Datasets is a separate install from TensorFlow, so be sure to install it before trying out any samples! If you are using Google Colab, it’s already preinstalled.

If you need to install it, you can do so with a pip command:

\SHELL{pip install tensorflow-datasets}


Once it's installed, you can use it to get access to a dataset with \PYTHON{tfds.load}, passing it the name of the desired dataset. For example, if you want to use Fashion MNIST, you can use code like this:

\begin{lstlisting}
import tensorflow as tf
import tensorflow_datasets as tfds
mnist_data = tfds.load("fashion_mnist")
for item in mnist_data:
    print(item)
\end{lstlisting}

Two very important concepts to learn with TensorFlow Datasets are the Splits API -- that gives you a flexible way of splitting up data into Training, Testing, Validation sets, and Mapping Functions, which allow you to do things like Augmentation. So, for example, you saw how to do Image Augmentation on a generator, but if you’re no longer using generators you’ll need an alternative! TFDS makes it simple with the mapping functions.

Here's code as an example:

\begin{lstlisting}
def augmentimages(image, label):
    image = tf.cast(image, tf.float32)
    image = (image/255)
    image = tf.image.random_flip_left_right(image)
    return image, label

data = tfds.load('horses_or_humans', split='train', as_supervised=True)
train = data.map(augmentimages)
\end{lstlisting}

In this case, tfds.load is used to get the horses or humans dataset. It's pre-split into a ``train'' subset with the training data, so you can request that. Then, once you have data, you can call it’s ‘map’ method, passing it a function like ``AugmentImages'' as shown, and from within there you can do your image augmentation.

There's a lot to learn with TFDS, and it’s a really powerful API. Visit: \URL{https://www.tensorflow.org/datasets} to go deeper!

\section{Overfitting - Avoiding Overfitting}

LAURENCE MORONEY: The horses or humans data set is quite small.
So when you're training on it, you might get
lulled into a false sense of security that your network is
more accurate than it is at seeing and recognizing images.
This concept is often called overfitting.
And we'll explore that in this video.
It is really important to keep this in mind, in particular, on TinyML devices.
Because ML can be power hungry, we want to make
sure we're as efficient as possible in classifying.
And misclassifications, due to overfitting,
might really hurt your app.
Imagine if the only shoes you had ever seen in your life were hiking boots.
Now, even though hiking boots come in lots of different shapes, and sizes,
and colors, and models, by looking at many of them,
you'll soon think that all shoes were hiking boots.
And then, you could get very good at identifying them.
And as you're using the features that make up a hiking
boot to determine that it's a shoe, things
like the size or the location within the image wouldn't faze you.
You'd have no trouble with an image like this.
And you can see that this image definitely has shoes in it.
But then you see a shoe like this.
And you don't recognize it as a shoe.
Even though it really is one, you've overfed yourself
into thinking that all shoes look like hiking boots.
To explore a technique to avoid overfitting,
let's consider how convolutional neural networks work.
Recall that the spot features in an image, like the ears of a cat,
but if you only train on cat images that are upright,
then the computer might only recognize ears
that are oriented in that way as ears.
So it may not recognize the image on the right as having the ears of a cat.
They are oriented differently.
Indeed, if your network for spotting cats
was only trained on images like this, it might not think the image on the right
is a cat.
But what if we rotate the image a little before training.
Now, we'll have a labeled image of a cat with ears
that look like the cat on the right.
So when we train a network with images like the one on the left,
we'll better be able to recognize images like the one on the right.
Earlier, we saw the image data generator.
And we use it to rescale the image, in order to normalize it.
We did this by using a parameter to the image data generator.
The good news is there's lots of other parameters we can use.
And here's a list of some of them, including rotation range,
which will randomly rotate every image up to 40 degrees
left or right, as it's read from the file system.
By tweaking this parameter, you can impact how much the image is rotated.
Similarly, what if all of our images have
the subject in the center of the image, like the training image on the left?
It might not recognize the picture on the right as human in the same way.
She's much too far to the left in her frame.
Now, we can't rely solely on extracting features.
Because if the location of the feature is way off, then it might not work.
So shifting images around within the frames when training can help
us to prevent overfitting here too.
Now, this is achieved using the width shift range and the height shift range
parameters.
And when set to 0.2 as shown here, the image
will be randomly shifted up to 20% on width and height
when it is streamed off the desk for training.
Another augmentation that can be done is skewing.
This can be useful in a scenario like this one, where the image on the right
is clearly a human.
But our training data has nothing like it.
It does have this image like the one on the left, which is a similar pose.
And our human brain recognizes it.
But extracted features might miss it.
But if we skew the left image, we now have
pixels that look like the right one.
And again, we artificially enhance our data set with more scenarios.
So hopefully, our neural network will be better at recognizing them.
This type of sharing is achieved using the shear range parameter.
So we can set it to shear images up to 20% when training by setting it to 0.2.
One more scenario was this one.
The woman on the right is obviously human.
But her legs are not visible in the image.
The horses or humans data set has full body poses of each human.
And the image on the left looks very like the one on the right.
But the classification might fail, because of the missing legs.
But if we zoom in the image on the left, it now
looks a lot like the one on the right.
So while training, if we had samples that looked like this,
the image on the right might classify better in future.
Again, we'd avoid overfitting for humans that have legs
and avoid mistakenly misclassifying the one without legs as not human.
The Zoom range parameter achieves this.
By setting it to 0.2, every image will be zoomed up to 20% when training.
And another really obvious one is flipping images.
So say our training set only has humans with their right hand raised, like we
can see in the image on the left.
Then, it might not be able to categorize pictures
like the one on the right, where the woman has her left hand raised.
Again, we could overfit for right hand raisers mistakenly.
But if we flip the image in our training data,
we now have a left hand raiser too.
So we can avoid overfitting for this circumstance.
And as before, it's a parameter to the generator.
By setting it's true, images will be randomly flipped while training.
There's many more parameters to explore.
And hopefully this whet your appetite for what's possible.
Now, one last thing to consider as the fill mode.
This is useful for when the augmentation removes pixels.
For example, in the shearing that we saw earlier.
And you want to fill those pixels back in.
You pick a fill mode for this.
And there's a number of different options.
I usually use nearest, which picks neighbor pixels
and fills the gaps in using them.
So that's a quick tour of some of the techniques to avoid overfitting.
Next up, you'll have a Colab where you can try this out for yourself.
It's good to experiment.
Augmentation can help with overfitting.
But like anything else, use it in moderation.
Too much of it can really slow down your training.
Because it makes your images unrecognizable.

\section{Coding exercise: Image Augmentation}

In this Colab, we’ll again explore the horses v. humans Colab. This time, however, we will train with all forms of data augmentation turned on. You’ll see how this generates quite a bit more data (which takes longer to train) but produces a much more generalizable model since we are avoiding overfitting. You’ll get to test this by again having the chance to upload your own horse or human image and see if the model can detect it correctly. Finally you’ll again get to visualize the model layers.

As many students have noted in the discussion forum, since we are training on a very limited dataset, with a very small model, the final model does make a decent amount of mistakes. You'll find that if you test with images that look closer to the training dataset it will work much better. This is because while we are avoiding some overfitting with data augmentation, all machine learning models will struggle to generalize far beyond the training dataset!

\URL{https://colab.research.google.com/github/tinyMLx/colabs/blob/master/2-4-6-HorseOrHumanWithAugmentation.ipynb}


\section{Exploring dropouts - Dropout Regularization}

You've been exploring overfitting, where a network may become too specialized in a particular type of input data and fare poorly on others. One technique to help overcome this is use of dropout regularization

When a neural network is being trained, each individual neuron will have an effect on neurons in subsequent layers. Over time, particularly in larger networks, some neurons can become overspecialized—and that feeds downstream, potentially causing the network as a whole to become overspecialized and leading to overfitting. Additionally, neighboring neurons can end up with similar weights and biases, and if not monitored this can lead the overall model to become overspecialized to the features activated by those neurons.

For example, consider this neural network, where there are layers of 2, 5, 5, and 2 neurons. The neurons in the middle layers might end up with very similar weights and biases.

\begin{figure}
  \includegraphics[width=\textwidth]{TinyML/1 Fundamentals/NN}
    \caption{A depiction of a fully connected neural network showing a series of 4 layers of neurons the first and last with 2 neurons and the middle two with 5 neurons. All neurons in are connected to all of the other neurons in the layer before and after them with lines.}
\end{figure}  



While training, if you remove a random number of neurons and connections, and ignore them, their contribution to the neurons in the next layer are temporarily blocked

\begin{figure}
  \includegraphics[width=\textwidth]{TinyML/1 Fundamentals/NNDropout}
    \caption{A deptiction of applying dropout to the above neural network. Some of the neurons are marked in a different color. Those are either completely disconnected from the rest of the neurons in both directions or  only from the previous layer.}
\end{figure}  

A deptiction of applying dropout to the above neural network. Some of the neurons are marked in a different color. Those are either completely disconnected from the rest of the neurons in both directions or  only from the previous layer.

This reduces the chances of the neurons becoming overspecialized. The network will still learn the same number of parameters, but it should be better at generalization—that is, it should be more resilient to different inputs.

The concept of dropouts was proposed by Nitish Srivastava et al. in their 2014 paper \HREF{https://jmlr.csail.mit.edu/papers/volume15/srivastava14a/srivastava14a.pdf}{``Dropout: A Simple Way to Prevent Neural Networks from Overfitting''}.

To implement dropouts in TensorFlow, you can just use a simple Keras layer like this:

\PYTHON{tf.keras.layers.Dropout(0.2),}

This will drop out at random the specified percentage of neurons (here, 20\%) in the specified layer. Note that it may take some experimentation to find the correct percentage for your network.

For a simple example that demonstrates this, consider the Fashion MNIST classifier you explored earlier.

If you change the network definition to have a lot more layers, like this:

\begin{lstlisting}
model = tf.keras.models.Sequential([
          tf.keras.layers.Flatten(input_shape=(28,28)),
          tf.keras.layers.Dense(256, activation=tf.nn.relu),
          tf.keras.layers.Dense(128, activation=tf.nn.relu),
          tf.keras.layers.Dense(64, activation=tf.nn.relu),
          tf.keras.layers.Dense(10, activation=tf.nn.softmax)])
\end{lstlisting}

Training this for 20 epochs gave around 94\% accuracy on the training set, and about 88.5\% on the validation set. This is a sign of potential overfitting.

Introducing dropouts after each dense layer looks like this:

\begin{lstlisting}
model = tf.keras.models.Sequential([
           tf.keras.layers.Flatten(input_shape=(28,28)),
           tf.keras.layers.Dense(256, activation=tf.nn.relu),
           tf.keras.layers.Dropout(0.2),
           tf.keras.layers.Dense(128, activation=tf.nn.relu),
           tf.keras.layers.Dropout(0.2),
           tf.keras.layers.Dense(64, activation=tf.nn.relu),
           tf.keras.layers.Dropout(0.2),
           tf.keras.layers.Dense(10, activation=tf.nn.softmax)])
\end{lstlisting}

When this network was trained for the same period on the same data, the accuracy on the training set dropped to about 89.5\%. The accuracy on the validation set stayed about the same, at 88.3\%. These values are much closer to each other; the introduction of dropouts thus not only demonstrated that overfitting was occurring, but also that adding them can help remove it by ensuring that the network isn’t overspecializing to the training data.

Keep in mind as you design your neural networks that great results on your training set are not always a good thing. This could be a sign of overfitting. Introducing dropouts can help you remove that problem, so that you can optimize your network in other areas without that false sense of security!


\section{Explore other approaches - Exploring Loss Functions and Optimizers}

To this point, you’ve been largely guided through different loss functions and optimizers that you can use when training a network.

It’s good to explore them for yourself, as well as understanding how to declare them in TensorFlow, particularly those that can accept parameters!

Note that there are generally 2 ways that you can declare these functions -- by name, in a string literal, or by object, by defining the class name of the function you want to use.

Here’s an example of doing it by name:

\begin{lstlisting}
optimizer = 'adam'
\end{lstlisting}

And one of doing it using the functional syntax

\begin{lstlisting}
from tensorflow.keras.optimizers import Adam
opt = Adam(learning_rate=0.001)

optimizer = opt
\end{lstlisting}

Using the former method is obviously quicker and easier, and you don’t need any imports, which can be easy to forget, in particular if you’re copying and pasting code from elsewhere! Using the latter has the distinct advantage of letting you set internal hyperparameters, such as the learning rate, giving you more fine-grained control over how your network learns.

You can learn more about the suite of optimizers in TensorFlow at \URL{https://www.tensorflow.org/api_docs/python/tf/keras/optimizers} -- to this point you’ve seen SGD, RMSProp and Adam, and I’d recommend you read up on what they do. After that, consider reading into some of the others, in particular the enhancements to the Adam algorithm that are available.

Similarly you can learn about the loss functions in TensorFlow at \URL{https://www.tensorflow.org/api_docs/python/tf/keras/losses}, and to this point you’ve seen Mean Squared Error, Binary CrossEntropy and Categorical CrossEntropy. Read into them to see how they work, and also look into some of the others that are enhancements to these.

When you're done, go back to some of the exercises or other colabs that you’ve done, and experiment with different loss functions or optimizers. In particular tweak things like the learning rate or other parameters to see if you can improve your models!

\section{Coding assignment: enhancing a CNN}


For this assignment you'll take what you've learned so far and build a classifier for bean disease. You'll be provided with training and validation data based on images taken of bean plants in Uganda. These images show healthy bean leaves as well as 2 types of common disease. Your job will be to build a neural network that can tell the difference between the healthy and diseased leaves.

\URL{https://colab.research.google.com/github/tinyMLx/colabs/blob/master/2-4-11-AssignmentQuestion.ipynb}


\section{Assignment solution}

We hope you enjoyed using Generators to help develop a model to solve a real world problem!

As usual, there are many possible solutions to this assignment.

For the generators, at a minimum you should have had both the training and validation Generators rescale the image data. This can be done using the training Generator as an example:

\begin{lstlisting}
train_datagen = ImageDataGenerator(
                   rescale=1./255
                )
\end{lstlisting}

If you wanted you could have gone further and applied additional image augmentation techniques to the training Generator to have a better dataset for training. A possible solution could be:

\begin{lstlisting}
train_datagen = ImageDataGenerator(
                     rescale=1./255,
                     rotation_range=40,
                     width_shift_range=0.2,
                     height_shift_range=0.2,
                     shear_range=0.2,
                     zoom_range=0.2,
                     horizontal_flip=True,
                     fill_mode='nearest'
                    )
\end{lstlisting}

You then needed to supply the following values to get the Generators to work:

\begin{lstlisting}
TRAIN_DIRECTORY_LOCATION = '/tmp/train'
VAL_DIRECTORY_LOCATION = '/tmp/validation'

TARGET_SIZE = (224,224)
CLASS_MODE = 'categorical'
\end{lstlisting}

For the model, again there isn’t one answer. One good approach (as suggested by the hints), is again to use convolutional and pooling layers to find high level features and then dense layers to classify those features. A possible solution to the model design is shown below:

\begin{lstlisting}
model = tf.keras.models.Sequential([
    # Find the features with Convolutions and Pooling
    tf.keras.layers.Conv2D(16, (3,3), activation='relu', input_shape=(224, 224, 3)),
    tf.keras.layers.MaxPooling2D(2, 2),
    tf.keras.layers.Conv2D(32, (3,3), activation='relu'),
    tf.keras.layers.MaxPooling2D(2,2),
    tf.keras.layers.Conv2D(64, (3,3), activation='relu'),
    tf.keras.layers.MaxPooling2D(2,2),
    tf.keras.layers.Conv2D(128, (3,3), activation='relu'),
    tf.keras.layers.MaxPooling2D(2,2),
    # Flatten the results to feed into a DNN
    tf.keras.layers.Flatten(),
    # 512 neuron hidden layer
    tf.keras.layers.Dense(512, activation='relu'),
    tf.keras.layers.Dense(3, activation='softmax')
  ])
\end{lstlisting}

For this assignment an optimizer with an adaptive learning rate often performs quite well and for the loss we want something that can help optimize for differences across categories. Again there are many possible options that will work. A possible solution is:

\begin{lstlisting}
OPTIMIZER = 'adam'
LOSS_FUNCTION = 'categorical_crossentropy'
\end{lstlisting}

Another optimizer we got to work well was:

\begin{lstlisting}
from tensorflow.keras.optimizers import RMSprop
OPTIMIZER = RMSprop(lr=0.0001)
\end{lstlisting}

As for Epochs as mentioned in the hint, something in the low tens is a good place to start. We found that as low as 20 Epochs worked well.

As always you’ll find the link to the assignment again below in case you want to explore it a bit more now that you have a solution.

\url{https://colab.research.google.com/github/tinyMLx/colabs/blob/master/2-4-11-AssignmentQuestion.ipynb}


\section{What am I building and what's the goal?}


SUSAN KENNEDY: Hi my name is Susan Kennedy
and I'll be your guest lecture on responsible AI design.
I have a PhD in philosophy with a specialization in ethics,
and my current research focuses on the ethical, social, and political impacts
of emerging technologies.
I'm here to talk to you about how you can proactively incorporate ethics
into your workflow so that you can pursue innovative applications of AI
in a responsible manner.
During your introduction to responsible AI,
you were shown the human-centered design framework.
Following this framework, we're now going
to focus on the ethical challenges that arise
during the first stage of innovation, the design phase.
The specific questions we'll think about are, what am I building
and what is the goal?
Currently, the hype surrounding AI is at an all time high
as it looks like a promising avenue to tackle
some of the world's biggest challenges.
As Nick Bostrom, a philosopher at the University of Oxford, has said,
machine intelligence is the last invention
that humanity will ever need to make.
Without a doubt, the capabilities of AI are promising.
But with that being said, as a designer, it's important to keep in mind
the AI is not always the best solution.
Earlier in this course, Lawrence explained the difference
between traditional programming and machine
learning in terms of how they operate.
To be reminded, with traditional programming,
we combine data inputs with rules to get answers, whereas with machine learning,
we combine data with answers so that our model can generate the rules.
Now we're going to take a closer look at the pros and cons of each
so that we can make a more informed decision about when
we should choose a machine learning model as a solution to some problem.
In some scenarios, traditional programming
will work just as well if not better than machine learning.
It's particularly suited to situations in which there
are lower volumes of data and the rules are relatively simple.
Some pros of traditional programming is that it's quicker
to build, easier to explain, easier to debug, easier to maintain,
and offers more consistent performance.
Some downsides of traditional programming is that it does not scale,
it does not adapt or update on its own, and it's not worked well for complex
tasks that would require a huge amount of rules or rules that domain experts
cannot specify.
In contrast, machine learning offers some important benefits
over traditional programming.
It can be used to solve complex problems,
it can scale adapt and update on its own,
and offer a personalized user experience, and improve over time.
Some downsides of machine learning models is that they're slower to build,
as they require the collection of vast amounts of data and time
to train the model.
They're harder to interpret and explain, also known as the Black Box Problem,
and relatedly, they're harder to debug for this reason.
Now that we've seen some trade offs between traditional programming
and machine learning, we want to keep this in mind
when we think about what our goal is.
The most comprehensive way to think about what our goal is
to answer these four questions.
What is the problem?
Why does it need to be solved?
How could the problem be solved manually with traditional programming?
And will AI offer unique value?
Let's look at what the answers to these questions
might look like by considering a case study
you watched a video about earlier, where AI
was used for the diagnosis of diabetic retinopathy or DR for short.
The problem is achieving an early diagnosis of DR progression
from color fundus photographs of the retina.
The problem needs to be solved because DR is reaching epidemic proportions,
and early diagnosis would allow for intervention and prevention
of further DR progression.
The problem could be solved manually if we collected data
from retina specialists or trained graders
and then wrote rules based on the presence of hemorrhages
and micro aneurysms.
But will AI offer unique value?
Seems like the answer is, yes.
Machine learning can improve predictions about the progression
of DR for individual patients and identify patterns and images that were
previously unrecognized by experts.
When you design your own applications, you
want to answer these questions to make sure you have a clearly defined goal
and you choose the right model, either machine
learning or traditional programming, to pursue this goal.
A final point worth taking into consideration
is how you can define your goals for positive social impact.
This is commonly referred to as AI for social good
by experts in both industry and academia, who
recognize the importance of using AI and machine learning to help address
the world's most pressing challenges.
For example, AI has been used for image classification
to spot elephant poachers and for sound detection
to stop illegal deforestation.
It's easy to see how these innovations are making the world a better place,
but you may be wondering, how would I be able to identify problems
at the societal level that need to be addressed?
The good news is the United Nations has put together a list of 17 goals
as part of their 2030 agenda for sustainable development,
and it spans issues ranging from ending poverty and world hunger,
improving health and education, and reducing inequality and injustice.
This will help you identify problems you care about
so you can spend your time creating a solution.
Now there are some challenges facing AI for social good,
such as learning with limited data, geographical imbalances,
biased data, and learning with limited memory and computation.
But now there's a window of opportunity thanks
to cutting edge innovations in TinyML.
it's pushing the boundaries for learning with limited memory and computation
and has other unique affordances like being
battery-operated, on-device computing, low latency, low cost, and small size.
Just take a look at this image of the world's smallest microprocessor that
fits on the dimple of a golf ball.
It will now be possible to explore the benefits of AI
in domains that were previously inaccessible.
We can sprinkle sensors across the forest to aid conservation efforts
or embed TinyML in everyday household items to make our lives easier.
With such exciting and new possibilities,
we can't wait to see what you will create.
Head into the discussion forum and let us
know how you imagine TinyML being designed for social good.

\section{Forum: Development and TinyML}

  \includegraphics[width=\textwidth]{TinyML/1 Fundamentals/Goals}


As part of the United Nations’ 2030 Agenda for Sustainable Development, a list of 17 goals has been put together. Read through the list of goals below:

\begin{enumerate}
  \item No poverty 
  \item Zero hunger 
  \item Good health and well being 
  \item Quality education 
  \item Gender equality 
  \item Clean water and sanitation 
  \item Affordable and clean energy 
  \item Decent work and economic growth  
  \item Industry, innovation and infrastructure 
  \item Reduced inequalities  
  \item Sustainable cities and communities 
  \item Responsible consumption and production  
  \item Climate action
  \item Life below water
  \item Life on land
  \item Peace, justice, and strong institutions
  \item Partnerships
\end{enumerate}


\begin{itemize}
  \item Which goal do you think is most important? Why?
  \item How do you think TinyML could be used to achieve this goal?
  \item What ideas have others posted that you are excited about?
\end{itemize}


\section{Who am I building this for?}

SUSAN KENNEDY: Building off of our human-centered design framework,
we want to ask ourselves, who am I building this for
and how are they going to use it?
If we think about these questions in the design phase,
we can ensure better uptake in our product
and identify potential problems so that we can improve our overall design
strategy.
The people who will potentially be impacted by the technology
are referred to as the stakeholders.
The direct stakeholders are the people who
will use your technology otherwise known as the users.
It's important to be specific when you consider the users of your technology.
Try to envision the characteristics that describe your target users like age,
gender, whether they're able-bodied or perhaps
have some disability, where they're located, and so forth.
Additionally, your technology may have indirect stakeholders.
And these are the people who are not users but are nonetheless
impacted by your technology.
To see what this might mean, suppose we're developing an image
classification for disease diagnosis.
The direct stakeholders would be the doctors
who use this technology to assist them in making a diagnosis.
The indirect stakeholders would be the patients who are being diagnosed.
Of course, there may be many different individuals or groups
who might be indirectly affected.
So be sure to think carefully about who your technology might impact.
Once you have a clear description of your stakeholders.
It will be easier to identify what values they might have.
Let's return to our example of an image classifier for disease diagnosis.
Suppose the doctor values the accuracy of this technology
so they can provide the best care for their patients.
They also care about whether they'll need additional training to utilize it,
how easy it will be to incorporate in their day-to-day workflow,
and whether it will allow for advances to be made in medical research.
The patient, on the other hand, values a personal level of care,
being fully informed about their diagnosis,
being able to trust their doctor, and ensuring
that their privacy is protected.
Now that we've identified the stakeholders' values,
you've probably noticed that the values for the doctor
look very different from those of the patient.
This is important because it helps us see where potential value tensions may
arise.
For example, the doctor may want the most accurate AI available.
But it's possible that the algorithm's decision-making process
will be opaque in such a way that the doctor can't explain
how it arrived at a certain diagnosis.
And this might interfere with the patient's desire to be fully informed.
Relatedly, this might interfere with the patient's ability
to trust that their doctor is making the right decision about their diagnosis.
Let's consider a different value tension.
The doctor values the advances in research
that would be made possible by the AI's collection of data.
But this collection of data might conflict
with the patient's desire for privacy.
So how can we resolve value tensions once we've identified them?
There are three questions we can ask.
Is one of the conflicting values considered a right?
If so, we should be careful to follow laws and regulations that
protect this right.
Do the stakeholders prioritize certain values?
Maybe the doctor values the AI's accuracy so much
that they'd be willing to undergo additional training that's
necessary to utilize it.
And does the value have a threshold for sufficiency?
For example, our patient's desire for privacy
might be satisfied if we employ deidentification techniques
on the personal data that's collected.
While value tensions require some critical thinking,
there's oftentimes a way to resolve them if we keep an open mind.
There are some other important factors we'll want
to think about during the design phase.
For instance, you'll want to consider how this technology will
look when analyzed from diverse perspectives, particularly
those that are historically marginalized or underrepresented.
Suppose we're designing a voice assistant
to respond to speech commands.
How would the choice of command phrases look
to people who are non-native English speakers, people
who speak African-American vernacular English, or people
with speech impairments?
Another thing to consider is the social context,
in particular what features of the social context
are likely to impact access, adoption, and use of the technology we're
designing?
Suppose we're designing a wearable health device.
And our intended users span different countries.
Will there be a digital divide where some people have restricted access
to this technology?
Do the countries have different rules and regulations for data protection?
And do customs impact how users prioritize values?
Once we have a picture of who we're building for,
we should spend some time thinking about how we're going to use it.
One of the earliest philosophers of technology, Don Ihde,
explained that technology can end up being used in ways that it
wasn't originally designed for.
He coined the term for this phenomenon as the multistability of technology.
Multistability matters from an ethical perspective.
Because even if we design technology for some specific purpose,
it may be compatible with malicious, immoral, or simply unintended purposes.
We can see the problems surrounding multistability
in these high-profile cases that made news headlines.
For example, Google said it would hold off
on selling facial recognition technology because it
could be used for inappropriate forms of surveillance.
And the research lab OpenAI said it would not
release its fake text generator GPT-2 because of worries
that it would result in the spread of fake news.
In these cases, worries about malicious and immoral use
resulted in the companies making a decision
to ban the technology by refusing to develop it
or after developing it refusing to deploy it.
But not all instances of nonintended use will play out this way.
It's important to anticipate the multistability of the technology we're
developing so we can proactively address the issue in the design phase.
If it's potentially very harmful like the cases we just saw,
we may need to stop developing it altogether.
But in other cases, it's possible that we just
need to implement some design changes to prevent problematic uses.
Lastly, there might even be some cases where
we should embrace nonintended use like if we discover
that the sensors we built to help athletes
could also be used to improve physical therapy treatment.
We may decide that we should redesign our technology for that purpose
instead.
If you keep in mind the stakeholders and potential uses of your technology
early on in the design phase, you'll be one step closer
to creating better, more successful products.


\section{User consequences - What are the consequences for when the user when it fails}


PROFESSOR: Revisiting our human-centered design framework,
the last question we want to explore in the design phase of AI
is, what are the consequences for the user when it fails?
It's important to keep in mind, AI isn't perfect.
Just look at this graph here that shows the accuracy of popular voice
assistants in terms of answering speech commands.
Google Assistant scored 86%, Siri was 93%, and Alexa 69%.
Even though they aren't perfect, these voice assistants
are still widespread and considered useful by many.
While perfect AI may not be possible or even
necessary to have a successful product, we'll
want to think carefully about what the consequences of our AI
will be for the user.
Suppose we develop image classification to predict whether a plant is diseased
or not.
When AI succeeds, it will correctly predict disease and cases
where the plant actually is diseased, a true positive, located
on the upper left.
It will also correctly predict the absence of disease
and cases where the plant is healthy, a true negative,
located in the lower right.
Now let's look at what happens when AI fails.
If our AI predicts disease when the plant actually
isn't, this incorrect prediction would be considered a false positive, located
in the upper right.
If it predicts no disease when the plant actually is,
this is a false negative, located in the lower left.
It's important to distinguish between these cases.
And that's why we specify a false positive as a type I
error and a false negative as a type II error.
The reason it's important to distinguish between these two types of error
is because it makes a difference for how we determine which one of the two ways
we should optimize our model.
If we optimize for high precision, we will minimize type I errors
and have fewer false positives, and this is good
when the costs of false positives are high.
But this comes at the price of increasing type II errors
and having more false negatives.
If we optimize for high recall, on the other hand,
we can minimize type II errors and have fewer false negatives,
and this is good when the cost of detections is high.
But again, this comes at a price where we increase type I errors,
and so we'll have more false positives.
We can visualize the tension between optimizing for precision
and recall with this graph.
Let's suppose we build a model to filter out spam email,
and this is the output of our linear regression model.
Actual spam is a green circle, false positives are gray circles,
actually not spam is a red X, and false negatives are gray Xs.
If we look at the chart on the left, we can
optimize for precision by increasing our classification threshold,
so less false positives but more false negatives.
If we look at the graph on the right, we've
optimized for recall by decreasing the classification threshold,
and now have less false negatives but more false positives.
So which measure is more important?
Well, the truth is, it depends.
We need to discover the right answer on a case by case basis.
But not to worry.
There are four key questions we can ask to figure out the right answer.
First, we'll identify the direct and indirect stakeholders.
Then we'll list the consequences of type I errors and we'll
do the same for type II errors, and, finally, compare our answers
and see which type of error would be most harmful.
If type I errors or false alarms are most harmful,
then precision is more important.
If type II errors are most harmful, then recall is more important.
Now let's revisit our example of developing a model for spam filtering.
The direct stakeholder is the user with the email account
and the indirect stakeholder will be the user's correspondents.
If our model makes a type I error, the user won't receive important mail
and her correspondents will be prevented from communicating with her.
If our model makes a type II error, the user will see spam in her inbox
and perhaps be slightly annoyed by the inconvenience.
In this case, it seems like type I error is our most harmful,
so we'll want to optimize our model for higher precision.
Let's consider another example of a smart smoke detector.
The direct stakeholder is the user, and indirect stakeholders
might be other members of the household and neighbors.
A type I error might result in the user and household members
being annoyed by noisy false alarms, but a type II error
might result in people failing to evacuate
during a fire and neighbors having less time to respond as the fire spreads.
Type II errors are most harmful in this case.
The best thing to do seems to be optimizing for higher recall.
You might be wondering, can't we take a shortcut
and simply generalize that all alarm systems should
optimize for high recall, for example?
While this might be easier, it just doesn't work that way.
Suppose we have a home surveillance security system.
It'll be the same application with the same classification
model that detects the sound of breaking glass, for instance.
The only difference will be that the feedback mechanism is different.
In one case, the user gets a text alert on their phone.
In the second case, the user gets a text,
but the police are immediately contacted as well.
In case 1, after going through our questions,
we determine type II errors are most harmful since the false alarm will
at most be considered a minor annoyance, so we
choose to optimize for higher recall.
But in the second case, where the police are notified,
a false alarm is no longer a minor nuisance
but a hazardous waste of valuable police resources.
Since our estimation of type I error is changed
because of different indirect stakeholders,
in this case, namely, the police, we get a different result,
where slightly higher precision might be preferred.
Of course, these cases are just short, simplified examples to help
illustrate the process.
In the real world, our cases might be significantly more complex
and require a more thorough analysis.
Now you can head into the discussion forum
and practice trying to work through these questions for a more
complicated case, and compare and discuss
your analysis with your fellow students in the course.

\section{Forum: Error types and ethics}

Imagine you are designing a voice assistant device to power on when it hears the designated wake word. The voice assistant is intended to be used in a person’s home. When the device recognizes the wake word, it will record audio for 20 seconds and send this data to the cloud. The direct stakeholder would be the user of this device, while the indirect stakeholders would be the other members of the household. 

Your task is to decide whether this voice assistant device should be optimized for high precision or high recall. To be reminded, you will want to consider what the consequences of type 1 errors and type 2 errors will be for both direct and indirect stakeholders. If you think type 1 errors are most harmful, you will want to optimize for high precision. If you think type 2 errors are most harmful, you will want to optimize for high recall.

\begin{itemize}
  \item Should you optimize this device for high precision or high recall? Please provide an explanation for your answer (i.e. which type of error you consider to be most harmful and which stakeholder(s) would be affected)
  \item Do you agree or disagree with what others have said in their posts?
\end{itemize}



\section{Recapping (Tiny) ML and Its Data-centric Role}

\subsection{Recapping (Tiny) ML and Its Data-centric Computing Role}
In Course 1, we have covered many of the fundamental topics of applied machine learning. We also looked at the broader computing ecosystem through the lens of an embedded platform where TinyML will be deployed. We start with these basics because they lay the groundwork for building TinyML applications. 

Hopefully, the power of supervised machine learning is now apparent: supervised machine learning provides us with the ability to develop data-driven function approximations. These approximations can be applied to essentially any data structure: images, sounds, sensor outputs, etc. As a primer for TinyML, it is helpful to understand some of the historical context of machine learning.

\subsection{How Do Today's Computers Compare to the Past?}

Machine learning has been around for decades. In fact, most of the algorithms in use today were developed in the 1980s and 1990s. Why were these algorithms not in widespread use at these times? The main bottleneck was a lack of available data and computing power. During these periods, floppy disks with a capacity of 1 MB were commonplace, and a state-of-the-art personal computer might have had the processing capacity of an Arduino. Given the space- and time-complexity of these algorithms, it was largely infeasible to run them on anything but small sets of data.

Over the course of time, improvements in processing and storage technologies, as well as rapid reductions in the cost of memory, have significantly reduced the barriers for performing machine learning. In the 2010s, it has become feasible to run machine learning models on local machines using one or more cores within the central processing unit (CPU). The demonstrable power of these techniques on larger datasets became apparent following the superior performance of AlexNet in the ImageNet Large Scale Visual Recognition Challenge in 2012.


\begin{figure}
    \GRAPHICSC{1.0}{1.0}{TinyML/1 Fundamentals/Arduino}
    \caption{An Arudino with an ARM Cortex M4 processor}
\end{figure}  

%An Arudino with an ARM Cortex M4 processor in 2020 that is 150,000x faster than a 1946 ENIAC supercomputer, 60x faster than a 1990's laptop and 4x faster than a 2007 iPhone. Showing typical compute specifications of 1-400MHz processor, 2-512 KB of main memory, 32KB-2MB of flash storage and 160uWatts-23.5mWats of power.

Shortly after, computation using graphics processing units (GPUs) became necessary to handle larger datasets and became more readily available due to introduction of cloud-based services such as SaaS platforms (e.g., Google Colaboratory) and IaaS (e.g., Amazon EC2 Instances). At this time, algorithms could still be run on single machines. Cloud-based services provide a convenient repository to localize data storage close to vast computing resources, a paradigm commonly referred to as the compute-centric paradigm.


\begin{figure}
    \GRAPHICSC{1.0}{1.0}{TinyML/1 Fundamentals/GPU}
    \caption{A picture of a NVIDIA Tesla GPU which is an example of the kinds of GPUs that unleash deep learning as well as a CloudTPU from google which represents th ekind of custom ML hardware being desgined by industry to meet their demands.}
\end{figure}  



More recently, we have seen the development of specialized application-specific integrated circuits (ASICs) and tensor processing units (TPUs), which can pack the power of ~8 GPUs and are readily available via cloud-based services. These devices have been augmented with the ability to distribute learning across multiple systems in an attempt to grow larger and larger models.

\subsection{Evolution of the Machine Learning Algorithms}

During 2020, we have seen the release of Turing-NLG and GPT-3 natural language models, the two largest machine learning models ever created. Natural language processing (NLP), or the ML task that focuses on teaching machines to handle human language is remarkably demanding. In March 2020, Turing-NLG was released, boasting a model size of 17 billion parameters, 10x larger than GPT-2 (at the time the largest model in existence). The insatiable demand for NLP model complexity is driven by the desire to enable a wide variety of applications such as question and answering, summarizing text, improving user experience with personal assistants such as Alexa, generating text to aid in sentence completion and so forth. 

The record set by Turing-NLG was short-lived. In May 2020, GPT-3 was released, a model with a whopping 175 billion parameters. Some estimates claim that GPT-3 cost around \$10 million dollars to train and used approximately 3 GWh of electricity (approximately the output of three nuclear power plants for an hour). Nowadays, training a single machine learning model can emit as much carbon as five cars. The GPT-3 model size is 45 TB, far beyond the capabilities of any single machine.

A graph showing the number of parameters in millions used in ML models agains time. This shows that from 2018-202 there has been exponential graph in ML model size. Starting in 2018 we see ELMO at 94, GPT at 110, VERT-Large at 340 then in 2019 we get to Transformer ELMo at 465, GPT-2 at 15000 MT\_DNN at 330, Grover-Mega at 1,500, MegatronicM at 8,300 and finally in 2020 GPT-3 at 175,000. These models come from Ai2, OpenAI, Google-AI, Microsoft, Nnidia, and the University of Washington.

\begin{figure}
  \includegraphics[width=\textwidth]{TinyML/1 Fundamentals/number_of_parameters}
    \caption{The Role of TinyML}
\end{figure}  



So where does TinyML come into the equation? The key to this answer is in the data, so to speak. The data used in machine learning applications are derived from a data source. This is often a camera, microphone, or some form of sensor tasked with capturing information about the physical world and transducing it into a format digestible by computers. For the past few years, there has been a surge of such devices connected to the cloud, colloquially known as the internet of things (IoT). 

These IoT devices periodically send data over the internet to the cloud and are often stored in a data warehouse which can easily be interfaced with via a cloud instance. Such warehousing of data from heterogeneously distributed devices requires a large amount of data to be transmitted through networking protocols from the IoT devices to the cloud. Some individuals have raised concerns with this infrastructure, namely related to (1) privacy, (2) latency, (3) storage, and (4) energy efficiency.

\begin{description}
  \item [Privacy.] Transmitting data creates potential for privacy violations. Such data could be intercepted by a malicious actor and becomes inherently less secure when warehoused in a singular location (such as a data warehouse).

  \item [Latency.] For standard IoT devices, such as Amazon Alexa, these devices transmit data to the cloud for processing and then return a response based on the algorithm’s output. In this sense, the device is just a convenient gateway to a cloud model, like a carrier pigeon between yourself and Amazon’s servers. The device is pretty “dumb” and fully dependent on the speed of the internet to produce a result. If you have a slow internet connection, Amazon Alexa will also become slow.

  \item [Storage.] For many IoT devices, the data they are obtaining and transmitting is of no merit. Imagine a security camera recording the entrance to a building for 24 hours a day. For a large portion of the day, the camera footage is of no utility, because nothing is happening.

  \item [Energy Efficiency.] Transmitting data (via wires or wirelessly) is energy-intensive, around an order of magnitude more so than onboard computations (specifically, multiply-accumulate units).
\end{description}


By empowering embedded systems with the ability to perform on-device machine learning, these problems are largely solved. IoT systems that can perform their own data processing require no data transmission, and are thus highly energy-efficient. Such devices would have minimal latency because there is reduced (if any) dependence on external communications. By keeping data primarily on the device and minimizing communications, security and privacy is improved. Additionally, for an intelligent system that only activates when necessary, lower storage capacity and fewer external communications are required. The industry term for such an intelligent IoT device is an edge device (devices at the “edge” of the cloud).

\subsection{The Future of TinyML}

As IoT devices become increasingly ubiquitous, the above-mentioned concerns are exacerbated. Billions of IoT devices are already in use, ranging from smart fridges and clocks, to thermostats, alarm systems, locks, and even children’s toys. At the start of 2020 there were roughly 8 billion active IoT devices. By 2030, it is anticipated that we will have nearly 50 billion IoT devices. 

\begin{figure}
  \includegraphics[width=\textwidth]{TinyML/1 Fundamentals/number_of_embedded}
    \caption{A graph showing the number of embedded IoT devices worldwide growing from 22 in 2018 to a projected 38.6 in 2025 to 50 in 2030.}
\end{figure}  

A graph showing the number of embedded IoT devices worldwide growing from 22 in 2018 to a projected 38.6 in 2025 to 50 in 2030.

Too many devices communicating over our networks simultaneously can quickly saturate the bandwidth, and present many attack vectors for hackers and malicious actors. These are some of the factors that have stimulated interest in an alternative computing paradigm for edge devices, the data-centric paradigm. This paradigm flips the compute-centric paradigm on its head, taking the computing power to the data, instead of the data to the computing power.

This brings us to a similar situation with that of the algorithmic developers of the 1980s and 1990s - we have our algorithms but are constrained by the hardware characteristics of microcontroller platforms, which rarely have flash memory or static RAM above 1 MB. The hierarchy of the computing system is shown below and there are vast differences (that we will learn a lot more about in the next courses). However, this situation is not exactly the same. Nowadays, we have access to our large scale models, they are just too big!

A figure showing the difference between tradition computing systems as compared to embeed ML systems as pyramids. Traditional systems have Flash (64GB), DRAM (4GB) off chip and L2 (4MB), L1 (48kb), and a Cortex-A processor on chip. The embedded ML systems instead rely on Cloud as the basis of the pyramid followed by on chip base of SRAM(320kb) and eFlash (1MB) with Cache (4KB) and a Cortex-M processor built on top.


\begin{figure}
  \includegraphics[width=\textwidth]{TinyML/1 Fundamentals/Copy_of_MCU-memory-Hierarchy__2_}
    \caption{Traditional computing systems versus embedded machine learning systems.}
\end{figure}  



The models used on edge devices need not necessarily be trained on edge devices - they can be trained in the cloud and subsequently ported to the edge device. The edge device need only perform model inference, which is significantly cheaper computationally. In many ways, running machine learning on an embedded device is similar to any production-ready machine learning model.

However, this is by no means a simple feat to achieve. Most microcontrollers do not have an operating system, and thus bare-metal and lightweight inference libraries (e.g., TFLite Micro, TVM, CMSIS-NN) are required, with essentially all unnecessary functionality (including debugging) removed. Similarly, model weights and computational graphs must be simplified to reduce their memory and computational overhead. This involves techniques such as model distillation, pruning, and quantization. This is the essence of TinyML.

Myriads of machine learning applications exist for edge devices, which are often located in remote areas which have sparse access to power and minimal networking capabilities. Edge devices constantly communicating via wireless transmitters consume approximately 10mW, whereas devices running TinyML models and only communicating when something of interest happens can function on less than 1mW, allowing them to run on a small coin battery continuously for more than a year. 

Combined with their low-cost, this presents countless opportunities for industrial applications. Sensor networks to monitor air quality throughout cities could be developed at relatively low cost and provide widespread coverage, focusing communication only in conditions of poor air quality. Similarly, industrial processes could be monitored for preventative maintenance purposes, reducing costs and mitigating hazards. Sensors could be placed along long-distance oil and gas pipelines to monitor for leaks or potential blockages. Intelligent wearables could allow end users to gain a better understanding of their own body, or to mitigate or solve medical problems such as neural communications or muscle atrophy.

We are at an interesting crossroads where machine learning is bifurcating between the two computing paradigms of compute-centric and data-centric computing. Although we appear to be quickly moving towards a ceiling in the compute-centric paradigm, it is clear that work in the data-centric paradigm has only just begun. 


\section{Why we are excited about TinyML}

VIJAY JANAPA REDDI: Welcome back, my friends.
It's good to see you again.
I hope you learned a good deal about the fundamentals of machine
learning from Laurence.
Now I'm here to tie things up in Course 1
and prepare you for what's coming next.
Now as you may recall, the goals of Course 1 were twofold.
First and foremost, we wanted to introduce the concept of TinyML.
Then we wanted to make sure that we all spoke
the same language of machine learning.
Now, different people might have different backgrounds.
Some of you might know a little bit of embedded systems.
Some of you might know a little bit about machine learning,
and we wanted to level the playing field.
So hopefully at this point you and I and everybody else taking this course right
now are all on the exact same page.
Now let's take a little step back and kind of recall why, for instance, we're
all so excited about TinyML.
Machine learning has been widely deployed for about the last 10 years
or so.
It's a very exciting area still.
Machine learning was originally being deployed
in big data centers that are as big as football fields that enable us to have
a wealth of application services.
For instance, when you do a Google search.
For instance, when you say you like something
or you don't like something on Netflix or Facebook.
All of those services are using machine learning, and they run up in the cloud.
But increasingly what we are seeing is that we're
witnessing that machine learning is moving more towards where
you and I are as physical beings.
Machine learning first started showing up on our smartphones.
Why?
That's because machine learning deployed at the endpoint
has the ability to provide more interactive and interesting
applications, and what we are seeing now is
that the capability that we tend to see in our phones
is now increasingly moving on to the endpoint devices like the Google
Assistant.
Why is this happening?
Well, first and foremost, data centers are large,
and they're made up of big, big servers that consume lots of energy.
They're extremely capable.
So for certain application tasks like, for instance,
when you search for information on the web,
there's a lot of benefit in going all the way up to the cloud
and looking at data across many servers and being
able to pull up a list of those results and sending them back to the end user.
That requires a remarkable amount of computational horsepower.
All these machines need to work collectively
to give us a search result, for instance.
While that is great, some of the penalties around that
is that the latency is high.
In other words, for instance, when you make a search query,
when you type something into Google search,
it has to go all the way from our device on our phone or on our computers
all the way up into the cloud.
There's a latency.
There's a high latency involved in there.
And when that latency shows up, it really
makes it hard to enable real-time machine-learning services
at the endpoint.
For instance, imagine you're taking a picture on your phone.
And from the time you take the picture, if it had to get saved locally,
sent to the cloud, and then brought back again with some enhancements,
you would actually feel that latency.
You would literally be able to see two phones, one that does it locally
and one that does it based in the cloud, and you
would say I want the phone that's able to do it locally.
And that's increasingly why we are seeing machine learning moving
more closely to where you and I are.
This does not, of course, imply that machine learning deployed in big data
centers is no longer useful.
It is absolutely useful.
It is an evolution that we're seeing.
We're seeing machine learning moving from big servers
out in remote cloud servers to endpoints like smartphones
and then onto devices that are all around us like the Google Assistant,
for instance, here.
Not only that, these devices are pervasive.
There are many devices like Google Assistant.
These endpoint devices are everywhere, and there are loads of them,
and they're all around us.
They're inside our homes.
They're on us like a smartwatch.
They're around us as well as they're also around our homes.
So I'm talking things like smart speakers, hearing aids, thermostats,
smart toothbrushes, even smart parking meters, for instance,
that know when you've actually parked your car automatically.
So TinyML is everywhere, and the fundamentals
that we have learned so far in Course 1 are the foundations for TinyML.
Based on that, we can unlock a whole new range of applications.
Now, what is the key thing for those applications?
It's not the device itself.
It's what's inside that device that gives you access to data,
to rich, rich real-time data, and that data is really,
my friend, what you want to capture.
Let's take a smartphone, for instance.
A smartphone that you and I probably have in our pockets
today is loaded with sensors.
For instance, there's an accelerometer.
An accelerometer measures the acceleration
that the handset is experiencing relative to freefall.
That allows us to determine a device's orientation along the three axes.
Then there's a gyroscope, for instance, that provides orientation information
with much greater precision.
There's a magnetometer that detects the magnetic fields, for instance,
when you open the compass app.
There's a proximity sensor that knows, for instance, when I pick up the phone
and put it right next to my ear that it can dim the screen down
and save battery life.
There's a light sensor that measures how bright
the ambient light is and then adjusts the screen display brightness
so that it can both be easier on the eyes as well as save battery.
There's a barometer inside that measures the pressure,
and that data is used to kind of figure out where exactly the device is
with respect to the sea level because that allows you to actually
have a much better position on GPS.
There's a thermometer to know how hot the phone is because sometimes they
do optimizations based on how hot the phone is.
When the phone gets very hot, they actually slow the phone down.
That's another sensor.
There's a pedometer inside, for instance, that
knows how many steps you have taken.
There's a heart-rate monitor that knows whether you're working out or not.
Just think about all these things that you
probably have been using on a daily basis
but not really been reflecting on it.
What I'm trying to get you to realize is that these sensors
that are on our devices already today have remarkable potential.
While I'm talking about these specific sensors
that I've just given you a variety of examples of, the reality of the fact
is that we're not fully exploiting these sensors to their maximum capability.
In fact, if we did that, there would be so much data being produced
and so much data being translated into rich knowledge and information
that we would have uncovered a whole new range of applications,
but we're nowhere close to that.
There's lots of good data that we don't actually leverage.
As you might recall, less than 1% of the data is being analyzed today,
and already that's changed the way we interact with machines
or what kind of application services we benefit from.
And imagine what happens when you touch the remaining 99% of the data
that we're yet to see.
So that's very exciting, right, because that can open up
a whole new range of applications, many of which
I've talked to you about right at the beginning of the course.
So I'm not going to go over them, but I want you to take a moment
and reflect back based on the fundamentals you've
learned what kind of different applications you can enable.
Now, in addition to looking at all these forward-looking applications
like robotics, drones, smart glasses, and all this stuff,
I also told you to take a moment to reflect
on how you can use this technology to actually impact the planet.
That sounds rather grandiose, but I told you,
you can do it one tiny step at a time.
You can use it for wildlife monitoring.
You can use TinyML, for instance, to monitor planets, the planet
forests, the oceans, and so forth.
There's a lot of different applications you can enable,
and we talked a little bit about this when we talked about the ElephantEdge
project, for instance.
So the key thing that we're getting to is that AI is everywhere.
As we go into Course 2, we're going to talk much more about AI everywhere.
As we do this, as exciting as it is, there
are lots of benefits we're going to get out of it,
but there are also serious risks that are about to come about,
and there are going to be some serious ethical challenges we
need to be concerned about.
Now, AI is fundamentally disruptive.
I don't mean to scare you, but it's a disruptive technology.
And whenever something disruptive comes about,
it changes, influences, or replaces things as we know them.
We need to be cognizant about what's going on.
And therefore companies are increasingly coming together,
and they're wanting us to think about what the implications are going
to be when we deploy this technology in the industry
like for predictive maintenance, for instance,
or when you use it for environmental applications
or, for instance, when we start putting all these smart devices on our bodies.
And I've talked about many different application use cases.
As we go into Course 2, we'll get much more deeper into these things.
The key thing is this.
As we start using all of this technology,
if we're going to be scared about it, we can completely
stop developing all this remarkable technology.
We don't need to learn about machine learning.
We don't need to get excited about TinyML possibilities.
We can do that.
But it's a double-edged sword.
I often like to say that you can use a knife to cut vegetables
or you can use a knife to harm someone.
It's a tool.
It's a question of how we use that tool.
And as long as we're responsible, we can use the technology for good.
And hopefully you learn something from Dr. Kennedy
when she talked about what are the considerations for AI when
we are deploying it.
In fact, before we get to the deployment,
we actually have to think about it at the design stage.
The design influences the development, and the development
influences the deployment.
Dr. Kennedy talked to us a lot about what the implications are
when we are at the design stage.
As we move into Course 2 and Course 3, we
will talk much more about what it means in terms of the development of the AI
algorithms and in terms of the deployment of the algorithms.
So what I'm trying to say is that if we shift our mindset from being
afraid of this technology being deployed everywhere and embrace it,
then companies and users will feel very good about actually enabling
this technology at a much wider scale, the scale at which we're talking
about when we talk about TinyML.
This is why you see that companies are all coming together and creating
partnerships or collaborations or consortiums
because they want their future engineers,
you and me, they want us to not only understand the technology
but to also be responsible when we're deploying that technology.
Remember the stat that I gave you earlier right at the beginning?
It was something like around 60% of people
just in America are afraid of all of these new devices
that we're talking about.
They don't know what information is being processed.
So what do they do?
They shy away from it.
Is that really what we want?
No.
We want people to feel comfortable when they're taking the technology on.
And to that end, companies for hiring the next generation of engineers,
they want us to be aware of this.
This is why Dr. Kennedy is teaching us all these different things
about responsible AI.
The cost of making a mistake when we deploy these things is really big.
It not only affects the product, one product within a company.
It affects the whole company.
It has grave consequences.
So ideally what we want to do is we want to nip these things in the bud.
The right way to do it is to learn about it
and to proactively think about it right as we are understanding technology.
And hopefully as we move into Course 2 and Course 3,
you will learn to appreciate the whole notion of responsible AI
and how it intersects with TinyML much more.
With that, my friends, I want to say I'm very much looking forward
to doing a quick recap of all the fundamental concepts
that we learned in Course 1, and I'm going to do that in the next video.
So stick with me, and then after that, we'll go on to Course 2.


\section{What we have learned thus far}

IJAY JANAPA REDDI: Hi.
What I want to do in this video is do a quick recap of where we started
and where we are wrapping up in this particular course.
I want to particularly focus on the fundamentals
that we tried to teach you.
Recall that this course is all about machine learning.
While machine learning is in itself a subfield within
artificial intelligence, what we specifically focus on in this course
is deep machine learning.
Deep machine learning is a type of machine learning
that leverages neural networks and lots of data
in order to learn interesting patterns about the data.
So these deep learning networks learn based
on data, based on patterns they have seen in the data,
and so when a similar pattern repeats in the future,
they can actually make an intelligent prediction about what's
likely to happen.
The key thing here is to remember that they learn from the data
that they have.
So if there is something in the data that they do not know about,
when the new data that they don't know about comes around,
they don't know what to make of it, so they will likely give you an error.
That's a trade off that we have to make, because when they work,
they work remarkably well.
And deep learning leverages lots of big data.
So as long as you do a good job in collecting the data,
you're going to be good.
Now, how do we train these neural networks that we talked about?
The fundamental machine learning paradigm
is this rinse and repeat cycle that we're looking at here.
We start by making a guess.
We see if we got the accuracy right.
In other words, is a cat a cat or a dog a dog,
or are we saying that a cat is actually a dog and that a dog is actually a cat?
We measure the accuracy often in terms of metrics known as top 1 or top 5.
And if we are right, then we reinforce that.
And if we are wrong, we correct that.
And so we repeat this cycle.
We continue to make a guess, we measure the accuracy,
and optimize based on the guess.
At the heart of this fundamental pipeline
is stochastic gradient descent, a really buzz word for something quite simple,
where the goal is you want to minimize the error, which Lawrence, as you might
recall, talked a lot to you about by using this ball that kind of rolls off
the hill.
What you want to do is you want to minimize the error.
And you're going to move to a place where
the ball is in a nice stable place, which is right at the bottom.
In order to do that, you can adjust the learning rates
so that you can actually move faster or you move slower,
and that's what the loss function is trying to help you do.
We talked a lot about this because there are
lots of nuances in terms of how the network learns
based on the loss function.
Sometimes you can overshoot and the ball goes all the way over to the other side
and then you can go back up.
Ultimately, you kind of do this back and forth motion
until you minimize the rocking that's happening on the ball.
So that's the fundamental notion, minimizing that error.
And out of that comes a neural network.
The topology of a neural network, in other ways,
the number of neurons you have and the way you
actually connect each of these individual neurons, that topology is
totally up to the developer's choice.
It's your choice.
It's my choice.
It's whoever is developing the neural network.
And the key function here is once you have a certain topology, what
you want to be able to do is to figure out what the weights in the network
are going to be and how you're going to adjust those weights.
And that's where the loss function comes in,
is that as you do a forward pass through this network
to make a prediction, what you want to do
is as you're training the network with the right weights,
you're trying to minimize the loss or the error
that the network is experiencing.
Now, Lawrence walked you through this by using the MS data set.
The MS data set has a whole bunch of images with handwritten formats
that allow you to kind of look at these images
and say whether it's a number that's between and 0 and 9.
And one of the things that you did by using the Google Colab
was to take these individual images, which
are 28 by 28 pixels in height and width, and feed that
into a densely connected network, sometimes known as MLPs,
or multilayer perceptrons.
And the key thing is that every neuron in there
has a connection to every other neuron on the output side.
And by doing that, there are lots of connections that show up in the middle.
That's why the mid region between the two layers of neurons
is extremely dense.
And effectively, this network will output a probability of 1
on whenever there's a strong indication that a certain digit matches
the output.
So, for example here, when it's a 7, the output neuron number 7
is actually going to go up.
Now, this is a very simple thing, but it's a very powerful concept
because you're going straight from an image to recognizing that image.
How do we do that?
One of the key things that we learn is about data augmentation or data
transformations.
Both are actually important concepts.
First we'll talk about data transformation.
You took a 28 by 28 two dimensional picture
and you flattened it into a single tensor.
28 by 28 is equal to 784 pixels.
You flatten it out.
That's how many inputs you have going in.
And that flattening can also be skewed in many different ways.
You take the 7 by 7, you can do data augmentation tricks,
which is extremely important because it is a nice way
to not have to collect lots of data.
For instance, a 7 that is really nice and straight
is one way of looking at it.
A 7 that's written at a slight angle is a different way of looking at it.
And you can do all these data transformations
without explicitly collecting pictures manually through human effort,
but instead using simple little imaging techniques,
you can skew the data set around so you can augment the data
and effectively sometimes double, triple your data set.
And by doing that, you teach your neurons
to be able to be a bit more flexible and generalized through the data set
that's coming in.
That's actually a very important concept that we learn.
So in addition to the data transformation of putting it
into the right format of the single tensor,
we also learned the data augmentation.
That's a critical concept.
And we're going to be building on this quite a lot once we
get into tinyML applications in the next course.
Now, we take this fundamental primitive that we have here
and then we rinse and repeat.
We feed in a whole bunch of different images, one after the other,
after the other, after the other, until the network reaches
a stable point where the accuracy is really high
or the loss is really, really low.
And we learned about how to split the data set up carefully between test,
training, and validation, which is extremely important in order
to be able to make sure that you are actually not overfitting to the data
that you have.
Rather, you're fitting to the data just enough.
Because overfitting is a very dangerous problem.
Because when the network sees a new behavior,
it's not going to recognize things.
It will only make predictions about the data you already have,
which is not a very useful model.
Now, after that, we moved on to extracting critical insights out
of the data set that's coming in or the input that's coming in.
For instance, in this picture that Lawrence showed you,
it's a very complex picture with a lot of rich information
and the picture on the right largely contains the same amount of information
in a very distilled format.
What are we looking at?
We're looking at the line edges for the characteristics that
are inside picture on the left.
Effectively by running this filter, which has got numbers,
and running that filter over the image on the left,
we're effectively distilling a new image on the right.
And this filter or the concept of a filter
is an extremely, extremely important concept
that you need to really internalize.
Because a filter, in effect, filters out noise.
Let's say if I'm only looking for the edges, the picture on the left
has a lot of edges.
The picture on the right has only edges.
So we're effectively removing all the noise
that we don't care about on the left hand side
and only keeping whatever we want on the right hand side,
and that is extremely important.
And neural networks effectively operate very well
when there are a whole range of these different filters
because each one of these filters is going to act and try and find
critical things.
For instance, this filter is trying to find all the edges that are there.
And another filter could be able to isolate certain levels of brightness
in the filter image.
So therefore, you tend to end up composing different filters together.
And we did that by writing code in Colab.
Lawrence just kind of walked you through a variety of different scenarios.
What we have seen is that you take TensorFlow and you learn
how to write using the Keras API calls.
And one of the key things you start off with
is knowing the syntax structure of how to program
in TensorFlow, which is a really powerful thing to be able to do.
Second thing you learn is how you actually
compose a neural network architecture.
For instance, in this particular model, you have convolutions,
then you have max pooling, and then you have dense layer networks at the end.
Now, the convolutions, it's not just a single convolution you have.
You actually have multiple convolutions, because each of these completions
is effectively trying to filter for a particular characteristic.
And so as you have multiple convolutions all kind of tied together,
they're filtering out data as it's moving through.
Nothing that you're trying to do is pulling.
If you recall, when you pull, what you do is you
minimize the information flow that's happening from the start of the network
to the end of the network.
You're effectively trying to pick out the most salient features from.
Every time you filter, you do a max pull in order to filter out
the most salient parts of that data.
And then you move it into the next filter, which
is looking for something different, and then the third filter
looks for something different.
And then you only continue to propagate the most salient features
after each filtering.
Finally, you get around to the dense network,
or sometimes it's also known as a fully connected network where
the fundamental idea is to be able to match
the filters to the labels of the output characteristics
that you and I really care about.
Now, with that fundamental network structure in place,
which you can end up doing is you repeat the network's training
process over and over and over.
The keyword here is epochs, which is something
that we will use a lot when we're training our tinyML applications.
Here, epochs 12 means that we are going to repeat the training
cycle over the data 12 times.
We're going to go over all the data we have 12 times.
All the data.
And so by doing that, we can actually get a much better accuracy
because the more time we train, the better it is.
Now, that's great.
But remember, training for longer is not always the right thing.
There's often a trade off overfitting to the data
or another thing is to remember that when you're training,
you're using very costly resources, so that drives up
the cost of training the network.
And that's not necessarily a good thing if you're strapped for cash.
One of the things that's happening behind the scenes
when you're training the networks that you might have taken for granted
is the fact that there are these big, beefy processors that you're
running these machine learnings on.
Training is a very, very, very costly process in reality.
So therefore, companies like Google have gone out
and built custom hardware just dedicated for training machine learning.
While in this particular course, you have yet
to use something like a cloud KPU, you are most definitely using
GPUs, which help you train faster.
Over the course of your experience in machine learning,
you will use CPUs, GPUs, KPUs, DSPs, NPUs.
There's a wide variety of different kinds of hardware.
I will try and teach you a little bit about this
as we get into the more advanced concepts way down in the later courses.
Now, remember that these TPUs are all packaged
into big data centers and these pods, as Google likes to call them,
are what are powering some of the machine learning services.
So when you train a network, it might very well
be running on one of these servers for all you know.
Now, that's all on the algorithm side.
Now, one of the key things is that that fundamental knowledge that you acquired
has been repurposed so many times over by big companies by building
deeper, and deeper, and deeper, and deeper, and deeper networks, so much so
that the number of parameters that are involved in these machine learning
models, sometimes algorithms, depending on how you refer to them,
are so big that it's in the order of billions of parameters now.
And that's what this chart is showing.
It's showing that the state of the art models
has so many billions of parameters in terms of the number of neurons
you have in the network and the number of connections
that are there between all the neurons.
It's so big that it can actually learn a lot of information.
So this computational capability and having a lot of deep networks
allows us to learn a lot of things.
That's great.
One of the things that's been happening, obviously,
is from AlexNet, which started out in 2012, as you might recall.
Models have been getting better.
In other words, they've been improving in terms of accuracy.
We talked about AlexNet moving into VGG as a big, major milestone.
Why did we do that?
Because as you saw, the VGG models improved
accuracy, which is on the y-axis here.
And the x-axis, if you don't recall, is actually
the amount of computational horsepower.
And the size of the circle here tells you
how big the model is, how big or lean and thin it is.
As models got bigger, we wanted to make them smaller because we were becoming
more aware that resources are costly.
It does cost a lot of money to train them all,
so we want to make them smaller.
But making them smaller does not mean that we lose accuracy.
We're actually able to boost accuracy.
For example, in going from EGG to ResNet,
we actually are able to boost the accuracy
while keeping the models smaller.
But then as smartphones came about, we felt a desperate need
in order to make the models much smaller,
and so we tried to shrink them down.
Now, the question is this.
As we move from big data centers, to mobile devices,
to the next major horizon and ecosystem, how
are we going to make these models very small?
How small?
Well, if you recall, a typical MobileNet model,
which is designed for smartphones, for instance,
occupies as much as around 17 megabytes, which is a lot of memory.
When, in fact, the microcontroller, for instance,
that we will use in course 3, which is a very typical microcontroller,
has only about 256 kilobytes.
That's a massive order of magnitude difference between where we are today
and where we need to be in the future.
And if we want to deploy machine learning, all
these billions of microcontrollers, we're
going to need to do something big.
And that means we have to innovate on the machine learning model stat.
We have to figure a way out not only to make the machine early model smaller,
but we also have to figure out how to cram those models
on tiny little hardware resources.
Now, that's all in the machine learning algorithm side.
But remember I also told you that there's an embedded systems side.
So we talked a lot about this and we tried
to really understand what are the fundamental constraints
in a microcontroller as compared to a microprocessor
that, for instance, is running one of those machine learning
models in the cloud or running it on your laptop.
From compute, to memory, to storage, and the amount of power
they consume, there is a significant order
of magnitude difference between the capabilities of a regular processor
and what the capabilities of all these tiny machine learning processors are.
Now, this seems really bad.
But remember, we want to do this because tiny devices are everywhere.
They're ubiquitous.
They're everywhere.
In the future, believe it or not, I won't be surprised
if we put machine learning on a tiny little device with a quantum battery
and we forget it.
We just forget it, and it's processing, and collecting information,
and doing interesting sensing things.
And someday it just biodegrades and so we're not polluting the Earth.
Rather, we have this remarkable technology
providing value add to all of us.
But all of that means we have very limited resources.
So one of the key things that happens, because power
is so critical in these systems, we understood, for instance,
how vendors have to make trade offs in how to build the device.
So in the embedded ecosystem, you tend to have
a lot of fragmentation, hardware fragmentation, software fragmentation.
We talked about how the floating point unit, a simple floating point
unit that sometimes is present on a device
might not be present on another device, which
means you have to recompile all the code for a special device.
That becomes a major problem.
Why do we do something like that which breaks compatibility?
We do it because it lowers the cost, it lowers
the amount of power consumption needed, and it also
reduces the engineering effort, which means
we can crank out these tiny little devices like no other.
But the trade off is portability, which introduces problems
to us that we will learn how to address in course 2 and course 3.
FP example is on the hardware side.
Remember I also talk to you about how the laptop and server
ecosystem has these rich operating systems that virtually
dominate all of these computers?
Same thing in the mobile ecosystem.
We have iOS and Android, and they're amazingly dominant,
and they basically run the entire world.
But what happens when you go to an embedded system?
On EmbedOS and FreeArc OS, these are operating systems that exist,
but not every embedded system uses them.
Sometimes there there's no operating system.
There's no concept of running multiple programs on the same embedded system.
Every device is different, so this creates a massive problem
for being able to systematically make models smaller
and being able to deploy them efficiently.
But don't worry.
All this said, we have solutions that are coming about
and that's what we're going to really learn next.
As we learn about models in this, when we go to the course,
you will learn much more about making the models smaller.
In addition, we'll also learn about the run times that need to be adapted,
and I'll talk about this in the next video.
And also, we need to understand what the capabilities of the hardware
are and do a codesign between the models, the runtimes, and the hardware.
Remember, I promised you one thing when we started off this course,
and it is that we're going to learn not only about each
of the individual layers, but we're actually
going to learn about the interactions that happen between the layers.
And that's extremely critical because that's
what's going to make you stand out from the rest.
Now, with that said, I'm going to leave you with some things to reflect
and I'll come back to you in the next video.

\section{What's coming next}

VIJAY REDDI: Hi there.
This is going to be the last video for Course 1.
And after this, we're going to be stepping into Course 2, which
I'm very excited about.
As we said in Course 1, it's all about the language and the fundamentals
of tiny machine learning.
That involves understanding what's tiny.
That's the embedded side.
What's TinyML?
Well, that's all about the ML fundamentals
that you need to know in order to be able to build a TinyML application.
As we go into Course 2, we're going to focus very much
on taking the bread and butter things that we've learned in Course 1
and actually building interesting applications.
And we will be using neural networks.
We will be using TensorFlow.
We will be using Google Colab.
Everything that you have learned so far, we're going to be building on it
and expanding our knowledge.
Let me just take you head in and say that we'll
be looking at some seriously interesting applications.
We'll be looking at building TinyML applications for speech, for vision,
and things that use different kinds of sensors,
like IMU, accelerometer, and gyroscope kind of things.
For instance, if you recall right at the beginning of Course 1,
we started talking about speech with OK Google or Alexa or Hey Siri.
In fact, it was one of the first things that I had in the very first slide deck
that I showed you.
In this particular example, with speech for instance,
we will be looking into training our own model that
responds to our own speech commands.
So you will have your own pick or choice of the model
that you want, trained specifically for the words that you care about.
You might care about waking the machine up when you say "hello,"
and it will respond.
We'll be using a variety of different sensors.
And that's how the assignments are all set up.
So when you'll be programming Google Colab,
for speech you'll be using the microphone,
for vision you'll be using the camera input,
for IMU you will be using different kinds of sensor data.
But given that you'll be programming in the Colab environment,
we will be doing all of this with the data already given to you or data
that you can collect, for instance, in the form of files.
We'll get much more into this, but for now, I just
want you to be excited knowing that you're actually going
to be building TinyML applications.
Because now you have the fundamentals that
are needed to be able to build the TinyML applications.
But the key thing that we'll also be focusing on
is, how do we make those small models not just be functional,
in terms of being able to do the speech, vision, or IMU use cases,
but actually making them really, really small
so they have the potential to fit in an embedded device?
That's going to be very exciting.
In order to do that, we'll be talking about different kinds of techniques,
like pruning, for instance-- pulling out certain networks without really
losing any accuracy in the network.
We'll be talking about using different kinds of numerical formats,
for instance.
Instead of using float32, which requires you to use four bytes of data,
we will be using int8, which only requires one byte of data.
So by compressing the network down by using different numerical formats,
and hopefully without losing accuracy-- there's some nuances in there that we
will talk about and learn about--
but by doing that, you can actually make the model very, very small.
And in doing so, you can have the potential
to put it onto a small, little, embedded device.
We'll also be looking at more advanced things like knowledge distillation,
for instance, where you can take a big network
and distill critical information out of it so that you can create a smaller
network that effectively has the same knowledge,
but has a totally different topology.
We'll talk about these kinds of different optimizations.
Now while we'll talk about the models a lot,
one of the key things that I also want to help you learn
is about the actual underneath-the-hood runtime system.
Because that's really critical for deploying machine learning efficiently.
What do I mean by that?
So far, when you have been training the model,
and when you have been programming things in Google Colab,
you've been focused on actually using things
like big GPUs that are in Google's data centers to get the trained model.
But once we start going into actually deploying the model
onto a variety of small devices--
for example, like trying to deploy it on Raspberry Pi
or on a tiny little microcontroller, for instance-- well,
that's going to be really challenging if you take the trained modeling
and put it right on it.
It will not fit.
We actually have to do a slew of different kinds of optimizations
and conversions.
First of all, we're going to learn what it
means to convert a trained TensorFlow model into a TFLite model.
We'll get into the details of what TFLite really is
and why we need to do that next.
Then we'll learn how to optimize that TFLite model so that it actually
is quantized, and so that it actually uses less amount of memory,
and so forth.
And after that, we will look deeper into what
are the APIs that are involved, because the APIs that you will need
as an application developer who is using TFLite
are slightly different from the APIs that you have been exposed to so far.
One of my goals out of Course 1, 2, and 3
is to really teach you how to program in these frameworks,
so that at the end of the day, you actually have a job prospect.
Because you will be able to say that I know how to program in TensorFlow,
I know how to program in TensorFlow Lite,
and I also know how to program in TensorFlow Lite Micro.
I'll get you there by exposing you to the variety of different API calls--
so similar to what Laurence did, in terms of teaching you
how to program in TensorFlow, I'll walk you
through a little bit in programming TensorFlow Lite.
And by the way, Laurence Moroney is going
to be with us right at the beginning of Course 2
to help us go from understanding TensorFlow into TensorFlow Lite.
Then as we get more into Course 3--
I'm looking two steps ahead now-- we'll be
talking about deploying the TinyML models that we have built in Course 2
onto the dedicated hardware.
That requires us to do a couple of more different transformations,
because we have to be a bit more cognizant about the embedded system.
To that end, we'll go even deeper into TensorFlow Lite
by specializing into TensorFlow Lite for Microcontrollers.
And ultimately, after the deployment, we'll
go into using that device to do interesting applications.
The last part is all going to be about the hardware in Course 3.
And so for that, you will have to start thinking about procuring a kit, which
we'll share the details for soon.
The kit effectively comes completely self-contained.
It comes with cables, a camera, a breadboard-- everything
you need in order to enable the different kinds of applications
that we're talking about.
In fact, I'll do even better than that.
I'll try and teach you about some potential applications
that you might be able to build based on the fundamentals
you pick up in Course 2, so that you can actually do things on your own.
And with that said, I want to leave you thinking
that there's an exciting new course coming up
that you want to be a part of, because in that course,
all the knowledge you have acquired so far
is actually going to be put into practice.
You will learn to develop tiny machine learning models.
And you will learn how to actually optimize them.
And you will learn the software that's actually involved behind the scenes.
And as we push that even further into Course 3,
you will learn to deploy it on the hardware,
and you'll be able to start building real, physical-world applications that
are in your hands and being able to show them off to everybody.
So I'm very excited to see you in Course 2.
With that, have a good time.

